% Point de fonctionnement d'un circuit et loi de Pouillet — Physique 1ère D — Cours signé OnBuch+
% Programme officiel MINESEC. Compiler avec : tectonic N15-point-fonctionnement-loi-pouillet.tex
\newcommand{\DOCMATIERE}{Physique}
\newcommand{\DOCNIVEAU}{1ère D}
\newcommand{\DOCMODULE}{Module 4 : Aspects énergétiques des circuits électriques}
\newcommand{\DOCLECON}{N15}
\newcommand{\DOCTITRE}{Point de fonctionnement d'un circuit et loi de Pouillet}
% OnBuch+ — préambule « cours signés » ENRICHI (compilé avec tectonic / XeLaTeX)
% Système visuel « Pop » d'OnBuch+ : fond crème (jamais blanc), bordures encre
% épaisses, ombres « dures » décalées, titres Archivo Black en capitales.
% Version MATHÉMATIQUES : graphes (pgfplots/tikz), boîtes pédagogiques
% (définition, propriété, méthode, attention, le savais-tu, exercice résolu,
% exercice, TP, bilan), page de garde et signature OnBuch+ ancrée en bas.
%
% Macros attendues dans main.tex AVANT \input{preamble} :
%   \DOCMATIERE  \DOCNIVEAU  \DOCMODULE  \DOCLECON (numéro)  \DOCTITRE
\documentclass[11pt]{article}

\usepackage{fontspec}
\usepackage[french]{babel}
\usepackage[a4paper,margin=2cm,top=3.4cm,bottom=2.8cm,headheight=1.4cm,footskip=1.3cm]{geometry}
\usepackage{amsmath,amssymb,amsfonts}
\usepackage{mathtools} % \xmapsto, \coloneqq... souvent utilisés par les modèles
\usepackage{cancel} % simplifications barrées (bilans de forces)
% \abs{z} (module, valeur absolue) et \norm{u} (norme), utilisés par les modèles
\providecommand{\abs}{}\let\abs\relax\DeclarePairedDelimiter{\abs}{\lvert}{\rvert}
\providecommand{\norm}{}\let\norm\relax\DeclarePairedDelimiter{\norm}{\lVert}{\rVert}
\usepackage[table]{xcolor} % [table] : fournit aussi \rowcolors (alternance de lignes)
\usepackage{pagecolor}
\usepackage{tikz}
\usetikzlibrary{arrows.meta,positioning,calc,shapes.geometric,shapes.misc,shapes.arrows,decorations.pathmorphing,decorations.pathreplacing,decorations.markings,patterns,shadows,mindmap,trees,fit,backgrounds,matrix,intersections,angles,quotes}
\usepackage{pgfplots}
\usepackage[version=4]{mhchem} % équations nucléaires \ce{^{235}_{92}U}
\usepackage[european,siunitx]{circuitikz} % schémas électriques
\pgfplotsset{compat=1.18}
\usepgfplotslibrary{fillbetween,groupplots} % aires entre courbes (calcul intégral)
\usepackage{enumitem}
\usepackage{fancyhdr}
% [most] charge toutes les bibliothèques tcolorbox (listings, minted, raster,
% documentation...) et gonfle inutilement la mémoire TeX sur un document long
% (~300 boîtes) : on ne charge que ce qui est réellement utilisé.
\usepackage[skins,breakable]{tcolorbox}
\usepackage{graphicx}
\usepackage{colortbl}
\usepackage{booktabs}
\usepackage{tabularx}
\usepackage{diagbox} % case d'en-tête diagonale des tableaux à double entrée
% Colonnes extensibles centrée / gauche / droite (C, L, R), souvent utilisées par les modèles
\newcolumntype{C}{>{\centering\arraybackslash}X}
\newcolumntype{L}{>{\raggedright\arraybackslash}X}
\newcolumntype{R}{>{\raggedleft\arraybackslash}X}
\usepackage{array}
\usepackage{multicol}
\usepackage{siunitx}
% parse-numbers=false : accepte aussi \SI{2,5 \times 10^{-3}}{...}
\sisetup{locale=FR,output-decimal-marker={,},group-minimum-digits=4,parse-numbers=false}
\DeclareSIUnit\ohm{\ensuremath{\symup{\Omega}}} % U+2126 absent de la police
\DeclareSIUnit\division{div} % réglages d'oscilloscope (ms/div, V/div)
\DeclareSIUnit\tr{tr} % tours (vitesse de rotation, ex. \SI{1500}{\tr\per\minute})
\DeclareSIUnit\ch{ch} % cheval-vapeur (puissance moteur, unité usuelle hors SI)
\DeclareSIUnit\franccfa{FCFA} % franc CFA (coûts, contexte camerounais)
\DeclareSIUnit\dioptre{\ensuremath{\delta}} % vergence d'une lentille (dioptrie)
\usepackage{qrcode}
% \degree : commande fréquemment utilisée par les modèles pour un angle ou une
% température en dehors de \SI{}{} (non fournie par les paquets chargés).
\providecommand{\degree}{^{\circ}}
% \qed (fin de démonstration), utilisé par les modèles sans amsthm
\providecommand{\qed}{\hfill$\square$}
% \barre : double barre des valeurs interdites dans un tableau de variations (tkz-tab)
\providecommand{\barre}{\ensuremath{\|}}
% environnement proof (amsthm non chargé) utilisé par les modèles
\ifdefined\proof\else\newenvironment{proof}[1][Démonstration]{\par\noindent\textit{#1.}\ }{\qed\par}\fi
\usepackage{lastpage}
\usepackage{hyperref}
\hypersetup{hidelinks}

% ============================================================
% Palette « Pop » OnBuch+ — mêmes hex que la classe P (plus_ui.dart)
% ============================================================
\definecolor{poporange}{HTML}{F59321}
\definecolor{popdark}{HTML}{E07A0C}
\definecolor{popcream}{HTML}{FFF6E8}
\definecolor{popink}{HTML}{1C1714}
\definecolor{poppurple}{HTML}{7A5AE0}
\definecolor{popgreen}{HTML}{2E9E6A}
\definecolor{poppink}{HTML}{E0457B}
\definecolor{popblue}{HTML}{2F7DE1}
\definecolor{popgold}{HTML}{C98A12}
\definecolor{popmuted}{HTML}{8A7F76}
\definecolor{popline}{HTML}{EADFD2}
\definecolor{popgray}{HTML}{8A7F76}
\colorlet{popgrey}{popgray}
% Alias tolérants (le modèle nomme parfois les couleurs différemment)
\colorlet{popwhite}{white}
\colorlet{popblack}{popink}
\colorlet{popred}{poppink}
\colorlet{popyellow}{popgold}
% Teintes claires pour fonds de tableaux / zones
\colorlet{popblueL}{popblue!10!white}
\colorlet{poporangeL}{poporange!14!white}
\colorlet{popgreenL}{popgreen!12!white}
\colorlet{poppurpleL}{poppurple!12!white}
\colorlet{poppinkL}{poppink!12!white}
\colorlet{popgoldL}{popgold!14!white}

\pagecolor{popcream}

% ============================================================
% Typographie
% ============================================================
\defaultfontfeatures{Ligatures=TeX}
\setmainfont{PlusJakartaSans-Medium.ttf}[Path=../../../fonts/,BoldFont=PlusJakartaSans-Bold.ttf]
% Maths en sans-serif (Fira Math) avec chiffres et lettres droites de Plus
% Jakarta, pour que les formules se fondent dans le texte.
\usepackage[math-style=ISO,bold-style=ISO]{unicode-math}
\setmathfont{FiraMath-Regular.otf}[Path=../../../fonts/]
\setmathfont{PlusJakartaSans-Medium.ttf}[Path=../../../fonts/,range={up/num,up/{Latin,latin},\mathup}]
\usepackage{icomma} % virgule décimale sans espace en mode math
% Caractères mathématiques Unicode absents des polices de texte (Plus Jakarta,
% Archivo Black) : rendus par leur équivalent mathématique, en texte comme en
% formule — sinon ils s'affichent en carré vide (ex. « Arithmétique dans ℤ »).
\usepackage{newunicodechar}
\newunicodechar{ℝ}{\ensuremath{\mathbb{R}}}
\newunicodechar{ℤ}{\ensuremath{\mathbb{Z}}}
\newunicodechar{ℂ}{\ensuremath{\mathbb{C}}}
\newunicodechar{ℕ}{\ensuremath{\mathbb{N}}}
\newunicodechar{ℚ}{\ensuremath{\mathbb{Q}}}
\newunicodechar{𝔻}{\ensuremath{\mathbb{D}}}
\newunicodechar{α}{\ensuremath{\alpha}}
\newunicodechar{β}{\ensuremath{\beta}}
\newunicodechar{γ}{\ensuremath{\gamma}}
\newunicodechar{δ}{\ensuremath{\delta}}
\newunicodechar{ε}{\ensuremath{\varepsilon}}
\newunicodechar{θ}{\ensuremath{\theta}}
\newunicodechar{λ}{\ensuremath{\lambda}}
\newunicodechar{μ}{\ensuremath{\mu}}
\newunicodechar{σ}{\ensuremath{\sigma}}
\newunicodechar{τ}{\ensuremath{\tau}}
\newunicodechar{φ}{\ensuremath{\varphi}}
\newunicodechar{ψ}{\ensuremath{\psi}}
\newunicodechar{ω}{\ensuremath{\omega}}
\newunicodechar{Σ}{\ensuremath{\Sigma}}
% majuscules produites par \MakeUppercase dans les titres (α-aminés -> Α)
\newunicodechar{Α}{\ensuremath{\alpha}}
\newunicodechar{Β}{\ensuremath{\beta}}
\newunicodechar{Γ}{\ensuremath{\Gamma}}
\newunicodechar{Δ}{\ensuremath{\Delta}}
\newunicodechar{Ε}{\ensuremath{\varepsilon}}
\newunicodechar{Θ}{\ensuremath{\Theta}}
\newunicodechar{Λ}{\ensuremath{\Lambda}}
\newunicodechar{Μ}{\ensuremath{\mu}}
\newunicodechar{Τ}{\ensuremath{\tau}}
\newunicodechar{Φ}{\ensuremath{\Phi}}
\newunicodechar{Ψ}{\ensuremath{\Psi}}
\newunicodechar{Ω}{\ensuremath{\Omega}}
\newunicodechar{↦}{\ensuremath{\mapsto}}
\newunicodechar{∈}{\ensuremath{\in}}
\newunicodechar{∉}{\ensuremath{\notin}}
\newunicodechar{∧}{\ensuremath{\wedge}}
\newunicodechar{⇔}{\ensuremath{\Leftrightarrow}}
\newunicodechar{⟺}{\ensuremath{\Longleftrightarrow}}
\newunicodechar{⇒}{\ensuremath{\Rightarrow}}
\newunicodechar{⟹}{\ensuremath{\Longrightarrow}}
\newunicodechar{∘}{\ensuremath{\circ}}
\newunicodechar{∪}{\ensuremath{\cup}}
\newunicodechar{∩}{\ensuremath{\cap}}
\newunicodechar{≡}{\ensuremath{\equiv}}
\newunicodechar{⊂}{\ensuremath{\subset}}
\newunicodechar{⊥}{\ensuremath{\perp}}
\newunicodechar{∃}{\ensuremath{\exists}}
\newunicodechar{∀}{\ensuremath{\forall}}
\newunicodechar{∛}{\ensuremath{\sqrt[3]{\,}}}
\newunicodechar{∜}{\ensuremath{\sqrt[4]{\,}}}
\newunicodechar{⃗}{\ensuremath{{}^{\to}}}
\newunicodechar{ⁿ}{\ifmmode^{n}\else\textsuperscript{\ensuremath{n}}\fi}
\newunicodechar{ˣ}{\ifmmode^{x}\else\textsuperscript{\ensuremath{x}}\fi}
\newunicodechar{ᵇ}{\ifmmode^{b}\else\textsuperscript{\ensuremath{b}}\fi}
\newunicodechar{ʳ}{\ifmmode^{r}\else\textsuperscript{\ensuremath{r}}\fi}
\newunicodechar{ᵘ}{\ifmmode^{u}\else\textsuperscript{\ensuremath{u}}\fi}
\newunicodechar{ʸ}{\ifmmode^{y}\else\textsuperscript{\ensuremath{y}}\fi}
\newunicodechar{ᵃ}{\ifmmode^{a}\else\textsuperscript{\ensuremath{a}}\fi}
\newunicodechar{ᵉ}{\ifmmode^{e}\else\textsuperscript{\ensuremath{e}}\fi}
\newunicodechar{ᵏ}{\ifmmode^{k}\else\textsuperscript{\ensuremath{k}}\fi}
\newunicodechar{ᵖ}{\ifmmode^{p}\else\textsuperscript{\ensuremath{p}}\fi}
\newunicodechar{ᴺ}{\ifmmode^{N}\else\textsuperscript{\ensuremath{N}}\fi}
\newunicodechar{ᶜ}{\ifmmode^{c}\else\textsuperscript{\ensuremath{c}}\fi}
\newunicodechar{ʰ}{\ifmmode^{h}\else\textsuperscript{\ensuremath{h}}\fi}
\newunicodechar{ᵠ}{\ifmmode^{q}\else\textsuperscript{\ensuremath{q}}\fi}
\newunicodechar{ₙ}{\ifmmode_{n}\else\textsubscript{\ensuremath{n}}\fi}
\newunicodechar{ₐ}{\ifmmode_{a}\else\textsubscript{\ensuremath{a}}\fi}
\newunicodechar{ᵢ}{\ifmmode_{i}\else\textsubscript{\ensuremath{i}}\fi}
\newunicodechar{ᵣ}{\ifmmode_{r}\else\textsubscript{\ensuremath{r}}\fi}
\newunicodechar{ₖ}{\ifmmode_{k}\else\textsubscript{\ensuremath{k}}\fi}
\newunicodechar{ᵦ}{\ifmmode_{\beta}\else\textsubscript{\ensuremath{\beta}}\fi}
\newunicodechar{ₓ}{\ifmmode_{x}\else\textsubscript{\ensuremath{x}}\fi}
\newfontfamily\popdisplay{ArchivoBlack-Regular.ttf}[Path=../../../fonts/]
\newfontfamily\poplogo{SpaceGrotesk-Bold.ttf}[Path=../../../fonts/]
\newfontfamily\popsemi{PlusJakartaSans-SemiBold.ttf}[Path=../../../fonts/]
\newfontfamily\popxbold{PlusJakartaSans-ExtraBold.ttf}[Path=../../../fonts/]
\newfontfamily\popxboldls{PlusJakartaSans-ExtraBold.ttf}[Path=../../../fonts/,LetterSpace=3.0]

\setlist[enumerate]{leftmargin=1.7em,itemsep=5pt,topsep=4pt}
\setlist[itemize]{leftmargin=1.7em,itemsep=4pt,topsep=4pt,label={\color{poporange}\textbullet}}
\setlength{\parindent}{0pt}
\setlength{\parskip}{5pt}
\renewcommand{\arraystretch}{1.25}

% pgfplots : style Pop par défaut
\pgfplotsset{
  every axis/.append style={axis line style={popink,line width=1pt},tick style={popink},
    grid style={popline,line width=0.5pt},label style={font=\small},tick label style={font=\footnotesize}},
  popcurve/.style={poporange,line width=1.8pt,no markers,smooth},
}
% Même style disponible dans un \draw TikZ simple (hors pgfplots)
\tikzset{popcurve/.style={poporange,line width=1.8pt}}
\tikzset{popgrid/.style={popline,very thin}}
\colorlet{popcurve}{poporange} % le nom du style est parfois utilisé comme couleur

% ============================================================
% Pill / squiggle
% ============================================================
\newcommand{\poppill}[3]{%
  \tikz[baseline=-0.55ex]\node[draw=popink,line width=1.2pt,rounded corners=2.6pt,%
    fill=#2,inner xsep=6.5pt,inner ysep=3pt]{%
    \popxboldls\fontsize{7}{7}\selectfont\color{#3}\MakeUppercase{#1}};%
}
\newcommand{\popsquiggle}[1]{%
  \tikz[baseline=-0.4ex]{
    \draw[#1,line width=2.6pt,line cap=round]
      plot [smooth] coordinates {(0,0.09) (0.34,0.01) (0.68,0.21) (1.02,0.01) (1.36,0.10)};
  }%
}
% Mot-clé surligné (marqueur orange sous le texte)
\newcommand{\cle}[1]{{\popxbold\color{popdark}#1}}

% ============================================================
% En-tête / pied
% ============================================================
\pagestyle{fancy}
\fancyhf{}
\renewcommand{\headrulewidth}{0pt}
\renewcommand{\footrulewidth}{0pt}
\lhead{\raisebox{-0.15ex}{\poplogo\fontsize{13}{13}\selectfont\color{popink}OnBuch\color{poporange}+}}
\rhead{\poppill{\DOCMATIERE\ \textbullet\ \DOCNIVEAU\ \textbullet\ Leçon \DOCLECON}{white}{popink}}
\lfoot{\popxbold\fontsize{7.6}{7.6}\selectfont\color{popmuted}\MakeUppercase{Cours signé OnBuch+}\ \textbullet\ {\popsemi\DOCTITRE}}
\rfoot{\tikz[baseline=-0.6ex]\node[rounded corners=8.5pt,fill=popink,minimum height=17pt,minimum width=17pt,inner xsep=5pt,inner sep=0]{\popxbold\fontsize{8}{8}\selectfont\color{popcream}\thepage\,/\,\pageref*{LastPage}};}

% ============================================================
% Titres
% ============================================================
\newcommand{\chaptitre}[2]{%
  \begin{tcolorbox}[enhanced,colback=poporange,colframe=popink,boxrule=2.6pt,arc=11pt,
      drop shadow=popink,left=15pt,right=15pt,top=12pt,bottom=11pt,before skip=3pt,after skip=18pt]
    \poppill{#2}{popink}{popcream}\par\vspace{9pt}
    {\raggedright\hyphenpenalty=10000\exhyphenpenalty=10000\popdisplay\fontsize{22}{24}\selectfont\color{popcream}\MakeUppercase{#1}\par}
  \end{tcolorbox}
}
% Section numérotée : pastille orange avec le numéro + titre Archivo
\newcounter{popsec}
\newcommand{\coursec}[1]{%
  \stepcounter{popsec}\setcounter{popsub}{0}%
  \par\vspace{16pt}\noindent
  \begin{minipage}{\linewidth}
  \tikz[baseline=-0.7ex]\node[draw=popink,line width=1.6pt,fill=poporange,rounded corners=4pt,minimum size=22pt,inner sep=0]{\popdisplay\fontsize{12}{12}\selectfont\color{popcream}\thepopsec};%
  \hspace{8pt}{\popdisplay\fontsize{15}{17}\selectfont\color{popink}\MakeUppercase{#1}}\par
  \vspace{4pt}\noindent{\color{poporange}\rule{36pt}{3.6pt}}
  \end{minipage}\par\nopagebreak\vspace{8pt}
}
\newcounter{popsub}
\newcommand{\courssub}[1]{%
  \stepcounter{popsub}%
  \par\vspace{9pt}\noindent{\popxbold\fontsize{11.5}{13}\selectfont\color{popink}{\color{poporange}\thepopsec.\thepopsub}\hspace{6pt}#1}\par\nopagebreak\vspace{4pt}
}

% ============================================================
% Boîtes pédagogiques — même recette : fond blanc, bordure épaisse colorée,
% ombre dure encre, badge plein. Titre optionnel [#1].
% ============================================================
\tcbset{popbase/.style={breakable,enhanced,colback=white,boxrule=2.3pt,arc=9pt,
  shadow={3pt}{-3pt}{0pt}{popink},left=14pt,right=14pt,top=11pt,bottom=11pt,before skip=11pt,after skip=15pt,
  fontupper=\popsemi\fontsize{10.6}{14.5}\selectfont\color{popink},
  fontlower=\popsemi\fontsize{10.6}{14.5}\selectfont\color{popink}}}
% Coche dessinée (le glyphe ✓ n'existe pas dans les polices du document)
\newcommand{\popcheck}[1]{\tikz[baseline=-0.5ex]\draw[#1,line width=1.6pt,line cap=round,line join=round](-0.13,0.0)--(-0.04,-0.09)--(0.13,0.1);}
\providecommand{\coche}{\popcheck{popgreen}}
\providecommand{\resultat}[1]{\par\smallskip\noindent\fcolorbox{popgreen}{popgreenL}{\parbox{\dimexpr\linewidth-2\fboxsep-2\fboxrule\relax}{\textbf{Résultat.} #1}}\par\smallskip}
\newcommand{\popboxhead}[3]{\poppill{#1}{#2}{white}\ifx\relax#3\relax\else\hspace{6pt}{\popxbold\fontsize{10.6}{12}\selectfont\color{popink}#3}\fi\par\vspace{7pt}}

\newtcolorbox{aretenir}[1][]{popbase,colframe=popgreen,before upper={\popboxhead{À retenir}{popgreen}{#1}}}
\newtcolorbox{exemplebox}[1][]{popbase,colframe=poporange,before upper={\popboxhead{Exemple}{poporange}{#1}}}
\newtcolorbox{definition}[1][]{popbase,colframe=popblue,before upper={\popboxhead{Définition}{popblue}{#1}}}
\newtcolorbox{propriete}[1][]{popbase,colframe=poppurple,before upper={\popboxhead{Propriété}{poppurple}{#1}}}
\newtcolorbox{methode}[1][]{popbase,colframe=popgold,before upper={\popboxhead{Méthode}{popgold}{#1}}}
\newtcolorbox{attention}[1][]{popbase,colframe=poppink,before upper={\popboxhead{Attention}{poppink}{#1}}}
\newtcolorbox{experience}[1][]{popbase,colframe=popblue,colback=popblueL,before upper={\popboxhead{Expérience}{popblue}{#1}}}
% « Le savais-tu ? » : carte violette pleine (contexte, culture, Cameroun)
\newtcolorbox{savaistu}[1][]{popbase,colback=poppurple,colframe=popink,
  fontupper=\popsemi\fontsize{10.4}{14.2}\selectfont\color{white},
  before upper={\poppill{Le savais-tu\,?}{white}{poppurple}\ifx\relax#1\relax\else\hspace{6pt}{\popxbold\color{white}#1}\fi\par\vspace{7pt}}}
% Exercice résolu : énoncé en haut, \tcblower puis la solution sur fond crème
\newcounter{popexo}\newcounter{popexr}
\newtcolorbox{exoresolu}[1][]{popbase,colframe=poporange,colbacklower=popcream,
  segmentation style={popink,line width=1.2pt,dashed},
  before upper={\stepcounter{popexr}\popboxhead{Exercice résolu \thepopexr}{poporange}{#1}},
  before lower={\poppill{Solution}{popink}{popcream}\par\vspace{6pt}}}
% Exercice (non résolu en place) — la correction est dans la partie Corrigés
\newtcolorbox{exercice}[1][]{popbase,colframe=popink,
  before upper={\stepcounter{popexo}\popboxhead{Exercice \thepopexo}{popink}{#1}}}
% Corrigé
\newtcolorbox{corrige}[1][]{popbase,colframe=popgreen,colback=popgreenL,
  before upper={\popboxhead{Corrigé}{popgreen}{#1}}}
% Objectifs / prérequis / bilan
\newtcolorbox{objectifs}{popbase,colframe=popink,colback=poporangeL,
  before upper={\popboxhead{Objectifs de la leçon}{popink}{}}}
\newtcolorbox{prerequis}{popbase,colframe=popink,colback=popblueL,
  before upper={\popboxhead{Prérequis}{popblue}{}}}
\newtcolorbox{bilan}{popbase,colframe=popink,colback=white,boxrule=2.8pt,
  before upper={\popboxhead{Fiche bilan}{poporange}{L'essentiel à connaître pour l'examen}}}
% Figure encadrée avec légende
\newtcolorbox{popfigure}[1][]{enhanced,breakable=false,colback=white,colframe=popline,boxrule=1.4pt,arc=8pt,
  left=10pt,right=10pt,top=10pt,bottom=8pt,before skip=10pt,after skip=12pt,halign=center,
  lower separated=false,halign lower=center,
  fontlower=\popsemi\fontsize{9}{11.5}\selectfont\color{popmuted},
  #1}
\newcommand{\legende}[1]{\par\vspace{6pt}{\popsemi\fontsize{9}{11.5}\selectfont\color{popmuted}#1}}

% ============================================================
% Signature OnBuch+
% ============================================================
\newcommand{\poplogoinline}[1]{{\poplogo\fontsize{#1}{#1}\selectfont\color{popink}OnBuch\color{poporange}+}}
\newcommand{\signatureonbuch}{%
  \begin{tcolorbox}[enhanced,colback=popink,colframe=popink,boxrule=2.2pt,arc=10pt,
      drop shadow=poporange,left=14pt,right=14pt,top=10pt,bottom=10pt,before skip=0pt,after skip=0pt]
    \tikz[baseline=-0.7ex]\node[circle,fill=popgreen,draw=popcream,line width=1.3pt,minimum size=22pt,inner sep=0]{\popcheck{white}};%
    \hspace{9pt}\begin{minipage}[c]{0.62\linewidth}
      {\popxbold\fontsize{12}{13}\selectfont\color{popcream}Signé\ \textbullet\ {\poplogo OnBuch\color{poporange}+}}\par\vspace{2pt}
      {\raggedright\popsemi\fontsize{8}{10.5}\selectfont\color{popline}\MakeUppercase{Équipe pédagogique OnBuch+}\\\MakeUppercase{Conforme au programme officiel MINESEC}\par}
    \end{minipage}\hfill
    {\poplogo\fontsize{22}{22}\selectfont\color{popcream}OnBuch\color{poporange}+}
  \end{tcolorbox}%
}

% ============================================================
% Page de garde
% #1 titre · #2 matière · #3 niveau · #4 module · #5 n° leçon · #6 durée ·
% #7 description / objectifs (texte ou itemize)
% ============================================================
\newcommand{\pagegarde}[7]{%
  \thispagestyle{empty}
  \newgeometry{a4paper,margin=1.7cm,top=1.6cm,bottom=1.4cm}%
  \begin{tikzpicture}[remember picture,overlay]
    \fill[poporange!18] ([xshift=-3.2cm,yshift=-3.6cm]current page.north east) circle (4.6cm);
    \fill[poppurple!14] ([xshift=2.4cm,yshift=7.6cm]current page.south west) circle (3.8cm);
  \end{tikzpicture}%
  {\poplogo\fontsize{30}{30}\selectfont\color{popink}OnBuch\color{poporange}+}\hfill
  \raisebox{0.6ex}{\poppill{Cours signé \textbullet\ Premium}{popink}{popcream}}\par
  \vspace{1pt}\popsquiggle{poppurple}\par
  \vspace{8pt}
  \begin{tcolorbox}[enhanced,colback=poporange,colframe=popink,boxrule=2.8pt,arc=13pt,
      drop shadow=popink,left=18pt,right=18pt,top=12pt,bottom=13pt,after skip=10pt]
    \poppill{#2 \textbullet\ #3}{popink}{popcream}\hspace{5pt}\poppill{#4}{white}{popink}\par\vspace{8pt}
    {\popdisplay\fontsize{14}{14}\selectfont\color{popink}LEÇON #5}\par\vspace{4pt}
    {\raggedright\hyphenpenalty=10000\exhyphenpenalty=10000\popdisplay\fontsize{25}{27}\selectfont\color{popcream}\MakeUppercase{#1}\par}
    \vspace{8pt}
    {\popxbold\fontsize{9.5}{9.5}\selectfont\color{popink}\MakeUppercase{Durée conseillée : #6}}
  \end{tcolorbox}
  \begin{tcolorbox}[enhanced,colback=white,colframe=popink,boxrule=2.2pt,arc=10pt,
      drop shadow=popink,left=16pt,right=16pt,top=10pt,bottom=10pt,after skip=8pt]
    \poppill{Dans cette leçon}{poppurple}{white}\par\vspace{5pt}
    \popsemi\fontsize{9.3}{12.6}\selectfont\color{popink}#7
  \end{tcolorbox}
  \noindent\begin{minipage}[t]{0.487\linewidth}
    \begin{tcolorbox}[enhanced,colback=popgreen,colframe=popink,boxrule=2.2pt,arc=10pt,
        drop shadow=popink,left=11pt,right=11pt,top=10pt,bottom=10pt,height=3.1cm,valign=center]
      \tcbox[colback=white,colframe=popink,boxrule=1.3pt,arc=5pt,boxsep=0pt,nobeforeafter,
        left=3pt,right=3pt,top=3pt,bottom=3pt,valign=center]{\qrcode[height=1.9cm]{https://play.google.com/store/apps/details?id=com.luvvix.onbuch&hl=fr}}%
      \hspace{8pt}\begin{minipage}[c]{0.5\linewidth}
        {\popdisplay\fontsize{11}{12.5}\selectfont\color{white}\MakeUppercase{Télécharge OnBuch}}\par\vspace{4pt}
        {\raggedright\popsemi\fontsize{8.4}{11}\selectfont\color{white}Gratuit \textbullet\ Google Play\\Tuteur IA, annales, concours, résultats\par}
      \end{minipage}
    \end{tcolorbox}
  \end{minipage}\hfill
  \begin{minipage}[t]{0.487\linewidth}
    \begin{tcolorbox}[enhanced,colback=poppurple,colframe=popink,boxrule=2.2pt,arc=10pt,
        drop shadow=popink,left=11pt,right=11pt,top=10pt,bottom=10pt,height=3.1cm,valign=center]
      \tikz[baseline=-0.7ex]\node[circle,fill=white,draw=popink,line width=1.3pt,minimum size=1.6cm,inner sep=0]{\popdisplay\fontsize{24}{24}\selectfont\color{poppurple}+};%
      \hspace{8pt}\begin{minipage}[c]{0.6\linewidth}
        {\popdisplay\fontsize{11}{12.5}\selectfont\color{white}\MakeUppercase{Passe à OnBuch+}}\par\vspace{4pt}
        {\raggedright\popsemi\fontsize{8.4}{11}\selectfont\color{white}Tous les cours signés, Tuteur IA illimité, fiches PDF — onglet OnBuch+ de l'app\par}
      \end{minipage}
    \end{tcolorbox}
  \end{minipage}
  \vspace{8pt}
  \popcontenu
  \vfill
  \signatureonbuch
  \restoregeometry
  \newpage
}

% Rangée « Ce cours contient » (page de garde)
\newcommand{\poptile}[3]{% #1 couleur #2 gros texte #3 légende
  \begin{tcolorbox}[enhanced,colback=#1,colframe=popink,boxrule=2pt,arc=9pt,shadow={2.5pt}{-2.5pt}{0pt}{popink},
      left=6pt,right=6pt,top=9pt,bottom=9pt,halign=center,valign=center,height=1.75cm,width=\linewidth,nobeforeafter]
    {\popdisplay\fontsize{12.5}{13}\selectfont\color{white}\MakeUppercase{#2}}\par\vspace{3pt}
    {\centering\popsemi\fontsize{7.8}{9.5}\selectfont\color{white}#3\par}
  \end{tcolorbox}}
\newcommand{\popcontenu}{%
  \par\noindent{\popxboldls\fontsize{7.5}{7.5}\selectfont\color{popmuted}CE COURS SIGNÉ CONTIENT}\par\vspace{7pt}
  \noindent\begin{minipage}[t]{0.235\linewidth}\poptile{popblue}{Cours}{complet, illustré, pas à pas}\end{minipage}\hfill
  \begin{minipage}[t]{0.235\linewidth}\poptile{poporange}{Méthodes}{et exercices résolus}\end{minipage}\hfill
  \begin{minipage}[t]{0.235\linewidth}\poptile{poppink}{Exercices}{type examen + corrigés}\end{minipage}\hfill
  \begin{minipage}[t]{0.235\linewidth}\poptile{popgreen}{Fiche bilan}{l'essentiel à retenir}\end{minipage}\par}

% Sommaire
\newtcolorbox{sommaire}{enhanced,colback=white,colframe=popink,boxrule=2.2pt,arc=10pt,
  drop shadow=popink,left=18pt,right=18pt,top=13pt,bottom=13pt,before skip=6pt,after skip=20pt,
  before upper={\poppill{Sommaire}{poppurple}{white}\par\vspace{9pt}\popsemi\fontsize{10.6}{14.5}\selectfont\color{popink}}}

% Signature de fin de cours — poussée tout en bas de la dernière page
\newcommand{\findecours}{%
  \par\vfill\nopagebreak
  \begin{center}\popsquiggle{poporange}\end{center}\vspace{4pt}
  \signatureonbuch
}
\AtBeginDocument{\def\llbracket{\mathopen{[\mkern-3.5mu[}}\def\rrbracket{\mathclose{]\mkern-3.5mu]}}} % intervalles d'entiers (glyphe ⟦ absent de la police)
% Glyphes absents de Fira Math : redéfinis à partir de glyphes disponibles
\AtBeginDocument{%
  \def\setminus{\mathbin{\backslash}}\let\smallsetminus\setminus
  \def\vdots{\mathord{\vbox{\baselineskip3.2pt\lineskiplimit0pt\kern4pt\hbox{.}\hbox{.}\hbox{.}}}}%
  \def\ddots{\mathinner{\mkern1mu\raise6.5pt\hbox{.}\mkern2mu\raise3.6pt\hbox{.}\mkern2mu\raise0.7pt\hbox{.}\mkern1mu}}%
  \def\triangle{\mathord{\scalebox{0.9}{$\symup{\Delta}$}}}%
  \def\nabla{\mathord{\raisebox{\depth}{\scalebox{1}[-1]{$\symup{\Delta}$}}}}%
  \def\checkmark{\popcheck{popdark}}%
  \def\star{\mathbin{\ast}}\def\bigstar{\mathord{\ast}}%
  \def\diamond{\mathbin{\scalebox{0.6}{$\Diamond$}}}%
  \def\bigcup{\mathop{\mathchoice{\raisebox{-0.3em}{\scalebox{1.7}{$\cup$}}}{\scalebox{1.25}{$\cup$}}{\cup}{\cup}}\displaylimits}%
  \def\bigcap{\mathop{\mathchoice{\raisebox{-0.3em}{\scalebox{1.7}{$\cap$}}}{\scalebox{1.25}{$\cap$}}{\cap}{\cap}}\displaylimits}%
  \def\lesseqqgtr{\mathrel{\lessgtr}}\def\bigoplus{\mathop{\oplus}\displaylimits}%
}
\providecommand{\ssi}{si et seulement si} % abréviation fréquente
\AtBeginDocument{\providecommand{\wideparen}{\overparen}} % arc de cercle

%% FIGFIX-BEGIN v4 — correctifs globaux des figures (docs/onbuch-generation/tools/figfix)
\usepackage{adjustbox}\usepackage{etoolbox}
% toute figure TikZ/pgfplots plus large que la ligne est réduite à la largeur disponible
\BeforeBeginEnvironment{tikzpicture}{\begin{adjustbox}{max width=\linewidth}}
\AfterEndEnvironment{tikzpicture}{\end{adjustbox}}
% légendes pgfplots lisibles par-dessus les courbes
\pgfplotsset{every axis/.append style={legend style={fill=white,fill opacity=0.92,text opacity=1,draw=black!15,inner sep=3pt}}}
% étiquettes d'arêtes (syntaxe edge["..."]) avec fond blanc
\tikzset{every edge quotes/.append style={fill=white,inner sep=1.2pt}}
%% FIGFIX-END
\begin{document}

\pagegarde{Point de fonctionnement d'un circuit et loi de Pouillet}{Physique}{1ère D}{Module 4 : Aspects énergétiques des circuits électriques}{N15}{4 h}{Cette leçon introduit la notion de point de fonctionnement d'un circuit électrique et la loi de Pouillet, en s'appuyant sur les composants usuels (résistor, diode, transistor). Elle permet de déterminer graphiquement et analytiquement le point de fonctionnement et de calculer la puissance dans des circuits simples.
\vspace{4pt}
\begin{multicols}{2}\small
\begin{itemize}
\item Identifier les caractéristiques courant‑tension des résistors, diodes et transistors.
\item Définir le point de fonctionnement d'un circuit et expliquer son unicité.
\item Tracer la courbe de charge d'un circuit et déterminer graphiquement le point de fonctionnement.
\item Énoncer et démontrer la loi de Pouillet (P = U·I) pour un dipôle.
\item Calculer la puissance dissipée par chaque composant dans un circuit série.
\item Résoudre des problèmes de détermination du point de fonctionnement pour un circuit R‑D et un transistor en commutation.
\end{itemize}\end{multicols}}

\begin{sommaire}
\begin{multicols}{2}
\begin{enumerate}
\item Rappels sur les composants de base
\item Caractéristiques courant‑tension I(U)
\item Point de fonctionnement : définition et méthode graphique
\item Loi de Pouillet : puissance et énergie dans un dipôle
\item Détermination du point de fonctionnement dans des circuits simples
\item Synthèse et préparation au Probatoire
\item Activité d'intégration
\item Méthodes et exercices résolus
\item Exercices
\item Corrigés des exercices
\item Fiche bilan
\end{enumerate}
\end{multicols}
\end{sommaire}

\begin{objectifs}
\begin{itemize}
\item Identifier les caractéristiques courant‑tension des résistors, diodes et transistors.
\item Définir le point de fonctionnement d'un circuit et expliquer son unicité.
\item Tracer la courbe de charge d'un circuit et déterminer graphiquement le point de fonctionnement.
\item Énoncer et démontrer la loi de Pouillet (P = U·I) pour un dipôle.
\item Calculer la puissance dissipée par chaque composant dans un circuit série.
\item Résoudre des problèmes de détermination du point de fonctionnement pour un circuit R‑D et un transistor en commutation.
\item Analyser les erreurs fréquentes (polarité diode, saturation transistor) et les éviter.
\item Appliquer les acquis à des exercices type Probatoire.
\end{itemize}
\end{objectifs}

\begin{prerequis}
\begin{itemize}
\item Loi d'Ohm (U = R·I) et notion de résistance équivalente.
\item Notion de tension, courant, puissance et énergie électrique (Seconde).
\item Représentation graphique de fonctions (courbes I(U)).
\item Symboles normalisés des composants électroniques de base.
\item Manipulation simple d'équations littérales.
\end{itemize}
\end{prerequis}

\newpage

\coursec{Rappels sur les composants de base}

Dans un quartier populaire de Yaoundé, un jeune bricoleur passionné d'électronique récupère une diode électroluminescente (LED) rouge sur un vieux chargeur de téléphone, une pile neuve de \SI{9}{V} et une résistance carbone dont il ignore la valeur. Il veut allumer cette LED pour éclairer sa table de révision le soir, sans la griller. Il sait que la LED a une tension de seuil autour de \SI{1,8}{V} et qu'elle supporte un courant maximal de \SI{20}{mA}. Quelle résistance doit-il choisir ? Combien de temps la pile tiendra-t-elle ? Cette situation concrète, que beaucoup d'entre vous ont vécue en bricolant un circuit pour une veilleuse ou un montage scolaire, illustre parfaitement le cœur de ce module : \textbf{connaître les composants, leurs caractéristiques courant-tension et savoir déterminer le point de fonctionnement d'un circuit pour en maîtriser l'énergie (loi de Pouillet)}. Avant d'aborder la méthode graphique du point de fonctionnement et le calcul de puissance, rappelons rigoureusement les trois acteurs principaux de nos circuits : le résistor, la diode et le transistor bipolaire.

\courssub{Le résistor : dipôle ohmique linéaire et effet Joule}

Le résistor (ou résistance) est le composant passif le plus fondamental. Il est constitué d'un matériau conducteur (carbone, film métallique, fil résistif) qui s'oppose au passage du courant électrique. Cette opposition se traduit par une dissipation d'énergie sous forme de chaleur : c'est l'\cle{effet Joule}.

\begin{definition}[Résistor – Dipôle ohmique linéaire]
Un \textbf{résistor} est un dipôle \textbf{passif} dont la tension $u$ aux bornes et le courant $i$ qui le traverse (convention récepteur) sont liés par la \textbf{loi d'Ohm} :
\[
u = R \, i
\]
où $R$ est la \textbf{résistance} (constante pour un résistor ohmique à température constante), exprimée en \textbf{ohms} (\si{\ohm}).
La caractéristique $i(u)$ est une droite passant par l'origine de pente $1/R$.
\end{definition}

\begin{popfigure}
\begin{tikzpicture}[scale=1.2, european]
% Resistor symbol
\draw (0,0) to[R=$R$, l_={$u$}, i={$i$}] (3,0);
% Labels
\node[above] at (1.5,0.8) {\textbf{Symbole normalisé (CEI 60617)}};
\node[below] at (1.5,-0.6) {Convention récepteur : $\vec{i}$ entre par la borne $+$ de $u$};
\end{tikzpicture}
\legende{Symbole du résistor, convention récepteur et loi d'Ohm $u = R\,i$.}
\end{popfigure}

\noindent\textbf{Modélisation énergétique (Bilan de puissance).}
La puissance instantanée absorbée par le résistor (effet Joule) s'écrit :
\[
p(t) = u(t)\,i(t) = R\,i^2(t) = \frac{u^2(t)}{R} \quad (\text{en watts, \si{W}})
\]
L'énergie dissipée en chaleur sur une durée $\Delta t$ est :
\[
W = \int_{t_1}^{t_2} p(t)\,dt \quad (\text{en joules, \si{J}})
\]
\begin{attention}[Piège : Convention récepteur vs générateur]
Au Probatoire, la convention \textbf{récepteur} est la norme pour les dipôles passifs (résistor, diode, lampe) : le courant $i$ entre par la borne de potentiel le plus élevé (borne $+$ de $u$). Si tu inverses la flèche de $i$ par rapport à $u$, la loi d'Ohm s'écrit $u = -R\,i$ et la puissance $p = u\,i$ devient \emph{négative} (le dipôle fournirait de l'énergie), ce qui est physiquement faux pour un résistor passif. \textbf{Toujours dessiner la flèche de $i$ entrant par la borne $+$ de $u$ sur tes schémas.}
\end{attention}

\courssub{La diode : dipôle passif non linéaire à conduction unidirectionnelle}

La diode est un composant \textbf{passif non linéaire} à base de jonction $PN$ (silicium ou germanium). Elle ne conduit le courant que dans un seul sens : le sens \textbf{direct} (anode $A$ vers cathode $K$). Elle ne fournit pas d'énergie au circuit (pas d'amplification), elle la dissipe ou la convertit (lumière pour une LED).

\begin{definition}[Diode – Caractéristique idéale et réelle]
Une \textbf{diode} idéale est un interrupteur commandé par la tension :
\begin{itemize}
\item Si $u_{AK} < 0$ (bloquée) : $i = 0$ (circuit ouvert).
\item Si $u_{AK} > 0$ (conduite) : $u_{AK} = 0$ (court-circuit idéal).
\end{itemize}
Une diode \textbf{réelle} (silicium) présente une \textbf{tension de seuil} $V_{\gamma} \approx \SI{0,7}{V}$ (germanium : \SI{0,3}{V}) :
\begin{itemize}
\item \textbf{Régime bloqué} ($u_{AK} < V_{\gamma}$) : courant quasi nul ($i \approx 0$), résistance très grande.
\item \textbf{Régime passant} ($u_{AK} \ge V_{\gamma}$) : la diode conduit, $u_{AK} \approx V_{\gamma}$ (quasi constante), le courant est limité par le circuit externe.
\end{itemize}
\end{definition}

\begin{popfigure}
\begin{tikzpicture}[scale=1.1, european]
% Diode symbol
\draw (0,0) to[D=$D$, l_={$u_{AK}$}, i={$i$}] (3,0);
% Anode/Cathode labels
\node[above left] at (0,0) {Anode ($A$)};
\node[above right] at (3,0) {Cathode ($K$)};
% Bar on cathode side
\draw[thick] (2.85,-0.2) -- (2.85,0.2);
\node[below] at (1.5,-0.8) {Sens direct : $A \to K$ ($i>0$, $u_{AK}>0$)};
\end{tikzpicture}
\legende{Symbole normalisé de la diode. La barre indique la cathode ($K$). Convention récepteur : $i$ entre par $A$, $u_{AK} = V_A - V_K$.}
\end{popfigure}

\begin{savaistu}[La diode au Cameroun : du poste de radio à la fibre optique]
Les premières radios à galène (années 1930-50) utilisaient une pointe de métal sur un cristal de galène (sulfure de plomb) comme diode de détection : c'était une jonction $PN$ rudimentaire ! Aujourd'hui, les diodes laser (InGaAsP) sont au cœur des liaisons à fibre optique qui relient Yaoundé, Douala et les capitales sous-régionales via le réseau \textbf{CAMTEL} et le câble sous-marin \textbf{WACS}. Les LED (diodes électroluminescentes) équipent l'éclairage public solaire des routes rurales (programme \textbf{SONATREL}/Banque Mondiale) et les écrans des téléphones \textbf{MTN}/\textbf{Orange}.
\end{savaistu}

\begin{attention}[Erreur fatale : Inversion anode/cathode]
Sur le schéma, le triangle pointe vers la barre. Le courant conventionnel ne peut circuler que \textbf{dans le sens de la flèche du triangle} (Anode $\to$ Cathode). Si tu inverses la diode dans un circuit (cathode vers le $+$ de la source), elle est \textbf{bloquée} : $i=0$, la LED ne s'allume pas, le transistor ne commande rien. \textbf{Vérifie toujours la polarité avant de souder ou de brancher !} Au Probatoire, une diode inversée dans un exercice de point de fonctionnement donne un circuit ouvert ($I=0$).
\end{attention}

\courssub{Le transistor bipolaire (BJT) : composant actif à trois bornes, amplificateur et interrupteur}

Le transistor bipolaire à jonction (BJT) est le composant \textbf{actif} roi de l'électronique analogique et de commutation. Il existe en deux polarités : \textbf{NPN} (le plus courant) et \textbf{PNP}. Il possède trois bornes : \textbf{Base (B)}, \textbf{Collecteur (C)}, \textbf{Émetteur (E)}.

\begin{definition}[Transistor bipolaire NPN – Structure et symbole]
Un transistor \textbf{NPN} est constitué de trois zones dopées : $N$ (collecteur), $P$ (base, très fine), $N$ (émetteur). Ses deux jonctions $PN$ sont :
\begin{itemize}
\item Jonction Base-Émetteur (BE) : diode $PN$ (seuil $V_{\gamma} \approx \SI{0,7}{V}$).
\item Jonction Base-Collecteur (BC) : diode $PN$ (généralement polarisée inverse en régime actif).
\end{itemize}
Le symbole normalisé montre une flèche sur l'émetteur \textbf{sortant} du boîtier pour le NPN (rentrant pour le PNP).
\end{definition}

\begin{popfigure}
\begin{tikzpicture}[scale=1.2, european]
% NPN Transistor
\draw (0,0) node[npn](npn){};
% Labels
\node[left] at (npn.B) {Base (B)};
\node[above right] at (npn.C) {Collecteur (C)};
\node[below right] at (npn.E) {Émetteur (E)};
% Arrow direction reminder
\draw[popblue, thick, ->] (npn.E) ++(0.3,0) -- ++(0.5,0) node[midway, above] {Flèche = sens courant conventionnel (NPN)};
\node[above] at (1.5,1.5) {\textbf{Transistor NPN (symbole CEI)}};
\end{tikzpicture}
\legende{Symbole du transistor bipolaire NPN. La flèche sur l'émetteur indique le sens du courant conventionnel dans la jonction BE (P$\to$N).}
\end{popfigure}

\noindent\textbf{Physique du fonctionnement (Modèle simplifié pour la Première).}
Le transistor est un \textbf{dispositif commandé en courant} : un faible courant de base $i_B$ contrôle un fort courant de collecteur $i_C$.
\begin{itemize}
\item Courant d'émetteur : $i_E = i_B + i_C$ (loi des nœuds).
\item \textbf{Gain en courant commun-émetteur} : $\beta = \dfrac{i_C}{i_B}$ (typiquement $50 \dots 300$ pour un transistor de petit signal type BC547, 2N2222).
\end{itemize}

\begin{propriete}[Modèle simplifié du transistor en commutation (Interrupteur commandé)]
Pour la détermination du point de fonctionnement en régime \textbf{tout-ou-rien} (logique, pilotage de LED, relais, moteur), on utilise le modèle \textbf{interrupteur commandé} à trois régimes :
\begin{enumerate}
\item \textbf{Bloqué (Cut-off)} : $i_B = 0 \Rightarrow i_C = 0$. Le transistor est un \textbf{circuit ouvert} entre C et E ($R_{CE} \to \infty$). Tension $u_{CE} = V_{CC}$ (tension d'alimentation).
\item \textbf{Actif (Linéaire)} : $i_B > 0$ et $i_C = \beta\,i_B$ (limité par le circuit externe). $u_{CE} > V_{CE(sat)} \approx \SI{0,2}{V}$. Le transistor agit comme une \textbf{source de courant commandée}.
\item \textbf{Saturé (Saturation)} : $i_B$ assez grand pour que le circuit externe ne puisse plus fournir $i_C = \beta\,i_B$. Le transistor est un \textbf{interrupteur fermé} : $u_{CE} \approx V_{CE(sat)} \approx \SI{0,2}{V}$ (quasi court-circuit). Condition : $i_B > \dfrac{i_{C(sat)}}{\beta}$ (sursaturation).
\end{enumerate}
\end{propriete}

\begin{aretenir}[Classification rigoureuse : Linéaire / Non-linéaire — Passif / Actif]
\begin{center}
\begin{tabularx}{\linewidth}{|l|X|X|X|}\hline
\textbf{Propriété} & \textbf{Résistor (Linéaire, Passif)} & \textbf{Diode (Non-linéaire, Passive)} & \textbf{Transistor (Non-linéaire, Actif)} \\ \hline
Énergie & Dissipe (effet Joule). $p = u\,i \ge 0$. & Dissipe (ou convertit en lumière). $p = u\,i \ge 0$. & \textbf{Amplifie} : puissance de sortie $>$ puissance de commande (alimentation externe requise). \\ \hline
Caractéristique $i(u)$ & \textbf{Linéaire} (droite passant par l'origine). & \textbf{Non linéaire} (seuil $V_{\gamma}$, conduction unidirectionnelle). & \textbf{Non linéaire} (famille de courbes $i_C(u_{CE})$ paramétrée par $i_B$). \\ \hline
Symétrie & Symétrique : $u \to -u \Rightarrow i \to -i$. & \textbf{Asymétrique} : conduction unidirectionnelle. & \textbf{Asymétrique} : polarités fixes (NPN/PNP). \\ \hline
Exemples & Résistance carbone, CMS, fil chauffant, LDR, thermistance. & Diode $PN$, LED, Zener, diode Schottky. & Transistor BJT, MOSFET, IGBT, Circuit intégré. \\ \hline
\end{tabularx}
\end{center}
\end{aretenir}

\courssub{Exemple résolu : Identification et polarisation d'un montage simple}

\begin{exoresolu}[Analyse d'un circuit de pilotage de LED par transistor NPN]
\textbf{Énoncé :} On considère le montage ci-dessous alimenté par $V_{CC} = \SI{9}{V}$. La LED rouge a une tension de seuil $V_{\gamma,LED} = \SI{1,8}{V}$ pour un courant nominal $I_{LED,nom} = \SI{10}{mA}$. Le transistor NPN (BC547) a $\beta = 200$ et $V_{CE(sat)} = \SI{0,2}{V}$, $V_{BE(on)} = \SI{0,7}{V}$. La résistance de base est $R_B = \SI{10}{k\ohm}$, la résistance de collecteur (pour la LED) est $R_C = \SI{470}{\ohm}$. L'interrupteur $K$ est fermé.
\begin{enumerate}
\item Redessiner le schéma en remplaçant chaque composant par son modèle équivalent (générateur de tension, résistance, interrupteur, diode).
\item Déterminer les courants $i_B$, $i_C$ et la tension $u_{CE}$. Le transistor est-il saturé ?
\item Calculer la puissance dissipée par $R_C$ et par le transistor. Vérifier le bilan de puissance.
\end{enumerate}

\begin{tikzpicture}[scale=0.9, european, every node/.style={font=\small}]
% Power supply
\coordinate (batneg) at (0,0);
\coordinate (batpos) at (0,1.6);
\draw (batneg) to[battery1={$V_{CC}=9V$}] (batpos);
% Switch
\draw (batpos) -- ++(1,0) to[nos=$K$] ++(1.5,0) coordinate (n1);
% Base resistor
\draw (n1) to[R={$R_B=10k\Omega$}, *-*] ++(0, -2) coordinate (nB);
% Transistor
\draw (nB) node[npn, anchor=B] (Q) {};
% Collector resistor + LED
\draw (Q.C) to[R={$R_C=470\Omega$}, *-] ++(0, 1.5) coordinate (nCtop);
\draw (nCtop) to[D=$LED$, *-] ++(0, 1) coordinate (nLEDtop);
\draw (nLEDtop) -- (nLEDtop -| batpos) -- (batpos);
% Emitter to ground
\draw (Q.E) -- ++(0, -0.5) node[ground] {};
% Base ground reference (implicit via Rb)
\draw (nB) -- ++(-0.5,0) node[left] {B};
\draw (Q.C) -- ++(0.5,0) node[right] {C};
\draw (Q.E) -- ++(0.5,0) node[right] {E};
\end{tikzpicture}

\tcblower
\textbf{Solution détaillée :}

\textbf{1. Modèles équivalents (Interrupteur $K$ fermé = fil) :}
\begin{itemize}
\item $V_{CC}$ : Générateur de tension idéal $E = \SI{9}{V}$.
\item $R_B, R_C$ : Résistances ohmiques.
\item Jonction BE (diode) : Générateur de tension $V_{BE} = \SI{0,7}{V}$ (seuil) + interrupteur fermé (si $i_B>0$).
\item LED (diode) : Générateur de tension $V_{\gamma,LED} = \SI{1,8}{V}$ + interrupteur fermé (si $i_C>0$).
\item Transistor (jonction BC) : En régime actif/saturé, polarisée inverse $\to$ circuit ouvert (courant de fuite négligé).
\item Transistor (commande $i_C$) : Source de courant commandée $i_C = \beta i_B$ (actif) ou Interrupteur fermé $u_{CE}=V_{CE(sat)}$ (saturé).
\end{itemize}

\textbf{2. Calcul des courants (Hypothèse : Transistor saturé $\to$ vérifier a posteriori).}
\textit{Maille Base-Émetteur :}
\[
V_{CC} = R_B\,i_B + V_{BE} \quad \Rightarrow \quad i_B = \frac{V_{CC} - V_{BE}}{R_B} = \frac{9 - 0,7}{10\,000} = \SI{0,83}{mA}
\]
\textit{Courant collecteur maximal imposé par le circuit externe (Maille Collecteur) :}
\[
V_{CC} = R_C\,i_C + V_{\gamma,LED} + u_{CE}
\]
Si saturé, $u_{CE} = V_{CE(sat)} = \SI{0,2}{V}$.
\[
i_{C(sat)} = \frac{V_{CC} - V_{\gamma,LED} - V_{CE(sat)}}{R_C} = \frac{9 - 1,8 - 0,2}{470} = \frac{7}{470} \approx \SI{14,9}{mA}
\]
\textit{Vérification de la saturation :}
Courant collecteur demandé par la base : $i_{C(\beta)} = \beta\,i_B = 200 \times 0,83\,\text{mA} = \SI{166}{mA}$.
Comme $i_{C(\beta)} \gg i_{C(sat)}$, la base fournit bien plus de courant que nécessaire. \textbf{Le transistor est bien SATURÉ.}
Donc : $i_C = i_{C(sat)} = \SI{14,9}{mA}$, $u_{CE} = \SI{0,2}{V}$.
\begin{attention}[Courant LED vs Nominal]
Le courant dans la LED (\SI{14,9}{mA}) dépasse le nominal (\SI{10}{mA}). Pour une durée de vie optimale, il faudrait augmenter $R_C$ à environ \SI{680}{\ohm} ($R_C = \frac{7}{0,01} = 700\,\Omega$). Ici, la LED brillera plus fort mais s'usera plus vite.
\end{attention}

\textbf{3. Puissances dissipées (Loi de Pouillet - anticipation Section 4) :}
\begin{align*}
P_{R_C} &= R_C \cdot i_C^2 = 470 \times (\num{0,0149})^2 \approx \SI{0,104}{W} = \SI{104}{mW} \quad (\text{Résistance } \frac{1}{4}W \text{ OK}) \\
P_{Transistor} &= u_{CE} \cdot i_C = 0,2 \times \num{0,0149} \approx \SI{3,0}{mW} \quad (\text{Très faible, pas de radiateur nécessaire}) \\
P_{LED} &= V_{\gamma,LED} \cdot i_C = 1,8 \times \num{0,0149} \approx \SI{26,8}{mW} \quad (\text{Lumière + chaleur}) \\
P_{R_B} &= R_B \cdot i_B^2 = 10000 \times (\num{0,00083})^2 \approx \SI{6,9}{\mu W} \quad (\text{Négligeable})
\end{align*}
\textbf{Bilan :} Puissance fournie par la pile $P_{bat} = V_{CC} \cdot (i_B + i_C) \approx 9 \times \num{0,01573} \approx \SI{141,6}{mW}$.
Somme des puissances absorbées $\approx 104 + 3 + 26,8 + 0,007 \approx \SI{133,8}{mW}$ (différence due aux arrondis sur $i_C$). \textbf{Le bilan de puissance est respecté.}
\end{exoresolu}

\courssub{Synthèse des symboles et conventions pour la suite}

\begin{center}
\begin{tabularx}{\linewidth}{|l|X|X|X|}\hline
\textbf{Composant} & \textbf{Symbole (CEI)} & \textbf{Grandeurs (Conv. Récepteur)} & \textbf{Modèle simplifié (Cours 1ère C)} \\ \hline
Résistor $R$ & \tikz[baseline=-0.5ex, european]{\draw (0,0) to[R=$R$] (1,0);} & $u = R\,i$ & Résistance constante $R$ \\ \hline
Diode $D$ (Si) & \tikz[baseline=-0.5ex, european]{\draw (0,0) to[D=$D$] (1,0);} & $u_{AK}, i$ & Seuil $V_{\gamma} \approx 0,7\,V$ ; $R_{on} \approx 0$ \\ \hline
Transistor NPN & \tikz[baseline=-0.5ex, european]{\draw (0,0) node[npn, scale=0.7] (q) {};} & $i_B, i_C, i_E$ ; $u_{BE}, u_{CE}$ & Interrupteur commandé : Bloqué / Actif ($\beta$) / Saturé ($V_{CE(sat)}$) \\ \hline
\end{tabularx}
\end{center}

\begin{methode}[Bon réflexes pour la suite (Point de fonctionnement)]
\begin{enumerate}
\item \textbf{Identifier} chaque composant et son modèle (linéaire / non-linéaire, passif / actif).
\item \textbf{Dessiner} le schéma complet avec \textbf{convention récepteur} sur chaque dipôle passif (flèche $i$ entrant par $+$ de $u$).
\item \textbf{Écrire} les équations caractéristiques de chaque composant ($u=Ri$, $u=V_{\gamma}$, $i_C=\beta i_B$ ou $u_{CE}=V_{CE(sat)}$).
\item \textbf{Appliquer} les lois de mailles et de nœuds (KVL, KCL).
\item \textbf{Résoudre} le système (souvent graphique pour les non-linéaires, voir Section 3).
\item \textbf{Vérifier} la cohérence (régime du transistor, sens des courants, puissances positives, bilan de puissance).
\end{enumerate}
\end{methode}

\noindent Nous avons maintenant tous les « briques » nécessaires. Dans la prochaine section, nous tracerons les caractéristiques $I(U)$ de ces composants et apprendrons la méthode graphique universelle pour trouver le point de fonctionnement d'un circuit : l'intersection de la caractéristique du composant et de la droite de charge du circuit.

\coursec{Caractéristiques courant‑tension I(U)}

Dans la section précédente, nous avons rappelé la constitution et le symbole normalisé de trois composants fondamentaux : le résistor, la diode et le transistor bipolaire. Nous savons maintenant \emph{ce qu'ils sont} physiquement. Pour comprendre \emph{comment ils se comportent} dans un circuit, il nous faut une relation mathématique ou graphique liant la tension $u$ à leurs bornes au courant $i$ qui les traverse. C'est l'objet de la \cle{caractéristique courant‑tension}, notée $i = f(u)$ ou $I(U)$ en régime continu. C'est la « carte d'identité électrique » de tout dipôle.

\courssub{Caractéristique du résistor : la loi d'Ohm}

Commençons par le composant le plus simple. Un résistor est un dipôle \cle{ohmique} : le courant qui le traverse est proportionnel à la tension appliquée.

\begin{experience}[Caractéristique $I(U)$ d'un résistor]
\textbf{Matériel :} Générateur de tension variable $E$ (0--10~V), résistor $R = \SI{100}{\ohm}$ (puissance \SI{1}{\watt}), ampèremètre, voltmètre, fils de connexion.\par
\textbf{Montage :} Circuit série simple. L'ampèremètre est en série avec $R$, le voltmètre en parallèle aux bornes de $R$.\par
\textbf{Protocole :} Faire varier $E$ par paliers de \SI{1}{\volt}. Relever $U$ (voltmètre) et $I$ (ampèremètre). Inverser la polarité de $E$ pour obtenir des tensions négatives.\par
\textbf{Observations :} Les points $(U, I)$ s'alignent sur une droite passant par l'origine. La pente de cette droite est constante.\par
\textbf{Interprétation :} Le résistor est un dipôle linéaire, passif et réversible. La relation est $U = R I$.
\end{experience}

\begin{popfigure}
\begin{tikzpicture}
\begin{axis}[
    width=\linewidth, height=7cm,
    axis lines=middle,
    xlabel={$U$ (V)}, ylabel={$I$ (mA)},
    xmin=-5, xmax=5, ymin=-50, ymax=50,
    xtick={-5,-2.5,0,2.5,5}, ytick={-50,-25,0,25,50},
    grid=major, grid style={popline},
    tick label style={font=\small},
    label style={font=\small},
]
\addplot[popcurve, thick, domain=-5:5, samples=2] {x*10}; % R=100 ohm -> I(mA) = U(V)/0.1 = 10*U
\end{axis}
\end{tikzpicture}
\legende{Caractéristique $I(U)$ d'un résistor $R = \SI{100}{\ohm}$. Droite passant par l'origine, pente $1/R = \SI{10}{\milli\ampere\per\volt}$.}
\end{popfigure}

\begin{propriete}[Loi d'Ohm — Dipôle ohmique]
Pour un résistor de résistance $R$ (en ohms $\Omega$), la tension $u$ à ses bornes (référence récepteur) et le courant $i$ qui le traverse sont liées par :
\[
u = R \, i \qquad \text{ou} \qquad i = \frac{1}{R}\, u = G\, u
\]
où $G = 1/R$ est la \cle{conductance} (en siemens~S).
\emph{Complément microscopique (non exigible au Probatoire) :} Dans un conducteur métallique, les électrons libres subissent un champ électrique $\vec{E} = \vec{u}/L$. La force électrique $\vec{F}_e = q\vec{E}$ les accélère, mais les collisions avec le réseau cristallin (phonons) créent une force de frottement proportionnelle à la vitesse de dérive $\vec{v}_d$ : $\vec{F}_f = -m\vec{v}_d/\tau$. À l'équilibre dynamique, $qE = m v_d/\tau$. Le courant est $i = n q S v_d$. En remplaçant $v_d$, on trouve $i = \frac{n q^2 \tau}{m} \frac{S}{L} u$. Le facteur $\frac{n q^2 \tau}{m}$ est la conductivité $\sigma$, et $\frac{S}{L}$ est la géométrie. Donc $i = \sigma \frac{S}{L} u = \frac{1}{\rho}\frac{S}{L} u = \frac{1}{R} u$.
\end{propriete}

\begin{aretenir}[Résistance statique vs dynamique]
Pour un dipôle ohmique, la \cle{résistance statique} $R = U/I$ est constante et égale à la \cle{résistance dynamique} $r_d = \dfrac{dU}{dI} = R$. La courbe $I(U)$ est une droite : la pente est constante partout.
\end{aretenir}

\courssub{Caractéristique de la diode : modèle de Shockley}

La diode est un dipôle \cle{non linéaire}, \cle{non réversible} (elle conduit dans un seul sens) et \cle{passif}. Sa caractéristique est décrite par la célèbre équation de \cle{Shockley} (1949), issue de la physique des jonctions $pn$.

\begin{experience}[Caractéristique $I(U)$ d'une diode 1N4148]
\textbf{Matériel :} Même montage que pour le résistor, mais on remplace $R$ par une diode 1N4148 (anode vers $+E$). On ajoute une résistance de protection $R_p = \SI{1}{\kilo\ohm}$ en série pour limiter le courant direct (la diode est quasi court-circuit pour $U > 0,7$~V).\par
\textbf{Protocole :} Balayer $E$ de $-\SI{10}{\volt}$ à $+\SI{2}{\volt}$. Noter qu'en inverse, le courant est quasi nul (nA). En direct, le courant « démarre » brutalement vers \SI{0,6}{\volt}--\SI{0,7}{\volt}.\par
\textbf{Observation :} Courbe exponentielle en direct, quasi plateau (courant de fuite $I_0$) en inverse.
\end{experience}

\begin{propriete}[Théorème — Caractéristique de Shockley (admis au Probatoire)]
Pour une jonction $pn$ idéale à la température $T$, le courant $i$ en fonction de la tension $u$ (référence récepteur) est :
\[
i(u) = I_0 \left( e^{\frac{q u}{n k T}} - 1 \right)
\]
où :
\begin{itemize}
    \item $I_0$ : \cle{courant de saturation inverse} (A), typiquement $10^{-12}$ à $10^{-9}$~A pour une diode Si. Dépend fortement de $T$ ($\propto T^3 e^{-E_g/kT}$).
    \item $q = \SI{1,602e-19}{\coulomb}$ : charge de l'électron.
    \item $k = \SI{1,381e-23}{\joule\per\kelvin}$ : constante de Boltzmann.
    \item $T$ : température absolue (K).
    \item $n$ : \cle{coefficient d'idéalité} ($1 \le n \le 2$). $n=1$ pour jonction idéale (diffusion), $n \approx 2$ si recombinaison dans la zone de déplétion domine. Pour 1N4148, $n \approx 1,9$.
\end{itemize}
La \cle{tension thermique} $V_T = \dfrac{kT}{q}$ vaut $\approx \SI{25,85}{\milli\volt}$ à $T = \SI{300}{\kelvin}$ ($27^\circ$C).
\end{propriete}

\begin{attention}[Référence récepteur obligatoire !]
L'équation de Shockley s'écrit \textbf{uniquement en convention récepteur} ($u = V_A - V_K$, $i$ entrant par l'anode). Si tu utilises la convention générateur ($u = V_K - V_A$), le signe change : $i = I_0(1 - e^{-qu/nkT})$. Au Probatoire, \textbf{précise toujours la convention} et dessine la flèche de $u$ et $i$ sur le schéma.
\end{attention}

\begin{popfigure}
\begin{tikzpicture}
\begin{axis}[
    width=\linewidth, height=8cm,
    axis lines=middle,
    xlabel={$U$ (V)}, ylabel={$I$ (mA)},
    xmin=-1.5, xmax=1.0, ymin=-0.5, ymax=15,
    xtick={-1.5,-1,-0.5,0,0.5,0.7,1.0}, ytick={0,5,10,15},
    grid=major, grid style={popline},
    tick label style={font=\small},
    label style={font=\small},
    clip=false,
]
% Shockley: I0=2nA=2e-6 mA, n=1.9, VT=0.02585 -> q/nkT = 1/(n*VT) = 1/(1.9*0.02585) = 20.33
\addplot[popcurve, very thick, domain=-1.5:0.85, samples=200] 
    { (2e-6)*(exp(x/(1.9*0.02585)) - 1) }; % I in mA
% Asymptote inverse
\addplot[popmuted, dashed, domain=-1.5:0] { -2e-6 }; 
\node[popdark, anchor=north east, font=\tiny] at (axis cs: -1.4, -0.002) {$I \approx -I_0$};
% Seuil visuel
\draw[poporange, dashed, thick] (axis cs:0.7,0) -- (axis cs:0.7,14) node[above, font=\small, poporange] {$U_\gamma \approx 0,7$~V};
\end{axis}
\end{tikzpicture}
\legende{Caractéristique $I(U)$ d'une diode 1N4148 ($I_0 = \SI{2}{\nano\ampere}$, $n=1,9$, $T=\SI{300}{\kelvin}$). Échelle linéaire : l'exponentielle paraît verticale après le seuil.}
\end{popfigure}

\begin{popfigure}
\begin{tikzpicture}
\begin{semilogyaxis}[
    width=\linewidth, height=7cm,
    axis lines=middle,
    xlabel={$U$ (V)}, ylabel={$I$ (A) -- échelle log},
    xmin=-1, xmax=0.85, ymin=1e-12, ymax=1e-1,
    xtick={-1,-0.5,0,0.2,0.4,0.6,0.7,0.8}, ytick={1e-12,1e-10,1e-8,1e-6,1e-4,1e-2,1e-1},
    grid=major, grid style={popline},
    tick label style={font=\small},
    label style={font=\small},
]
\addplot[popcurve, very thick, domain=-1:0.85, samples=200] 
    { 2e-9*(exp(x/(1.9*0.02585)) - 1) };
% Linéaire en direct => droite en log
\draw[poppurple, dashed] (axis cs:0,1e-9) -- (axis cs:0.85, 0.1) node[right, font=\tiny, poppurple] {Pente $\frac{q}{nkT}$};
\end{semilogyaxis}
\end{tikzpicture}
\legende{Même diode en échelle semi-logarithmique. Le régime direct devient une droite (pente $q/nkT$), le régime inverse montre le plateau $-I_0$.}
\end{popfigure}

\textbf{Zoom sur la région de seuil (0,6--0,8 V).}
C'est la zone cruciale pour le point de fonctionnement. La notion de « tension de seuil » $U_\gamma$ est une \emph{approximation ingénieur} : on dit que la diode « conduit » si $U > U_\gamma$ ($I$ devient significatif, ex. $> \SI{1}{\milli\ampere}$). Pour du Si, $U_\gamma \in [0,6 ; 0,8]$~V selon le courant visé.

\begin{popfigure}
\begin{tikzpicture}
\begin{axis}[
    width=\linewidth, height=7cm,
    axis lines=middle,
    xlabel={$U$ (V)}, ylabel={$I$ (mA)},
    xmin=0.55, xmax=0.85, ymin=0, ymax=5,
    xtick={0.55,0.6,0.65,0.7,0.75,0.8,0.85}, ytick={0,1,2,3,4,5},
    grid=major, grid style={popline},
    tick label style={font=\small},
    label style={font=\small},
]
\addplot[popcurve, very thick, domain=0.55:0.85, samples=100] 
    { (2e-6)*(exp(x/(1.9*0.02585)) - 1) };
% Tangente en 0.7V pour résistance dynamique (I=3.08mA, rd=15.8 ohm -> slope 63.3 mA/V)
\addplot[poporange, dashed, domain=0.55:0.85] { 63.3*(x-0.7) + 3.08 }; 
\node[poporange, anchor=south west, font=\small] at (axis cs:0.56, 0.5) {Tangente en $U=0,7$~V : $r_d \approx 15,8~\Omega$};
% Points repères
\addplot[only marks, mark=*, popblue] coordinates {(0.6, 0.4) (0.7, 3.08)};
\node[popblue, anchor=south] at (axis cs:0.6, 0.4) {$I \approx 0,4$~mA};
\node[popblue, anchor=south] at (axis cs:0.7, 3.08) {$I \approx 3,1$~mA};
\end{axis}
\end{tikzpicture}
\legende{Zoom sur le « genoux » de la diode 1N4148 ($I_0=\SI{2}{\nano\ampere}, n=1,9$). La pente (conductance dynamique) augmente exponentiellement. La résistance dynamique $r_d = dU/dI$ chute de $\sim \SI{120}{\ohm}$ à 0,6~V à $\sim \SI{16}{\ohm}$ à 0,7~V.}
\end{popfigure}

\begin{exoresolu}[Calcul du courant direct d'une 1N4148]
\textbf{Énoncé :} Une diode 1N4148 a pour paramètres $I_0 = \SI{2,0}{\nano\ampere}$ et $n = 1,9$. La température est $T = \SI{300}{\kelvin}$ ($V_T = \SI{25,85}{\milli\volt}$). Calculer le courant $I$ pour une tension directe $U = \SI{0,70}{\volt}$.\par
\textbf{Solution détaillée :}
On utilise la loi de Shockley en régime direct ($U \gg nV_T$, le « $-1$ » est négligeable) :
\[
I \approx I_0 \, e^{\frac{U}{n V_T}}
\]
Calcul de l'exposant :
\[
\frac{U}{n V_T} = \frac{0,70}{1,9 \times 0,02585} = \frac{0,70}{0,049115} \approx 14,25
\]
Calcul de l'exponentielle :
\[
e^{14,25} \approx 1,54 \times 10^6
\]
Courant :
\[
I \approx 2,0 \times 10^{-9} \times 1,54 \times 10^6 = 3,08 \times 10^{-3}~\text{A} = \boxed{3,08~\text{mA}}
\]
(Valeur exacte avec le $-1$ : $3,08$~mA, différence négligeable).
\textbf{Analyse dimensionnelle :} $I_0$ (A) $\times$ $\exp(\text{sans dimension})$ = A. Correct.
\end{exoresolu}

\begin{aretenir}[Résistance dynamique de la diode]
En régime direct ($U > U_\gamma$), la résistance dynamique (ou différentielle) est :
\[
r_d = \frac{dU}{dI} = \frac{n V_T}{I} \approx \frac{n \times 26~\text{mV}}{I(\text{mA})}~\Omega
\]
Exemple : pour $I = \SI{1}{\milli\ampere}$, $n=1$, $r_d \approx \SI{26}{\ohm}$. Pour $I = \SI{10}{\milli\ampere}$, $r_d \approx \SI{2,6}{\ohm}$. \textbf{La diode se comporte comme une résistance variable commandée par le courant.}
\end{aretenir}

\courssub{Caractéristiques de sortie du transistor bipolaire (BJT)}

Le transistor bipolaire (NPN ou PNP) est un dipôle \textbf{à trois bornes} (Base B, Collecteur C, Émetteur E). Il faut deux familles de courbes pour le décrire complètement. Nous nous concentrons ici sur les \cle{courbes de sortie} (caractéristiques $I_C = f(V_{CE})$ pour $I_B$ imposé), essentielles pour le point de fonctionnement en régime commun-émetteur.

\begin{experience}[Famille de courbes $I_C(V_{CE})$ d'un transistor 2N2222 (NPN)]
\textbf{Montage :} Générateur $V_{BB}$ + $R_B$ pour imposer $I_B$ (échelonné : 0, 10, 20, 50, 100~$\mu$A). Circuit collecteur : $V_{CC}$ variable (0--15~V) + $R_C$ (protection) + Ampèremètre $I_C$ + Voltmètre $V_{CE}$.\par
\textbf{Protocole :} Pour chaque $I_B$ fixe, faire varier $V_{CC}$ et relever $(V_{CE}, I_C)$.\par
\textbf{Observations :} Trois zones distinctes apparaissent.
\end{experience}

\begin{popfigure}
\begin{tikzpicture}
\begin{axis}[
    width=\linewidth, height=8cm,
    axis lines=left,
    xlabel={$V_{CE}$ (V)}, ylabel={$I_C$ (mA)},
    xmin=0, xmax=12, ymin=0, ymax=25,
    xtick={0,0.2,2,4,6,8,10,12}, ytick={0,5,10,15,20,25},
    grid=major, grid style={popline},
    tick label style={font=\small},
    label style={font=\small},
    legend style={at={(0.98,0.98)}, anchor=north east, font=\small, fill=popcream},
]
% Ib = 0 (coupure)
\addplot[popmuted, thick, domain=0:12, samples=2] {0.02}; % ICEO ~ 0
\addlegendentry{$I_B = 0~\mu\text{A}$ (Coupure)}
% Ib = 20 uA
\addplot[popblue, very thick] coordinates {(0,0) (0.15, 2.5) (0.3, 4.8) (12, 5.0)};
\addlegendentry{$I_B = 20~\mu\text{A}$}
% Ib = 50 uA
\addplot[popgreen, very thick] coordinates {(0,0) (0.15, 6.5) (0.3, 12.2) (12, 12.5)};
\addlegendentry{$I_B = 50~\mu\text{A}$}
% Ib = 100 uA
\addplot[poporange, very thick] coordinates {(0,0) (0.15, 13) (0.3, 23) (12, 24)};
\addlegendentry{$I_B = 100~\mu\text{A}$}
% Zone labels
\node[popdark, rotate=90, font=\small] at (axis cs:0.1, 12) {Saturation};
\node[popdark, font=\small] at (axis cs:6, 22) {Zone active (linéaire)};
\node[popdark, font=\small] at (axis cs:6, 2) {Coupure};
% Droite de charge exemple (pointillé)
\addplot[poppurple, dashed, thick, domain=0:12] { (12-x)/470 * 1000 }; % Vcc=12V, Rc=470ohm -> Ic max = 25.5mA
\node[poppurple, anchor=south, font=\small] at (axis cs:10, 5) {Droite de charge $R_C=\SI{470}{\ohm}$};
\end{axis}
\end{tikzpicture}
\legende{Famille de courbes de sortie $I_C = f(V_{CE})$ pour un transistor NPN (type 2N2222). Paramètre : $I_B$. Trois zones : saturation ($V_{CE} < \SI{0,2}{\volt}$), active ($V_{CE} > \SI{0,5}{\volt}$, courbes plates), coupure ($I_B=0$). La droite de charge (pointillée) intersecte les courbes : les intersections sont les points de fonctionnement possibles.}
\end{popfigure}

\begin{definition}[Zones de fonctionnement du transistor BJT (NPN)]
\begin{itemize}
    \item \textbf{Zone de coupure :} $I_B \le 0$ (jonction BE bloquée). $I_C \approx I_{CEO}$ (très faible, $\mu$A). Le transistor est un \textbf{interrupteur ouvert} ($R_{CE} \to \infty$).
    \item \textbf{Zone active (ou linéaire) :} $I_B > 0$ et $V_{CE} > V_{CE(sat)} \approx \SI{0,2}{\volt}$ (jonction BE directe, BC bloquée). Les courbes sont quasi horizontales : $I_C \approx \beta \, I_B$ (indépendant de $V_{CE}$). Le transistor est un \textbf{générateur de courant commandé en courant} (amplificateur).
    \item \textbf{Zone de saturation :} $I_B > 0$ et $V_{CE} < V_{CE(sat)}$ (jonctions BE et BC directes). $I_C$ chute brutalement, limitée par le circuit externe ($V_{CC}, R_C$). $I_C < \beta I_B$. Le transistor est un \textbf{interrupteur fermé} ($V_{CE} \approx \SI{0,1}{\volt}$--\SI{0,2}{\volt}, $R_{CE} \approx 0$).
\end{itemize}
\end{definition}

\begin{methode}[Lecture du gain $\beta$ (ou $h_{FE}$) sur les courbes de sortie]
\begin{enumerate}
    \item Choisir un point de fonctionnement dans la \textbf{zone active} (courbe plate, ex. $V_{CE} = \SI{5}{\volt}$).
    \item Lire l'ordonnée $I_C$ (en mA) sur cette courbe.
    \item Identifier la valeur du paramètre $I_B$ (en $\mu$A) de cette courbe.
    \item Calculer $\beta = \dfrac{I_C(\text{mA}) \times 10^{-3}}{I_B(\mu\text{A}) \times 10^{-6}} = \dfrac{I_C(\text{mA})}{I_B(\mu\text{A})} \times 10^3$.
\end{enumerate}
\textbf{Exemple sur la figure :} Courbe $I_B = \SI{50}{\micro\ampere}$, à $V_{CE}=\SI{5}{\volt}$, $I_C \approx \SI{12,5}{\milli\ampere}$. $\beta = \frac{12,5}{50} \times 1000 = 250$.
\end{methode}

\begin{attention}[$\beta$ n'est pas une constante !]
$\beta = h_{FE}$ dépend de :
\begin{itemize}
    \item $I_C$ (pic à courant moyen, chute à bas et haut courant).
    \item $V_{CE}$ (effet Early : pente légère en zone active $\Rightarrow \beta$ augmente avec $V_{CE}$).
    \item Température $T$ (augmente avec $T$).
    \item Dispersion de fabrication (étalement 100--500 pour un même modèle).
\end{itemize}
Au Probatoire, si une valeur de $\beta$ n'est pas donnée, \textbf{utilise celle lue sur la courbe au point de fonctionnement considéré}, pas une valeur « typique » du datasheet.
\end{attention}

\courssub{Influence de la température (Point de programme Probatoire)}

La température $T$ modifie les caractéristiques de \textbf{tous} les semi-conducteurs. C'est un sujet classique au Probatoire (question de cours ou analyse de résultats).

\begin{propriete}[Effets de la température $T$]
\begin{center}
\begin{tabularx}{\linewidth}{|l|X|X|}\hline
\textbf{Composant} & \textbf{Effet sur la caractéristique} & \textbf{Conséquence pratique} \\ \hline
\textbf{Résistor (métal)} & $R$ augmente avec $T$ (coefficient $\alpha > 0$). $R(T) = R_0[1 + \alpha(T-T_0)]$. & Dérive du point de fonctionnement, auto-échauffement (effet Joule). \\ \hline
\textbf{Diode (Si)} & $I_0$ \textbf{augmente très fort} ($\propto T^3 e^{-E_g/kT}$). $V_T = kT/q$ augmente. \newline \textbf{À $I$ constant :} $U_\gamma$ \textbf{diminue} d'environ \textbf{-2 mV/°C}. \newline \textbf{À $U$ constant :} $I$ \textbf{augmente} exponentiellement. & Risque d'\textbf{emballement thermique} (thermal runaway) si polarisée en tension fixe. Nécessite une résistance de limitation en série. \\ \hline
\textbf{Transistor (BJT)} & $\beta$ \textbf{augmente} avec $T$. $V_{BE}$ (seuil) \textbf{diminue} ($\approx -2$~mV/°C). $I_{CBO}$ (courant de fuite collecteur-base) \textbf{double tous les 10°C}. & Dérive du point de repos ($I_C$ augmente si $V_{BE}$ fixe). Instabilité thermique en classe A. \\ \hline
\end{tabularx}
\end{center}
\end{propriete}

\begin{savaistu}[Le barrage de Lom Pangar et la température des composants]
Au Cameroun, les centrales hydroélectriques (Lom Pangar, Nachtigal, Edéa) alimentent le réseau via des postes de transformation où des thyristors et diodes de puissance (redresseurs HVDC, FACTS) commutent des centaines d'ampères. Leur caractéristique $I(U)$ dépend de $T$. Si le refroidissement (radiateurs, huile, air forcé) tombe en panne, $T$ monte, $U_\gamma$ baisse, le courant augmente $\rightarrow$ encore plus de chaleur $\rightarrow$ destruction. C'est pourquoi la \textbf{protection thermique} (sondes PT100, disjoncteurs thermiques) est critique sur les postes SONATREL/ENEO. Un ingénieur doit connaître la dérivée $\partial U/\partial T$ pour dimensionner les radiateurs !
\end{savaistu}

\begin{exoresolu}[Détermination graphique de $r_d$ et $\beta$]
\textbf{Énoncé :} On donne la caractéristique $I_C(V_{CE})$ pour $I_B = \SI{30}{\micro\ampere}$ (courbe bleue ci-dessous). En zone active ($V_{CE}=\SI{6}{\volt}$), $I_C = \SI{8,5}{\milli\ampere}$. La pente de la courbe en zone active est $\Delta I_C / \Delta V_{CE} = \SI{0,1}{\milli\ampere\per\volt}$.
\begin{enumerate}
    \item Calculer $\beta$ (gain de courant) en ce point.
    \item Calculer la résistance de sortie $r_{ce}$ (résistance dynamique collecteur-émetteur).
    \item Quelle serait la tension $V_{CE}$ de saturation si $I_B = \SI{30}{\micro\ampere}$ et $R_C = \SI{1}{\kilo\ohm}$, $V_{CC} = \SI{12}{\volt}$ ?
\end{enumerate}
\textbf{Solution :}
\begin{enumerate}
    \item $\beta = \dfrac{I_C}{I_B} = \dfrac{8,5 \times 10^{-3}}{30 \times 10^{-6}} = \boxed{283}$.
    \item $r_{ce} = \dfrac{dV_{CE}}{dI_C} = \dfrac{1}{0,1~\text{mA/V}} = \dfrac{1}{10^{-4}~\text{S}} = \boxed{10~\text{k}\Omega}$. (C'est la pente inverse de l'effet Early).
    \item En saturation, $I_C$ est limité par le circuit : $I_{C(sat)} \approx \dfrac{V_{CC}}{R_C} = \dfrac{12}{1000} = \SI{12}{\milli\ampere}$. Comme $I_{C(sat)} = \SI{12}{\milli\ampere} > \beta I_B = \SI{8,5}{\milli\ampere}$, le transistor \textbf{sature} (forcé par $R_C$). $V_{CE(sat)} \approx \SI{0,15}{\volt}$ (lu sur la courbe $I_B=\SI{30}{\micro\ampere}$ pour $I_C=\SI{12}{\milli\ampere}$, extrapolation vers $V_{CE}$ bas).
\end{enumerate}
\end{exoresolu}

\begin{exercice}[Caractéristiques composants — Entraînement Probatoire]
\begin{enumerate}
    \item Une diode 1N4007 ($I_0 = \SI{5}{\nano\ampere}$, $n=1,8$) est polarisée sous $U = \SI{0,65}{\volt}$ à $T=\SI{300}{\kelvin}$. Calculer $I$. Que vaut $r_d$ en ce point ?
    \item Sur les courbes de sortie d'un BC547 (NPN), la courbe $I_B = \SI{20}{\micro\ampere}$ passe par $I_C = \SI{4,0}{\milli\ampere}$ à $V_{CE}=\SI{5}{\volt}$. La pente en zone active est $\SI{0,05}{\milli\ampere\per\volt}$. Déterminer $\beta$ et $r_{ce}$.
    \item \textbf{QCM :} La température d'une diode Si polarisée en direct constant ($U = \SI{0,7}{\volt}$) augmente de $20^\circ$C. Le courant : (a) diminue (b) augmente (c) ne change pas. Justifier en une phrase.
    \item \textbf{Analyse :} Pourquoi une résistance de base $R_B$ est-elle indispensable pour piloter un transistor en commutation (saturation) depuis une sortie logique 3,3~V (ex. microcontrôleur STM32) ?
\end{enumerate}
\end{exercice}

\begin{corrige}[Exercice Caractéristiques]
\begin{enumerate}
    \item $V_T = \SI{0,02585}{\volt}$. $U/(nV_T) = 0,65/(1,8\times0,02585) = 13,96$. $I \approx 5\times10^{-9} \times e^{13,96} = 5\times10^{-9} \times 1,15\times10^6 = \boxed{5,75~\text{mA}}$. $r_d = nV_T/I = (1,8\times0,02585)/0,00575 = \boxed{8,1~\Omega}$.
    \item $\beta = 4,0~\text{mA} / 20~\mu\text{A} = 200$. $r_{ce} = 1 / (0,05~\text{mA/V}) = 1 / (5\times10^{-5}) = \boxed{20~\text{k}\Omega}$.
    \item \textbf{(b) augmente.} À $U$ constant, $I \propto e^{qU/nkT} \times T^3$ (via $I_0$). L'exponentielle domine : $I$ croit exponentiellement avec $T$. C'est l'emballement thermique.
    \item La sortie MCU fournit $U_{OH} \approx \SI{3,3}{\volt}$. La jonction BE du transistor a $U_\gamma \approx \SI{0,7}{\volt}$. Sans $R_B$, la tension \SI{3,3}{\volt} s'appliquerait quasi directement sur la jonction BE (diode), provoquant un courant destructeur $I_B \gg I_{B(max)}$. $R_B$ limite $I_B$ à $(3,3 - 0,7)/R_B$. Pour saturer un transistor pilotant $I_C = \SI{100}{\milli\ampere}$ avec $\beta_{min}=50$, il faut $I_B > \SI{2}{\milli\ampere}$. $R_B < (2,6~\text{V})/2~\text{mA} = \SI{1,3}{\kilo\ohm}$. On choisit $R_B = \SI{1}{\kilo\ohm}$.
\end{enumerate}
\end{corrige}

\courssub{Synthèse des caractéristiques $I(U)$}

\begin{center}
\begin{tabularx}{\linewidth}{|l|X|X|X|}\hline
\textbf{Dipôle} & \textbf{Équation / Forme} & \textbf{Paramètres clés} & \textbf{Résistance dynamique} \\ \hline
Résistor & $U = R I$ (Droite origine) & $R$ (cst) & $r_d = R$ (cst) \\ \hline
Diode (Si) & $I = I_0(e^{U/nV_T}-1)$ & $I_0, n, V_T$ & $r_d = \dfrac{nV_T}{I} \approx \dfrac{26n}{I(\text{mA})}~\Omega$ \\ \hline
Transistor (Zone active) & $I_C = \beta I_B$ (Horizontale) & $\beta(I_C, V_{CE}, T)$ & $r_{ce} = \dfrac{\Delta V_{CE}}{\Delta I_C}$ (Early, $\sim 10$--$100~\text{k}\Omega$) \\ \hline
Transistor (Saturation) & $V_{CE} \approx 0,15$~V & $V_{CE(sat)}, \beta_{forced} = I_C/I_B < \beta$ & $r_{ce(sat)} \approx 0,1$--$1~\Omega$ \\ \hline
\end{tabularx}
\end{center}

\begin{aretenir}[Check-list avant la section 3 (Point de fonctionnement)]
\begin{itemize}
    \item Savoir tracer/identifier les 3 caractéristiques de base.
    \item Savoir lire $\beta$ sur courbes de sortie (zone active seulement !).
    \item Connaître la formule $r_d = nV_T/I$ pour la diode.
    \item Retenir : $U_\gamma$ diode $\downarrow$ quand $T \uparrow$ (-2 mV/°C) ; $\beta$ transistor $\uparrow$ quand $T \uparrow$.
    \item \textbf{Toujours} dessiner les flèches $U, I$ (convention récepteur) sur le schéma du dipôle avant d'écrire une équation.
\end{itemize}
\end{aretenir}

La section 3 utilisera ces courbes pour résoudre graphiquement (et analytiquement) l'intersection entre la \textbf{droite de charge} (imposée par le circuit externe : source + résistances) et la \textbf{caractéristique du composant} (imposée par sa physique). C'est là que naît le \textbf{point de fonctionnement}.

\coursec{Point de fonctionnement : définition et méthode graphique}

\courssub{Transition depuis les caractéristiques I(U)}

Dans la section précédente, nous avons appris à \emph{caractériser} un dipôle par sa courbe $I = f(U)$, c'est-à-dire à décrire son comportement intrinsèque, indépendamment du circuit dans lequel il est inséré. Nous savons maintenant tracer la caractéristique d'une résistance (droite passant par l'origine), d'une diode (seuil vers $\SI{0,6}{V}$) ou d'un transistor (famille de courbes $I_C = f(U_{CE})$ pour $I_B$ fixé).

Mais un composant ne vit jamais seul : il est alimenté par une source (pile, générateur, alimentation secteur) et souvent limité par une résistance de protection ou de polarisation. La question centrale devient donc : \textbf{quel couple tension/courant $(U, I)$ le dipôle va-t-il \emph{réellement} atteindre dans ce circuit précis ?} C'est toute la problématique du \textbf{point de fonctionnement} (PF).

\courssub{Définition rigoureuse du point de fonctionnement}

Considérons le circuit série le plus simple : une source de tension idéale $E$ (force électromotrice), une résistance $R$ (résistance interne de la source ou résistance ajoutée en série) et un dipôle $\mathcal{D}$ (non linéaire en général) dont on connaît la caractéristique $I = f(U)$. Le circuit est parcouru par un courant $I$ et la tension aux bornes du dipôle est $U$ (convention récepteur).

\begin{center}
\begin{popfigure}
\begin{tikzpicture}[european, scale=1.1]
% Source
\draw (0,0) to[battery1, l={$E$}, invert] (0,2);
% Resistance
\draw (0,2) to[R=$R$, i={$i$}] (3,2);
% Dipole D
\draw (3,2) to[generic, l={$\mathcal{D}$}] (3,0);
% Ground/return
\draw (3,0) -- (0,0);
% Voltage arrow for E
\draw[<->, popblue, thick] (0.5,0.2) -- (0.5,1.8) node[midway, right] {$E$};
% Voltage arrow for U
\draw[<->, poporange, thick] (3.5,0.2) -- (3.5,1.8) node[midway, right] {$U$};
% Current arrow
\draw[->, popgreen, thick] (1.5, 2.3) -- (2.5, 2.3) node[midway, above] {$I$};
\end{tikzpicture}
\legende{Circuit série source-résistance-dipôle : convention récepteur pour $\mathcal{D}$.}
\end{popfigure}
\end{center}

Appliquons la \textbf{loi des mailles} (somme des tensions = 0) dans le sens du courant $I$ :
\[ E - R I - U = 0 \quad \Longrightarrow \quad U = E - R I. \]
Cette équation relie $U$ et $I$ \emph{imposée par le circuit d'alimentation}. Elle représente une droite dans le repère $(U, I)$ (abscisses $U$, ordonnées $I$ pour coller à la convention $I=f(U)$).

\begin{definition}[Point de fonctionnement (PF)]
Le \textbf{point de fonctionnement} d'un dipôle $\mathcal{D}$ dans un circuit donné est le couple $(U_0, I_0)$ qui satisfait \emph{simultanément} :
\begin{enumerate}
    \item La \textbf{caractéristique intrinsèque} du dipôle : $I_0 = f(U_0)$.
    \item La \textbf{loi du circuit} (appelée \textbf{courbe de charge}) : $U_0 = E - R I_0$ (ou $I_0 = \frac{E - U_0}{R}$).
\end{enumerate}
Graphiquement, c'est le point d'intersection de la courbe $I = f(U)$ et de la droite $I = \frac{E - U}{R}$.
\end{definition}

\courssub{Construction de la courbe de charge}

La courbe de charge (ou droite de charge) est l'ensemble des couples $(U, I)$ compatibles avec la source $E$ et la résistance $R$, \emph{quel que soit le dipôle $\mathcal{D}$}. C'est une droite affine décroissante.

Pour la tracer rapidement, on détermine ses deux intercepts :
\begin{itemize}
    \item \textbf{Point A (court-circuit du dipôle) :} $U = 0 \Rightarrow I_{cc} = \frac{E}{R}$. C'est le courant maximal que le circuit peut délivrer (court-circuit sur $\mathcal{D}$).
    \item \textbf{Point B (circuit ouvert) :} $I = 0 \Rightarrow U_{voc} = E$. C'est la tension maximale aux bornes de $\mathcal{D}$ (circuit ouvert).
\end{itemize}
La droite passe par $A(0, \frac{E}{R})$ et $B(E, 0)$. Son équation sous forme réduite (avec $U$ en abscisse) est :
\[ I = -\frac{1}{R} U + \frac{E}{R}. \]
La pente est $-\frac{1}{R}$ (négative) et l'ordonnée à l'origine est $\frac{E}{R}$.

\begin{attention}[Ne pas confondre les deux courbes !]
C'est l'erreur \textbf{n°1} au Probatoire.
\begin{itemize}
    \item La \textbf{caractéristique $I=f(U)$} : propriété \emph{propre} au dipôle (résistance, diode, transistor...). Elle ne dépend \emph{pas} de $E$ ni de $R$.
    \item La \textbf{courbe de charge} : propriété du \emph{circuit d'alimentation} ($E$ et $R$). Elle ne dépend \emph{pas} de la nature du dipôle $\mathcal{D}$.
\end{itemize}
Le PF est \emph{l'intersection} des deux. Si tu changes $R$, la courbe de charge pivote (pente change) ; si tu changes $E$, elle se traduit parallèlement. La caractéristique du dipôle, elle, ne bouge pas.
\end{attention}

\courssub{Méthode graphique : détermination pas à pas}

\begin{methode}[Détermination graphique du Point de Fonctionnement]
\textbf{Étape 1 : Tracer la caractéristique du dipôle $I = f(U)$.}
\begin{itemize}
    \item Repère orthonormé : $U$ (V) en abscisses, $I$ (A ou mA) en ordonnées.
    \item Tracer la courbe connue (droite pour $R$, courbe exponentielle pour diode, famille pour transistor).
\end{itemize}

\textbf{Étape 2 : Tracer la courbe de charge (droite).}
\begin{itemize}
    \item Calculer $I_{cc} = E/R$ (ordonnée à l'origine).
    \item Noter $U_{voc} = E$ (abscisse à l'origine).
    \item Tracer la droite passant par $(0, I_{cc})$ et $(E, 0)$.
\end{itemize}

\textbf{Étape 3 : Identifier l'intersection.}
\begin{itemize}
    \item Le point d'intersection $P(U_0, I_0)$ est le Point de Fonctionnement.
    \item Lire graphiquement $U_0$ (abscisse) et $I_0$ (ordonnée).
    \item \emph{Vérification :} $U_0 + R I_0 \approx E$ (loi des mailles).
\end{itemize}
\end{methode}

\courssub{Illustration : Diode en série avec une résistance}

Prenons l'exemple classique : une diode Silicon (seuil $U_{seuil} \approx \SI{0,6}{V}$) alimentée par $E = \SI{5,0}{V}$ à travers une résistance $R = \SI{1,0}{k\Omega}$.

\begin{exemplebox}[Point de fonctionnement d'une diode 1N4148]
\textbf{Données :} $E = \SI{5,0}{V}$, $R = \SI{1,0}{k\Omega}$. Diode Si : caractéristique modélisée par $I = 0$ pour $U < \SI{0,6}{V}$ et $I = I_S(e^{U/U_T}-1)$ pour $U \ge \SI{0,6}{V}$ (avec $I_S \approx \SI{e-12}{A}$, $U_T \approx \SI{0,026}{V}$ à $\SI{25}{\celsius}$).

\textbf{1. Courbe de charge :}
$I_{cc} = 5,0 / 1000 = \SI{5,0}{mA}$. $U_{voc} = \SI{5,0}{V}$.
Droite passant par $(0, \SI{5,0}{mA})$ et $(\SI{5,0}{V}, 0)$.

\textbf{2. Intersection graphique (voir figure ci-dessous) :}
La droite coupe la caractéristique de la diode vers $U_0 \approx \SI{0,72}{V}$.
Lecture de l'ordonnée : $I_0 \approx \SI{4,3}{mA}$.

\textbf{3. Vérification numérique (modèle de Shockley) :}
On résout $I = I_S(e^{U/U_T}-1)$ et $U = 5 - 1000 I$.
Par itération ou calculatrice : $U_0 \approx \SI{0,718}{V}$, $I_0 \approx \SI{4,28}{mA}$.
La lecture graphique est cohérente (précision $\pm \SI{0,02}{V}$, $\pm \SI{0,1}{mA}$).
\end{exemplebox}

\begin{center}
\begin{popfigure}
\begin{tikzpicture}
\begin{axis}[
    width=0.9\linewidth,
    height=8cm,
    axis lines=middle,
    xlabel={$U\ (\mathrm{V})$},
    ylabel={$I\ (\mathrm{mA})$},
    xmin=-0.2, xmax=5.5,
    ymin=-0.2, ymax=5.5,
    xtick={0,0.6,0.72,1,2,3,4,5},
    xticklabels={$0$,$0{,}6$,$U_0$,$1$,$2$,$3$,$4$,$E=5$},
    ytick={0,1,2,3,4,4.28,5},
    yticklabels={$0$,$1$,$2$,$3$,$4$,$I_0$,$I_{cc}=5$},
    grid=major,
    grid style={popline},
    tick label style={font=\small},
    label style={font=\small},
    clip=false,
    domain=0:5.5,
    samples=200
]
% Load Line
\addplot[popblue, thick, domain=0:5] {5 - x} node[pos=0.1, above left, font=\small, popblue] {$I = \frac{E-U}{R}$};
% Diode Characteristic (Shockley approx)
% Is = 1e-12 A = 1e-9 mA, Ut = 0.026 V
% I(mA) = 1e-9 * (exp(U/0.026) - 1)
\addplot[poporange, thick, domain=0:0.85, restrict y to domain=-0.5:6] ({x}, {1e-9*(exp(x/0.026)-1)}) node[pos=0.9, above right, font=\small, poporange] {$I = f(U)$ Diode};
% Intersection point
\addplot[only marks, mark=*, mark size=3, popdark] coordinates {(0.718, 4.282)};
\node[popdark, font=\small, anchor=south west] at (axis cs:0.718, 4.282) {$PF(U_0, I_0)$};
% Dashed lines to axes
\draw[dashed, popmuted] (axis cs:0.718, 0) -- (axis cs:0.718, 4.282);
\draw[dashed, popmuted] (axis cs:0, 4.282) -- (axis cs:0.718, 4.282);
\end{axis}
\end{tikzpicture}
\legende{Point de fonctionnement d'une diode Si ($E=\SI{5}{V}, R=\SI{1}{k\Omega}$). Intersection de la caractéristique (orange) et de la courbe de charge (bleue).}
\end{popfigure}
\end{center}

\courssub{Unicité du point de fonctionnement : théorème et démonstration guidée}

Pour un dipôle \emph{passif} (qui ne délivre pas d'énergie, $U \cdot I \ge 0$) dont la caractéristique est \emph{strictement croissante} (résistance ohmique, diode polarisée direct, transistor en zone active), le point de fonctionnement est \textbf{unique}.

\begin{propriete}[Unicité du PF pour caractéristique strictement croissante]
Soit un dipôle passif de caractéristique $I = f(U)$ dérivable et \textbf{strictement croissante} ($f'(U) > 0$ pour tout $U$). Soit une courbe de charge $I = g(U) = \frac{E}{R} - \frac{1}{R}U$ (strictement décroissante car $R > 0$). Alors le système
\[ \begin{cases} I = f(U) \\ I = g(U) \end{cases} \]
admet \textbf{une et une seule solution} $(U_0, I_0)$.
\end{propriete}

\begin{experience}[Démonstration guidée de l'unicité]
\textbf{Objectif :} Prouver qu'il existe un unique point d'intersection.

\textbf{1. Existence (Théorème des valeurs intermédiaires).}
On considère la fonction $h(U) = f(U) - g(U) = f(U) + \frac{1}{R}U - \frac{E}{R}$.
\begin{itemize}
    \item $h$ est continue (somme de fonctions continues).
    \item En $U=0$ : $h(0) = f(0) - \frac{E}{R}$. Pour un dipôle passif, $f(0) \ge 0$ (souvent $f(0)=0$). Donc $h(0) \le 0$ (souvent $<0$ si $E>0$).
    \item En $U=E$ : $h(E) = f(E) + \frac{E}{R} - \frac{E}{R} = f(E) \ge 0$.
\end{itemize}
Comme $h(0) \le 0$ et $h(E) \ge 0$, d'après le TVI, il existe au moins un $U_0 \in [0, E]$ tel que $h(U_0) = 0$, c'est-à-dire $f(U_0) = g(U_0)$. \textbf{Il existe au moins un PF.}

\textbf{2. Unicité (Monotonie stricte).}
La dérivée de $h$ est $h'(U) = f'(U) + \frac{1}{R}$.
\begin{itemize}
    \item Par hypothèse, $f'(U) > 0$ (caractéristique strictement croissante).
    \item $R > 0 \Rightarrow \frac{1}{R} > 0$.
\end{itemize}
Donc $h'(U) > 0$ pour tout $U$. La fonction $h$ est \textbf{strictement croissante} sur $\mathbb{R}^+$.
Une fonction strictement croissante ne peut s'annuler \textbf{qu'une seule fois}.
Donc l'équation $h(U)=0$ a une solution unique. \textbf{Le PF est unique.}
\end{experience}

\begin{savaistu}[Cas particuliers : caractéristiques non monotones]
Certains dipôles (diode tunnel, diode Gunn, certains transistors en régime particulier) présentent une \textbf{région à résistance différentielle négative} : la caractéristique $I=f(U)$ \emph{décroît} sur un intervalle ($f'(U) < 0$).
Dans ce cas, $h'(U) = f'(U) + 1/R$ peut changer de signe. La courbe de charge (droite décroissante) peut couper la caractéristique \textbf{en 3 points} (ex: diode tunnel).
\begin{itemize}
    \item 2 points sont \textbf{stables} (fonctionnement normal, amplificateur ou oscillateur).
    \item 1 point est \textbf{instable} (le circuit ne s'y stabilise pas spontanément).
\end{itemize}
Ce phénomène est exploité dans les \textbf{oscillateurs à diode tunnel} (fréquences micro-ondes, radars, télécoms). Au programme de Première, on retient : \emph{« Pour les dipôles usuels (R, D, T en régime normal), le PF est unique. »}
\end{savaistu}

\courssub{Exercice résolu : Choix de la résistance de polarisation}

\begin{exoresolu}[Polarisation d'une LED rouge]
\textbf{Énoncé :} On souhaite faire fonctionner une LED rouge (seuil $U_{seuil} \approx \SI{1,8}{V}$, courant nominal $I_{nom} = \SI{20}{mA}$) à partir d'une pile $E = \SI{9,0}{V}$ (neuve). Déterminer la valeur de la résistance $R$ à placer en série pour obtenir le point de fonctionnement nominal. La résistance normalisée la plus proche est-elle adaptée ?

\textbf{Solution détaillée :}

\textbf{1. Analyse du Point de Fonctionnement visé (PF cible).}
On veut $I_0 = I_{nom} = \SI{20}{mA} = \SI{0,020}{A}$.
La tension aux bornes de la LED à ce courant est $U_0 \approx U_{seuil} = \SI{1,8}{V}$ (caractéristique quasi verticale au-delà du seuil).

\textbf{2. Loi des mailles (Courbe de charge).}
$E = U_0 + R I_0 \quad \Rightarrow \quad R = \frac{E - U_0}{I_0}$.

\textbf{3. Calcul numérique.}
\[ R = \frac{9,0 - 1,8}{0,020} = \frac{7,2}{0,020} = \SI{360}{\Omega}. \]

\textbf{4. Choix normalisé (série E12 : 10, 12, 15, 18, 22, 27, 33, 39, 47, 56, 68, 82 $\times$ puissances de 10).}
Valeurs voisines : $\SI{330}{\Omega}$ et $\SI{390}{\Omega}$.

\textbf{5. Vérification du PF réel avec $R = \SI{330}{\Omega}$ (valeur inférieure $\rightarrow$ courant supérieur).}
$I_0 \approx \frac{9,0 - 1,8}{330} = \frac{7,2}{330} \approx \SI{21,8}{mA}$.
Dépassement de $9\%$ par rapport au nominal ($\SI{20}{mA}$). \textbf{Risque d'échauffement et durée de vie réduite.}

\textbf{6. Vérification avec $R = \SI{390}{\Omega}$ (valeur supérieure $\rightarrow$ courant inférieur).}
$I_0 \approx \frac{7,2}{390} \approx \SI{18,5}{mA}$.
Courant $7,5\%$ en dessous du nominal. Luminosité légèrement moindre, mais \textbf{sécurité assurée}.

\textbf{Conclusion :} On choisira \textbf{$R = \SI{390}{\Omega}$} (ou $\SI{470}{\Omega}$ pour plus de marge). La puissance dissipée dans $R$ : $P_R = R I_0^2 \approx 390 \times (0,0185)^2 \approx \SI{0,13}{W}$. Une résistance $\frac{1}{4}\ \mathrm{W}$ (standard) convient.
\tcblower
\textbf{Note pédagogique :} Toujours choisir la résistance normalisée \emph{supérieure} à la valeur calculée pour un composant sensible au courant (LED, diode laser). Pour une résistance de tirage (pull-up) sur une entrée logique, on viserait plutôt la valeur inférieure pour garantir le niveau haut.
\end{exoresolu}

\courssub{Synthèse de la section}

\begin{aretenir}[Essentiel à retenir sur le Point de Fonctionnement]
\begin{itemize}
    \item Le \textbf{Point de Fonctionnement (PF)} $(U_0, I_0)$ est l'unique couple tension/courant compatible à la fois avec le \textbf{dipôle} (caractéristique $I=f(U)$) et le \textbf{circuit} (courbe de charge $U = E - RI$).
    \item \textbf{Méthode graphique universelle :} Tracer la caractéristique $\rightarrow$ Tracer la droite de charge (intercepts $E$ et $E/R$) $\rightarrow$ Lire l'intersection.
    \item \textbf{Unicité garantie} pour tout dipôle passif à caractéristique strictement croissante ($R$, $D$ direct, $T$ zone active) car la courbe de charge est strictement décroissante ($R>0$).
    \item \textbf{Attention :} La courbe de charge dépend de $E$ et $R$ (le circuit) ; la caractéristique dépend du dipôle seul. Ne jamais les confondre !
    \item \textbf{Application concrète :} Choix de $R$ pour polariser une LED, un transistor (polarisation base/collecteur), dimensionnement d'un circuit d'alimentation.
\end{itemize}
\end{aretenir}

\begin{exercice}[Entraînement : PF d'une diode Zener]
On monte une diode Zener (tension de rupture $V_Z = \SI{5,1}{V}$) en série avec une résistance $R = \SI{220}{\Omega}$ et une source variable $E$.
\begin{enumerate}
    \item Tracer la caractéristique complète de la Zener (direct + inverse) et la courbe de charge pour $E = \SI{12}{V}$.
    \item Déterminer graphiquement le PF. La Zener est-elle en régime de régulation ?
    \item Quelle valeur minimale de $E$ pour que la Zener fonctionne en régulation (courant inverse $|I_Z| \ge \SI{1}{mA}$) ?
\end{enumerate}
\end{exercice}

\begin{corrige}[Exercice : PF d'une diode Zener]
\begin{enumerate}
    \item \textbf{Caractéristique :} Direct : comme diode Si ($U_{seuil}\approx \SI{0,7}{V}$). Inverse : quasi verticale à $U = -\SI{5,1}{V}$ pour $I < 0$ (convention récepteur).
    \textbf{Courbe de charge ($E=\SI{12}{V}, R=\SI{220}{\Omega}$) :} $I_{cc} = 12/220 \approx \SI{54,5}{mA}$. $U_{voc} = \SI{12}{V}$. Droite passant par $(0, \SI{54,5}{mA})$ et $(\SI{12}{V}, 0)$.
    \item \textbf{Intersection :} La droite coupe la branche inverse au voisinage de $U \approx -\SI{5,1}{V}$. Lecture : $I_0 \approx -\SI{31}{mA}$ (courant inverse). Oui, la Zener est en régulation ($|I| > \SI{1}{mA}$). PF : $(-\SI{5,1}{V}, -\SI{31}{mA})$.
    \item \textbf{Condition $|I_Z| \ge \SI{1}{mA}$ en régime inverse.} Convention récepteur : $U_Z = -\SI{5,1}{V}$, $I_Z = -\SI{1}{mA}$. Loi des mailles : $E = U_Z + R I_Z$.
    Pour polariser la Zener en inverse, la source doit être branchée « à l'envers » par rapport au schéma de la définition (anode de la Zener au + de la source). On écrit les grandeurs en valeur absolue : $E = V_Z + R |I_Z|$.
    $E_{min} = 5,1 + 220 \times 0,001 = 5,1 + 0,22 = \SI{5,32}{V}$.
    Donc $E_{min} \approx \SI{5,3}{V}$.
\end{enumerate}
\end{corrige}

\coursec{Loi de Pouillet : puissance et énergie dans un dipôle}

Dans la section précédente, nous avons appris à déterminer le \cle{point de fonctionnement} d'un circuit : ce couple de valeurs $(U_0, I_0)$ qui fixe l'état électrique du dipôle. Mais savoir \emph{où} fonctionne un composant ne suffit pas encore à l'ingénieur ou au technicien : il faut aussi savoir \emph{combien d'énergie} ce composant consomme ou dissipe chaque seconde. C'est cette question — cruciale pour dimensionner une installation, choisir un câble, éviter la surchauffe d'un transistor ou estimer la facture ENEO d'un foyer à Yaoundé — qui fait l'objet de la présente section : la \cle{loi de Pouillet} et le bilan énergétique d'un dipôle.

\courssub{Puissance électrique instantanée}

\textbf{Observation et intuition.}\par
Branche un fer à repasser et une lampe LED sur le réseau \SI{220}{\volt}. En quelques minutes, le fer devient brûlant alors que la lampe reste tiède. Pourtant, les deux appareils sont soumis à la même tension ! La différence vient de la \emph{vitesse} à laquelle chacun transforme l'énergie électrique : le fer convertit environ \SI{1000}{\joule} chaque seconde, la LED seulement \SI{9}{\joule} par seconde. Cette « vitesse de conversion d'énergie » est précisément ce que la physique appelle la \cle{puissance}.

\begin{definition}[Puissance électrique instantanée]
La \cle{puissance électrique instantanée} $p(t)$ reçue (ou absorbée) par un dipôle traversé par un courant $i(t)$ et soumis à une tension $u(t)$ est :
\[ p(t) = u(t)\cdot i(t) \]
avec $u$ en volts ($\mathrm{V}$), $i$ en ampères ($\mathrm{A}$) et $p$ en watts ($\mathrm{W}$). En convention récepteur (flèches de $u$ et $i$ en sens opposés), $p>0$ signifie que le dipôle \emph{absorbe} de l'énergie ; $p<0$ signifie qu'il en \emph{fournit} (cas d'un générateur).
\end{definition}

\begin{experience}[Démonstration guidée : de la définition du travail électrique à $p = u\,i$]
Cette démonstration est formatrice et sa structure est exigible au Probatoire. Suivons-la pas à pas.
\begin{enumerate}
\item \textbf{Travail de la force électrique.} Une charge $q$ qui se déplace d'un point $A$ vers un point $B$ d'un circuit, entre lesquels règne la tension $u_{AB}$, reçoit (ou cède) le travail électrique :
\[ W_{AB} = q\cdot u_{AB}. \]
\item \textbf{Charge transitant pendant $\Delta t$.} Si le dipôle est traversé par un courant d'intensité $i$, la quantité de charge qui le traverse pendant la durée $\Delta t$ est :
\[ q = i\cdot \Delta t \qquad (\text{définition de l'intensité}). \]
\item \textbf{Énergie échangée.} En reportant : $W = i\,\Delta t \cdot u = u\,i\,\Delta t$.
\item \textbf{Passage à la puissance.} La puissance étant l'énergie échangée par unité de temps :
\[ p = \frac{W}{\Delta t} = u\cdot i. \]
\end{enumerate}
\textbf{Vérification dimensionnelle :} $[u\,i] = \mathrm{V}\cdot\mathrm{A} = \dfrac{\mathrm{J}}{\mathrm{C}}\cdot\dfrac{\mathrm{C}}{\mathrm{s}} = \dfrac{\mathrm{J}}{\mathrm{s}} = \mathrm{W}$. La relation est homogène.
\end{experience}

\courssub{Régime continu : la loi de Pouillet}

En régime continu (courant et tension constants dans le temps), la puissance instantanée ne dépend plus du temps : on note simplement $P = U\cdot I$.

\begin{definition}[Loi de Pouillet]
En régime continu, la \cle{puissance absorbée par un dipôle} soumis à la tension $U$ et traversé par le courant $I$ est :
\[ \boxed{P = U\cdot I} \]
$U$ en volts ($\mathrm{V}$), $I$ en ampères ($\mathrm{A}$), $P$ en watts ($\mathrm{W}$). Cette relation, dite \cle{loi de Pouillet}, est valable pour \emph{tout} dipôle (résistor, diode, transistor, moteur, électrolyseur...), pas seulement pour les résistors.
\end{definition}

\begin{propriete}[Formes équivalentes pour un résistor ohmique]
Pour un \cle{résistor} de résistance $R$ vérifiant la loi d'Ohm $U = R\,I$, la loi de Pouillet s'écrit sous trois formes équivalentes :
\[ P = U\cdot I = R\cdot I^2 = \frac{U^2}{R} \]
\textbf{Démonstration (exigible) :} en injectant $U = RI$ dans $P = UI$, on obtient $P = (RI)I = RI^2$ ; en injectant $I = U/R$, on obtient $P = U\cdot\dfrac{U}{R} = \dfrac{U^2}{R}$.
\textbf{Vérification dimensionnelle :} $[RI^2] = \Omega\cdot\mathrm{A}^2 = \mathrm{V}\cdot\mathrm{A} = \mathrm{W}$. Homogène.
\end{propriete}

\begin{attention}[Pièges fréquents au Probatoire]
\begin{itemize}
\item Les formes $P = RI^2$ et $P = U^2/R$ ne sont valables \textbf{que pour un résistor ohmique}. Pour une diode ou un transistor, seule $P = U\,I$ est utilisable, avec les valeurs $U_0$ et $I_0$ du point de fonctionnement.
\item Dans $P = U^2/R$, $U$ est la tension \emph{aux bornes du résistor}, pas nécessairement la tension du générateur.
\item Ne pas confondre puissance ($\mathrm{W}$) et énergie ($\mathrm{J}$ ou $\mathrm{kWh}$) : une puissance est un débit d'énergie, une énergie est une quantité accumulée.
\item Attention aux unités : convertir systématiquement les $\mathrm{mA}$ en $\mathrm{A}$ avant tout calcul ($10\ \mathrm{mA} = 0{,}010\ \mathrm{A}$).
\end{itemize}
\end{attention}

\begin{exemplebox}[Exemple chiffré : bouilloire électrique]
Une bouilloire utilisée dans un foyer de Douala est modélisée par un résistor de résistance $R = \SI{48,4}{\ohm}$ branché sur le réseau $U = \SI{220}{\volt}$.
\begin{align*}
I &= \frac{U}{R} = \frac{220}{48{,}4} \approx \SI{4,55}{\ampere}\\
P &= \frac{U^2}{R} = \frac{220^2}{48{,}4} = \frac{48\,400}{48{,}4} = \SI{1000}{\watt}
\end{align*}
Vérification croisée : $P = RI^2 = 48{,}4\times(4{,}55)^2 \approx \SI{1000}{\watt}$. Les trois formes donnent bien le même résultat : la bouilloire convertit \SI{1000}{\joule} d'énergie électrique en chaleur chaque seconde.
\end{exemplebox}

\begin{popfigure}
\begin{center}
\begin{tikzpicture}
\begin{axis}[
width=0.85\linewidth, height=6.5cm,
xlabel={$U$ (V)}, ylabel={$P$ (W)},
xmin=0, xmax=12, ymin=0, ymax=6,
grid=both, grid style={popline!40},
legend style={at={(0.03,0.97)}, anchor=north west, font=\small},
]
\addplot[popcurve, domain=0:11, samples=200, popblue, very thick] {x^2/24};
\addlegendentry{Résistor $R=\SI{24}{\ohm}$ : $P=\dfrac{U^2}{R}$ (parabole)}
\addplot[popcurve, domain=0:2.5, samples=200, poporange, very thick] {x*0.02*(exp(3.5*(x-0.6)))};
\addlegendentry{Diode : $P = U\,I(U)$ (croissance quasi exponentielle)}
\draw[dashed, popmuted] (axis cs:0.7,0) -- (axis cs:0.7,0.02);
\end{axis}
\end{tikzpicture}
\end{center}
\legende{Puissance dissipée en fonction de la tension. Pour le résistor, $P=U^2/R$ est une parabole ; pour la diode, la puissance reste quasi nulle sous la tension de seuil ($U_s \approx \SI{0,7}{\volt}$ pour le silicium) puis croît très rapidement.}
\end{popfigure}

\courssub{Énergie dissipée pendant une durée $\Delta t$}

\begin{definition}[Énergie électrique]
L'\cle{énergie} électrique reçue ou dissipée par un dipôle de puissance $P$ constante pendant la durée $\Delta t$ est :
\[ \boxed{W = P\cdot \Delta t = U\cdot I\cdot \Delta t} \]
avec $P$ en watts ($\mathrm{W}$), $\Delta t$ en secondes ($\mathrm{s}$) et $W$ en joules ($\mathrm{J}$). Pour un résistor : $W = R\,I^2\,\Delta t$ (\cle{effet Joule}).
\end{definition}

Dans la vie courante et sur les factures d'ENEO, on utilise le \cle{kilowattheure} ($\mathrm{kWh}$) : c'est l'énergie consommée par un appareil de puissance \SI{1}{\kilo\watt} fonctionnant pendant \SI{1}{\hour} :
\[ 1\ \mathrm{kWh} = 1000\ \mathrm{W}\times 3600\ \mathrm{s} = 3{,}6\times 10^{6}\ \mathrm{J} = \SI{3,6}{\mega\joule}. \]

\begin{exemplebox}[Coût d'utilisation de la bouilloire]
La bouilloire de \SI{1000}{\watt} fonctionne \SI{10}{\minute} par jour. Énergie quotidienne :
\[ W = P\cdot\Delta t = 1000\times 600 = \SI{6,0e5}{\joule} = \SI{0,167}{\kilo\watt\hour}. \]
Sur un mois ($30$ jours) : $W_{\text{mois}} \approx \SI{5,0}{\kilo\watt\hour}$. À environ \SI{100}{\mathrm{FCFA}} le $\mathrm{kWh}$ (tarif domestique), cela représente environ \SI{500}{\mathrm{FCFA}} par mois.
\end{exemplebox}

\courssub{Application aux composants connus : résistor, diode, transistor}

\textbf{Le résistor : effet Joule pur.}\par
Toute l'énergie électrique absorbée par un résistor est convertie en \emph{chaleur} : c'est l'\cle{effet Joule}. Il est utile (chauffage, cuisson, lampes à incandescence) ou parasité (échauffement des câbles, des composants). La puissance dissipée vaut $P = RI^2$ : elle croît avec le \emph{carré} du courant, ce qui explique pourquoi SONATREL transporte l'électricité des barrages de Lom Pangar ou d'Edéa sous très haute tension : à puissance transportée fixée, augmenter $U$ diminue $I$, donc diminue fortement les pertes $RI^2$ dans les lignes.

\textbf{La diode : puissance de seuil.}\par
Une diode au silicium en conduction présente une tension quasi constante $U_s \approx \SI{0,7}{\volt}$. La puissance qu'elle dissipe au point de fonctionnement $(U_s, I_0)$ est :
\[ P_D = U_s\cdot I_0. \]
Par exemple, pour $I_0 = \SI{20}{\milli\ampere}$ : $P_D = 0{,}7\times 0{,}020 = \SI{0,014}{\watt} = \SI{14}{\milli\watt}$. En dessous du seuil, le courant est quasi nul et la puissance dissipée est négligeable ; au-dessus, elle croît très vite (voir la courbe ci-dessus), d'où la nécessité de la résistance de protection en série.

\textbf{Le transistor : pertes collecteur-émetteur.}\par
En régime de commutation ou d'amplification, un transistor bipolaire est traversé par le courant collecteur $I_C$ et présente une tension $U_{CE}$ entre collecteur et émetteur. La puissance dissipée (pertes qui échauffent le composant) vaut :
\[ P_T = U_{CE}\cdot I_C. \]
Exemple : un transistor de point de fonctionnement $U_{CE} = \SI{6}{\volt}$, $I_C = \SI{50}{\milli\ampere}$ dissipe $P_T = 6\times 0{,}050 = \SI{0,30}{\watt}$. Si le boîtier ne peut évacuer que \SI{0,25}{\watt}, il faut ajouter un radiateur ou choisir un autre transistor.

\begin{methode}[Calculer la puissance dissipée par chaque élément d'un circuit série]
Une fois le point de fonctionnement déterminé (section précédente) :
\begin{enumerate}
\item Relever le courant $I_0$ du point de fonctionnement (le même pour tous les dipôles en série).
\item Relever la tension aux bornes de \emph{chaque} dipôle : $U_R$, $U_D$, $U_{CE}$... (la somme doit redonner la f.é.m. du générateur, loi des mailles).
\item Calculer pour chaque dipôle $P_k = U_k\cdot I_0$. Pour un résistor, vérifier avec $P_R = R\,I_0^2$.
\item Vérifier le bilan : la somme des puissances absorbées par les récepteurs doit égaler la puissance fournie par le générateur $P_g = E\cdot I_0$.
\item Comparer chaque puissance dissipée à la puissance maximale admissible du composant (donnée du constructeur) pour conclure sur la sécurité du montage.
\end{enumerate}
\end{methode}

\begin{exoresolu}[Circuit série : générateur, résistor et diode]
\textbf{Énoncé.} Un générateur de f.é.m. $E = \SI{9,0}{\volt}$ alimente en série un résistor $R = \SI{415}{\ohm}$ et une diode au silicium de tension de seuil $U_s = \SI{0,7}{\volt}$ (supposée parfaitement conductrice au-delà). Le point de fonctionnement donne $I_0 = \SI{20}{\milli\ampere}$.
\begin{enumerate}
\item Vérifier la tension aux bornes du résistor.
\item Calculer la puissance dissipée par le résistor, par la diode et la puissance fournie par le générateur.
\item Conclure sur le bilan énergétique.
\end{enumerate}
\tcblower
\textbf{Solution détaillée.}
\begin{enumerate}
\item Loi d'Ohm : $U_R = R\,I_0 = 415\times 0{,}020 = \SI{8,3}{\volt}$. Vérification par la loi des mailles : $U_R + U_s = 8{,}3 + 0{,}7 = \SI{9,0}{\volt} = E$. Cohérent.
\item Puissances :
\begin{align*}
P_R &= R\,I_0^2 = 415\times (0{,}020)^2 = \SI{0,166}{\watt}\\
P_D &= U_s\cdot I_0 = 0{,}7\times 0{,}020 = \SI{0,014}{\watt}\\
P_g &= E\cdot I_0 = 9{,}0\times 0{,}020 = \SI{0,180}{\watt}
\end{align*}
\item Bilan : $P_R + P_D = 0{,}166 + 0{,}014 = \SI{0,180}{\watt} = P_g$. La puissance fournie par le générateur est intégralement dissipée par les récepteurs : l'énergie se conserve. Le résistor absorbe environ $92\ \%$ de la puissance totale.
\end{enumerate}
\end{exoresolu}

\courssub{Lien avec le rendement d'un circuit d'alimentation}

Dans un circuit d'alimentation (chargeur de téléphone, alimentation d'un poste radio dans un village électrifié par un mini-barrage), seule une partie de la puissance fournie par la source est réellement \emph{utile} ; le reste est perdu par effet Joule (résistances internes, transistors, diodes).

\begin{definition}[Rendement d'un circuit d'alimentation]
Le \cle{rendement} $\eta$ d'un circuit d'alimentation est le rapport de la puissance utile $P_u$ (reçue par la charge) à la puissance totale $P_t$ fournie par la source :
\[ \eta = \frac{P_u}{P_t} = \frac{P_u}{P_u + P_{\text{pertes}}} \qquad (0 \leq \eta \leq 1) \]
On l'exprime souvent en pourcentage : $\eta(\%) = 100\times\dfrac{P_u}{P_t}$.
\end{definition}

\begin{exemplebox}[Rendement d'une alimentation simple]
Un générateur fournit $P_t = \SI{0,180}{\watt}$ au circuit précédent ; la diode est le composant utile (par exemple une LED témoin) et reçoit $P_u = \SI{0,014}{\watt}$. Le rendement vaut :
\[ \eta = \frac{0{,}014}{0{,}180} \approx 0{,}078 \approx \SI{7,8}{\percent}. \]
Rendement très faible : la majeure partie de l'énergie est perdue dans le résistor. C'est pourquoi les alimentations modernes (chargeurs de téléphones MTN/Orange) utilisent des circuits à découpage dont le rendement dépasse \SI{80}{\percent}.
\end{exemplebox}

\begin{attention}[Conditions d'application]
\begin{itemize}
\item Le rendement se calcule toujours avec des \emph{puissances} (ou des énergies sur une même durée), jamais avec des tensions ou des courants seuls.
\item $\eta$ est un nombre sans unité ; ne jamais écrire $\eta = 78\ \mathrm{W}$.
\item Un rendement supérieur à $1$ (ou $100\ \%$) est le signe certain d'une erreur de calcul.
\end{itemize}
\end{attention}

\begin{savaistu}[La loi de Pouillet et les radiateurs de puissance]
La loi de Pouillet est à la base du dimensionnement des \cle{radiateurs} en électronique embarquée. Un transistor de puissance (dans un amplificateur audio, un onduleur solaire ou le variateur d'un ventilateur) qui dissipe $P_T = U_{CE}\,I_C$ transforme cette puissance en chaleur : sa température de jonction s'élève et peut détruire le composant au-delà de \SI{150}{\celsius} environ. Les ingénieurs calculent donc $P_T$, puis choisissent un radiateur en aluminium (plaque ailetée) capable d'évacuer cette chaleur vers l'air. C'est exactement le même raisonnement qui préside au refroidissement des gros transformateurs SONATREL que l'on voit, baignant dans l'huile, dans les postes de distribution : l'énergie perdue par effet Joule, $P = RI^2$, doit impérativement être évacuée. À l'échelle d'un pays, réduire ces pertes, c'est économiser l'équivalent de la production de plusieurs turbines du barrage de Nachtigal !
\end{savaistu}

\begin{aretenir}[L'essentiel de la section]
\begin{itemize}
\item Puissance instantanée reçue par un dipôle : $p(t) = u(t)\,i(t)$ ; en continu : $P = U\,I$ (loi de Pouillet, valable pour tout dipôle).
\item Pour un résistor ohmique : $P = RI^2 = \dfrac{U^2}{R}$ (effet Joule).
\item Énergie sur une durée $\Delta t$ : $W = P\,\Delta t$ ; $1\ \mathrm{kWh} = \SI{3,6}{\mega\joule}$.
\item Diode : $P_D = U_s\,I_0$ ; transistor : $P_T = U_{CE}\,I_C$.
\item Bilan énergétique : la somme des puissances absorbées égale la puissance fournie ; rendement $\eta = P_u/P_t \leq 1$.
\end{itemize}
\end{aretenir}

\coursec{Détermination du point de fonctionnement dans des circuits simples}

La loi de Pouillet nous a appris à calculer la puissance mise en jeu dans un dipôle : $P = U \times I$. Mais pour appliquer cette loi, encore faut-il connaître le couple $(U, I)$ auquel le dipôle fonctionne réellement une fois branché dans le circuit. Ce couple, c'est précisément le \cle{point de fonctionnement}. Dans cette section, nous passons à la pratique : tu vas apprendre à déterminer ce point dans les deux circuits les plus classiques du programme — le circuit source–résistance–diode et le circuit à transistor en commutation — puis à vérifier que les puissances dissipées restent compatibles avec les limites des composants. C'est exactement le travail de l'ingénieur qui choisit les composants d'un circuit, par exemple pour l'éclairage à LED d'une maison rurale alimentée par un petit panneau solaire.

\courssub{Circuit série source--résistance--diode}

\textbf{Situation et mise en équation.}\par
Considérons le circuit le plus simple comportant un composant non linéaire : un générateur de f.é.m. $E$ (supposé idéal), un résistor de résistance $R$ et une diode $D$, tous en série.

\begin{popfigure}
\begin{center}
\begin{circuitikz}[european, scale=1.1, transform shape]
\draw (0,0) to[battery1, l={$E$}] (0,3)
      to[R=$R$, v={$U_R$}] (3,3)
      to[D=$D$, v={$U_D$}] (3,0)
      -- (0,0);
\draw[-{Stealth}, popblue, thick] (1.5,3.6) -- (2.2,3.6) node[midway, above, popblue]{$I$};
\fill[popink] (1.5,3) circle (1.5pt) node[below left=2pt]{A};
\fill[popink] (3,1.5) circle (1.5pt) node[right=2pt]{B};
\end{circuitikz}
\end{center}
\legende{Circuit série source--résistance--diode. Le point de fonctionnement est le couple $(U_D, I)$ réellement établi.}
\end{popfigure}

La loi des mailles donne immédiatement :
\[ E = U_R + U_D = R\,I + U_D. \]
D'autre part, la diode obéit à sa caractéristique $I = f(U_D)$, dont le modèle de Shockley (hors programme, mais utile pour comprendre) s'écrit $I = I_S\left(e^{U_D/U_T}-1\right)$ avec $U_T \approx \SI{25}{mV}$ à température ambiante. Le point de fonctionnement est donc la solution du système :
\[
\begin{cases}
I = \dfrac{E - U_D}{R} & \text{(droite de charge, imposée par le circuit)}\\[6pt]
I = f(U_D) & \text{(caractéristique de la diode)}
\end{cases}
\]

\begin{attention}[Une équation qu'on ne sait pas résoudre à la main]
En injectant le modèle de Shockley dans la loi des mailles, on obtient $E = R\,I_S\left(e^{U_D/U_T}-1\right) + U_D$ : c'est une \cle{équation transcendante}, qui mélange $U_D$ et son exponentielle. Elle ne possède \textbf{pas} de solution littérale exprimable avec les fonctions usuelles. Deux stratégies s'offrent alors à toi : la \textbf{méthode graphique} (intersection des courbes) et le \textbf{modèle simplifié} de la diode (tension de seuil constante $V_f$), qui donne une solution analytique immédiate.
\end{attention}

\begin{methode}[Détermination du point de fonctionnement d'un circuit R--D]
\begin{enumerate}
\item \textbf{Écrire la droite de charge} : à partir de la loi des mailles, exprimer $I = \dfrac{E - U_D}{R}$.
\item \textbf{Tracer cette droite} sur le graphe de la caractéristique $I(U_D)$ de la diode. Deux points suffisent : $(U_D = 0\,;\, I = E/R)$ et $(U_D = E\,;\, I = 0)$.
\item \textbf{Lire l'intersection} avec la caractéristique de la diode : c'est le point de fonctionnement $Q(U_{D0}, I_0)$.
\item \textbf{Affiner si besoin par itération numérique} : partir d'une valeur estimée de $U_D$, calculer $I = (E-U_D)/R$, lire sur la courbe la tension correspondant à ce courant, et recommencer jusqu'à stabilisation.
\item \textbf{Vérifier les puissances} : calculer $P_D = U_{D0} I_0$ et $P_R = R I_0^2$, puis comparer aux valeurs maximales admissibles du fabricant (\emph{\cle{datasheet}}).
\end{enumerate}
\end{methode}

\begin{exoresolu}[Circuit complet : source \SI{9}{V}, résistor \SI{470}{\ohm}, LED rouge]
Une LED rouge de \SI{5}{mm} (tension directe $V_f = \SI{2,0}{V}$ dans le modèle à seuil) est branchée en série avec un résistor $R = \SI{470}{\ohm}$ sur une pile de $E = \SI{9}{V}$. Déterminer le point de fonctionnement, puis les puissances dissipées par chaque dipôle.
\tcblower
\textbf{Solution.} Dans le modèle à seuil, la diode passante se comporte comme une source de tension $V_f$ : la maille s'écrit
\[ E = R\,I + V_f \quad\Longrightarrow\quad I = \frac{E - V_f}{R} = \frac{9 - 2{,}0}{470} = \SI{1,49e-2}{A} \approx \SI{15}{mA}. \]
Le point de fonctionnement est donc $Q\,(U_D = \SI{2,0}{V}\,;\, I = \SI{15}{mA})$. Puissances mises en jeu :
\begin{align*}
P_{\text{LED}} &= U_D \times I = 2{,}0 \times \num{1,49e-2} \approx \SI{30}{mW},\\
P_R &= R I^2 = 470 \times (\num{1,49e-2})^2 \approx \SI{104}{mW},\\
P_{\text{source}} &= E \times I = 9 \times \num{1,49e-2} \approx \SI{134}{mW}.
\end{align*}
On vérifie le bilan énergétique : $P_{\text{source}} = P_R + P_{\text{LED}}$ ($134 \approx 104 + 30$ mW). Un résistor standard $\SI{0,25}{W}$ convient largement, et la LED supporte couramment \SI{20}{mA} : le montage est sain.
\end{exoresolu}

\begin{exemplebox}[Choisir la résistance de protection d'une LED]
Tu veux alimenter une LED blanche de \SI{5}{mm} ($V_f = \SI{2,1}{V}$, courant nominal $I_f = \SI{20}{mA}$) sous $E = \SI{9}{V}$. La résistance de limitation se calcule en \emph{imposant} le point de fonctionnement souhaité :
\[ R = \frac{E - V_f}{I_f} = \frac{9 - 2{,}1}{\num{20e-3}} = \SI{345}{\ohm}. \]
Cette valeur n'existe pas dans la série normalisée E12 : on choisit la valeur standard la plus proche \textbf{disponible}, $R = \SI{330}{\ohm}$, ce qui donne $I = \dfrac{6{,}9}{330} \approx \SI{21}{mA}$, légèrement au-dessus du nominal mais généralement toléré. Choisir \SI{390}{\ohm} donnerait \SI{18}{mA} : la LED éclairerait un peu moins mais durerait plus longtemps. C'est un vrai choix d'ingénieur !
\end{exemplebox}

\courssub{Transistor NPN en commutation et droite de charge}

\textbf{Situation.}\par
Le transistor bipolaire NPN est l'interrupteur commandé de l'électronique : un faible courant de base $I_B$ contrôle un courant de collecteur $I_C$ beaucoup plus grand. En \cle{commutation}, on l'utilise en tout-ou-rien : soit \emph{bloqué} ($I_C = 0$, interrupteur ouvert), soit \emph{saturé} ($V_{CE} \approx 0$, interrupteur fermé). C'est le principe des commandes de relais, de moteurs ou de lampes dans les cartes électroniques.

\begin{popfigure}
\begin{center}
\begin{circuitikz}[european, scale=1.05, transform shape]
\draw (0,0) to[battery1, l={$E$}] (0,4) to[R=$R_C$, v={$U_{R_C}$}] (3.5,4) -- (3.5,3.2);
\draw (3.5,1.6) node[npn](Q){};
\draw (3.5,3.2) -- (Q.C);
\draw (Q.E) -- (3.5,0) -- (0,0);
\draw (Q.B) -- (1.8,1.15) to[R=$R_B$] (-1.2,1.15) to[battery1, l_={$E_B$}] (-1.2,0) -- (0,0);
\draw[-{Stealth}, popblue, thick] (4.4,2.9) -- (4.4,2.1) node[midway, right, popblue]{$I_C$};
\draw[-{Stealth}, poporange, thick] (0.6,1.6) -- (1.2,1.3) node[midway, above left, poporange]{$I_B$};
\node[popdark] at (4.6,0.9) {$V_{CE}$};
\draw[-{Stealth}, popdark] (4.3,1.7) -- (4.3,0.15);
\end{circuitikz}
\end{center}
\legende{Transistor NPN en commutation : la base est pilotée via $R_B$, la charge $R_C$ est alimentée sous $E$.}
\end{popfigure}

\textbf{Droite de charge du transistor.}\par
Dans la maille collecteur--émetteur, la loi des mailles donne :
\[ E = R_C I_C + V_{CE} \quad\Longrightarrow\quad \boxed{\,I_C = \frac{E - V_{CE}}{R_C}\,} \]
C'est l'équation de la \cle{droite de charge} dans le plan $(V_{CE}, I_C)$. Elle passe par les points remarquables $(V_{CE} = 0\,;\, I_C = E/R_C)$ — le \emph{courant de saturation maximal} — et $(V_{CE} = E\,;\, I_C = 0)$ — le \emph{blocage}. Le point de fonctionnement est l'intersection de cette droite avec la courbe $I_C(V_{CE})$ correspondant au courant de base $I_B$ imposé.

\begin{popfigure}
\begin{center}
\begin{tikzpicture}
\begin{axis}[
  width=0.82\linewidth, height=6.2cm,
  xlabel={$V_{CE}$ (V)}, ylabel={$I_C$ (mA)},
  xmin=0, xmax=13, ymin=0, ymax=14,
  grid=both, grid style={popline!40},
  axis line style={popdark},
  legend style={font=\small, at={(0.97,0.97)}, anchor=north east}
]
\addplot[popcurve, domain=0:12, samples=100, popblue] {12 - x};
\addlegendentry{Droite de charge $I_C=(E-V_{CE})/R_C$}
\addplot[popcurve, domain=0:12, samples=120, poporange] {11*(1-exp(-x/0.35))};
\addlegendentry{$I_C(V_{CE})$ pour $I_B$ saturant}
\addplot[popcurve, domain=0:12, samples=120, popgreen] {5*(1-exp(-x/0.5))};
\addlegendentry{$I_C(V_{CE})$ pour $I_B$ faible (régime linéaire)}
\addplot[only marks, mark=*, mark size=2pt, popink] coordinates {(0.2,11.8)};
\node[popink, anchor=west] at (axis cs:0.5,12.6) {$Q_{\text{sat}}$};
\addplot[only marks, mark=*, mark size=2pt, popink] coordinates {(7.4,4.6)};
\node[popink, anchor=south west] at (axis cs:7.4,4.9) {$Q_{\text{lin}}$};
\end{axis}
\end{tikzpicture}
\end{center}
\legende{Droite de charge et réseau de caractéristiques du transistor. Selon $I_B$, le point de fonctionnement est en régime linéaire ($Q_{\text{lin}}$) ou en saturation ($Q_{\text{sat}}$).}
\end{popfigure}

\textbf{Recherche du point de saturation.}\par
En régime linéaire, $I_C = \beta\, I_B$. Mais cette relation ne peut pas dépasser la limite imposée par la droite de charge : $I_C \leq E/R_C$. Le transistor \emph{sature} lorsque le courant « demandé » $\beta I_B$ excède ce que le circuit peut fournir :
\[ \beta I_B \geq \frac{E - V_{CE,\text{sat}}}{R_C} \quad\Longrightarrow\quad \text{saturation, avec } V_{CE} = V_{CE,\text{sat}} \approx \SI{0,2}{V}. \]
Le courant de saturation vaut alors $I_{C,\text{sat}} = \dfrac{E - V_{CE,\text{sat}}}{R_C}$.

\begin{attention}[Le piège classique : oublier $V_{BE}$ dans le calcul de $I_B$]
Quand la base est pilotée à travers $R_B$ par une source $E_B$, le courant de base n'est \textbf{pas} $E_B/R_B$ ! La jonction base--émetteur se comporte comme une diode passante : il faut retrancher sa tension de seuil $V_{BE} \approx \SI{0,7}{V}$ :
\[ I_B = \frac{E_B - V_{BE}}{R_B}. \]
Oublier $V_{BE}$ surestime $I_B$, parfois du simple au double si $E_B$ est faible (par exemple $E_B = \SI{5}{V}$ : l'erreur relative vaut $0{,}7/5 = 14\,\%$ ; pour $E_B = \SI{2}{V}$, elle atteint $35\,\%$ !).
\end{attention}

\begin{exemplebox}[Transistor 2N2222 en commutation : étude complète]
Soit le montage avec $E = \SI{12}{V}$, $R_C = \SI{1}{k\ohm}$, un 2N2222 de gain $\beta = 100$, et un courant de base $I_B = \SI{20}{\micro A}$. Le transistor est-il saturé ?
\begin{itemize}
\item Courant « demandé » en régime linéaire : $I_C = \beta I_B = 100 \times \num{20e-6} = \SI{2}{mA}$.
\item Courant maximal du circuit : $I_{C,\text{sat}} = \dfrac{E - V_{CE,\text{sat}}}{R_C} = \dfrac{12 - 0{,}2}{1000} = \SI{11,8}{mA}$.
\end{itemize}
Comme $\beta I_B = \SI{2}{mA} < \SI{11,8}{mA}$, le transistor est en \textbf{régime linéaire} : $I_C = \SI{2}{mA}$ et $V_{CE} = E - R_C I_C = 12 - 2 = \SI{10}{V}$. Pour saturer, il faudrait $I_B \geq \dfrac{I_{C,\text{sat}}}{\beta} = \SI{118}{\micro A}$ ; en pratique on impose $I_B \approx 2$ à $5$ fois cette valeur (sursaturation) pour garantir la commutation malgré la dispersion de $\beta$. En saturation, la puissance dissipée par le transistor est remarquablement faible :
\[ P_T = V_{CE,\text{sat}} \times I_{C,\text{sat}} = 0{,}2 \times \num{11,8e-3} \approx \SI{2,4}{mW}, \]
très inférieure à la limite du \emph{datasheet} ($P_{\max} \approx \SI{500}{mW}$ pour un 2N2222 boîtier TO-92). C'est tout l'intérêt de la commutation : l'interrupteur fermé ne chauffe presque pas, alors qu'au point $Q_{\text{lin}}$ ci-dessus il dissiperait $P_T = 10 \times \num{2e-3} = \SI{20}{mW}$, et jusqu'à $P_T = \dfrac{E^2}{4R_C} = \SI{36}{mW}$ au milieu de la droite de charge.
\end{exemplebox}

\courssub{Vérification systématique des puissances dissipées}

Déterminer le point de fonctionnement ne suffit pas : il faut encore s'assurer qu'aucun composant ne sera détruit. Tout \emph{datasheet} indique des valeurs maximales à ne jamais dépasser.

\begin{aretenir}[Bilan de puissance à vérifier pour tout montage]
Une fois le point de fonctionnement $Q(U_0, I_0)$ déterminé :
\begin{itemize}
\item \textbf{Résistor} : $P_R = R I_0^2 = \dfrac{U_R^2}{R} < P_{R,\max}$ (souvent \SI{0,25}{W} ou \SI{0,5}{W}) ;
\item \textbf{Diode / LED} : $P_D = U_0 I_0 < P_{D,\max}$ et $I_0 < I_{F,\max}$ ;
\item \textbf{Transistor} : $P_T = V_{CE}\, I_C < P_{T,\max}$, $I_C < I_{C,\max}$, $V_{CE} < V_{CEO}$ ;
\item \textbf{Bilan global} : la puissance fournie par la source $E I_0$ égale la somme des puissances reçues (loi de Pouillet généralisée au circuit).
\end{itemize}
En électronique de puissance (alimentations, variateurs), on ajoute un \emph{dissipateur thermique} si $P_T$ devient notable.
\end{aretenir}

\begin{exercice}[Type Probatoire : point de fonctionnement d'une diode]
On branche en série un générateur $E = \SI{6}{V}$, un résistor $R = \SI{200}{\ohm}$ et une diode dont la caractéristique relevée en TP est :
\begin{center}
\begin{tabular}{|l|c|c|c|c|c|c|}
\hline
\rowcolor{popblueL} $U_D$ (V) & 0 & 0{,}4 & 0{,}5 & 0{,}6 & 0{,}65 & 0{,}7 \\
\hline
$I$ (mA) & 0 & 0{,}5 & 2 & 8 & 15 & 28 \\
\hline
\end{tabular}
\end{center}
\begin{enumerate}
\item Écrire l'équation de la droite de charge et la placer sur le graphe $I(U_D)$.
\item Déterminer graphiquement le point de fonctionnement $Q$.
\item Calculer la puissance dissipée par la diode et celle dissipée par le résistor. Vérifier le bilan de puissance.
\end{enumerate}
\end{exercice}

\begin{corrige}[Exercice type Probatoire]
\textbf{1.} Loi des mailles : $I = \dfrac{E - U_D}{R} = \dfrac{6 - U_D}{200}$ (en ampères), soit $I(\text{mA}) = 30 - 5\,U_D(\text{V})$. Points de tracé : $(0\,;\,\SI{30}{mA})$ et $(\SI{6}{V}\,;\,0)$.

\textbf{2.} On cherche l'intersection avec la caractéristique. Essayons $U_D = \SI{0,65}{V}$ : la droite de charge donne $I = 30 - 5\times 0{,}65 = \SI{26,75}{mA}$, alors que la diode ne passe que \SI{15}{mA} : on est au-dessus de la courbe. Essayons $U_D = \SI{0,7}{V}$ : la droite donne $I = 30 - 3{,}5 = \SI{26,5}{mA}$, alors que la diode passe \SI{28}{mA} : on est sous la courbe. L'intersection se situe donc entre ces deux valeurs. Par interpolation graphique, on estime $U_{D0} \approx \SI{0,69}{V}$ et $I_0 \approx \SI{27}{mA}$. Le point de fonctionnement est $Q\,(\SI{0,69}{V}\,;\,\SI{27}{mA})$ environ.

\textbf{3.} Puissances (avec les valeurs interpolées) :
\[ P_D = 0{,}69 \times \num{27e-3} \approx \SI{18,6}{mW} \;;\quad P_R = 200 \times (\num{27e-3})^2 \approx \SI{146}{mW}. \]
Bilan : $P_{\text{source}} = 6 \times \num{27e-3} = \SI{162}{mW} \approx 18{,}6 + 146 = \SI{165}{mW}$ (l'écart provient de la lecture graphique). Le bilan est vérifié : l'énergie fournie par la source se répartit intégralement entre le résistor et la diode.
\end{corrige}

\begin{savaistu}[Des LED partout au Cameroun]
L'électrification rurale s'appuie de plus en plus sur des kits solaires domestiques : un petit panneau, une batterie \SI{12}{V}, et des lampes à LED. Le dimensionnement exact que tu viens d'apprendre — choix de $R$ ou d'un régulateur pour fixer le point de fonctionnement de chaque LED, vérification des puissances — est celui qu'effectuent les techniciens qui conçoivent ces kits. Une LED bien polarisée éclaire plus de dix ans ; mal polarisée (courant excessif), elle grille en quelques semaines. La physique du point de fonctionnement éclaire littéralement les villages !
\end{savaistu}

\coursec{Synthèse et préparation au Probatoire}

Dans la section précédente, tu as appris à déterminer le point de fonctionnement de circuits simples : un générateur et un résistor, un générateur et une diode, puis un étage à transistor. Tu as vu que la méthode graphique — intersection de la droite de charge avec la caractéristique du dipôle — est l'outil universel, tandis que le calcul analytique est plus rapide quand tous les composants sont linéaires. Le moment est venu de consolider : cette dernière section te donne une méthode récapitulative infaillible, un arbre de décision pour choisir la bonne stratégie, la liste des pièges classiques du Probatoire, et un aperçu des applications concrètes qui donnent tout son sens à cette leçon.

\courssub{La démarche complète en cinq étapes}

Face à n'importe quel exercice de point de fonctionnement, le candidat qui réussit est celui qui suit toujours le même plan de bataille. L'intuition est simple : un circuit série impose une seule contrainte commune (le même courant traverse tous les dipôles, et les tensions s'ajoutent) ; le point de fonctionnement est le seul couple $(I, U)$ qui satisfait à la fois la loi du générateur et la loi du récepteur. Tout le reste n'est que technique de résolution.

\begin{methode}[Checklist : les 5 points pour tout problème de point de fonctionnement]
\begin{enumerate}
\item \textbf{Identifier les composants.} Liste les dipôles du circuit : générateur de f.é.m. $E$ et de résistance interne $r$, résistor $R$, diode (tension de seuil $V_S$), transistor (jonction base-émetteur $V_{BE} \approx \SI{0,7}{V}$, gain $\beta$). Note les valeurs numériques avec leurs unités SI.
\item \textbf{Tracer (ou écrire) les caractéristiques.} Pour chaque dipôle non linéaire, la caractéristique $I(U)$ est donnée par un graphique ou une table de mesures. Pour un résistor, c'est la droite $I = U/R$.
\item \textbf{Construire la droite (ou courbe) de charge.} Applique la loi des mailles au circuit complet. Pour un générateur $(E, r)$ alimentant un dipôle $\mathcal{D}$ :
\[ E = rI + U_{\mathcal{D}} \quad \Longrightarrow \quad I = \frac{E - U_{\mathcal{D}}}{r}. \]
C'est une droite décroissante reliant les points remarquables $(U=0\,;\,I = E/r)$ et $(U = E\,;\,I = 0)$.
\item \textbf{Déterminer l'intersection : le point de fonctionnement.} Graphiquement, lis les coordonnées $(U_0, I_0)$ du point d'intersection ; analytiquement, résous le système de deux équations à deux inconnues. Vérifie la cohérence (signes, ordres de grandeur).
\item \textbf{Calculer les puissances et l'énergie (loi de Pouillet).} Puissance fournie par le générateur $P_g = E I_0$, puissance utile $P_u = U_0 I_0$, puissance dissipée par effet Joule dans le générateur $P_J = r I_0^2$ et dans un résistor $P = R I_0^2$. Vérifie le bilan : $P_g = P_u + P_J$.
\end{enumerate}
\end{methode}

\begin{attention}[Le bilan de puissances est ton meilleur contrôle]
Au Probatoire, une erreur de lecture graphique se paie cash... sauf si tu vérifies le bilan énergétique. La loi de Pouillet impose $E I_0 = U_0 I_0 + r I_0^2$, c'est-à-dire $E = U_0 + r I_0$ : tu retrouves la loi des mailles. Si ton point de fonctionnement lu sur le graphique ne vérifie pas cette relation à quelques pour cent près, c'est que ta lecture est fausse. Prends l'habitude de faire ce contrôle systématiquement : il prend vingt secondes et peut sauver tout l'exercice.
\end{attention}

\courssub{Arbre de décision : quelle méthode choisir ?}

Tous les circuits ne se traitent pas de la même façon. La question clé est la \emph{linéarité} : un composant est linéaire si sa caractéristique $I(U)$ est une droite passant par l'origine (loi d'Ohm). Dans ce cas, l'algèbre suffit. Sinon, il faut le graphique.

\begin{popfigure}
\begin{center}
\begin{tikzpicture}[
  font=\small,
  decision/.style={diamond, draw=popdark, fill=poppinkL, text width=3.2cm, align=center, inner sep=1pt, aspect=2},
  bloc/.style={rectangle, rounded corners, draw=popdark, fill=popblueL, text width=3.6cm, align=center, inner sep=4pt},
  fleche/.style={-{Stealth[length=2.5mm]}, thick, popdark}
]
\node[decision] (q1) {Caractéristique linéaire ?};
\node[bloc, right=2.2cm of q1] (ohm) {Tous dipôles linéaires : loi d'Ohm directe, calcul analytique $I = \dfrac{E}{r+R}$};
\node[decision, below=1.4cm of q1] (q2) {Diode ou dipôle non linéaire ?};
\node[bloc, right=2.2cm of q2] (graph) {Méthode graphique : caractéristique $I(U)$ + droite de charge $\Rightarrow$ intersection};
\node[decision, below=1.4cm of q2] (q3) {Transistor présent ?};
\node[bloc, right=2.2cm of q3] (trans) {Réseau de courbes de sortie $I_C(U_{CE})$ paramétré par $I_B$ + droite de charge $I_C = \dfrac{V_{CC}-U_{CE}}{R_C}$};
\node[bloc, below=1.4cm of q3, fill=popgreenL] (pf) {Point de fonctionnement $(U_0, I_0)$ puis puissances (loi de Pouillet)};

\draw[fleche] (q1) -- node[above]{oui} (ohm);
\draw[fleche] (q1) -- node[left]{non} (q2);
\draw[fleche] (q2) -- node[above]{oui} (graph);
\draw[fleche] (q2) -- node[left]{non} (q3);
\draw[fleche] (q3) -- node[above]{oui} (trans);
\draw[fleche] (ohm.south) |- (pf.east);
\draw[fleche] (graph.south) |- (pf.east);
\draw[fleche] (trans.south) |- (pf.east);
\draw[fleche] (q3) -- node[left]{cas général} (pf);
\end{tikzpicture}
\end{center}
\legende{Arbre de décision pour le choix de la méthode de détermination du point de fonctionnement.}
\end{popfigure}

\begin{aretenir}[Règle d'or du choix de méthode]
\begin{itemize}
\item Circuit composé uniquement de \cle{générateurs} et de \cle{résistors} : méthode \textbf{analytique} (loi d'Ohm généralisée, loi de Pouillet pour l'intensité $I = \dfrac{E}{r + R_{\text{éq}}}$).
\item Présence d'une \cle{diode} ou de tout dipôle à caractéristique courbe : méthode \textbf{graphique} (droite de charge), ou modèle approché de la diode (source de tension $V_S$ en série avec $r_d$) si l'énoncé le fournit.
\item Présence d'un \cle{transistor} : méthode graphique sur le réseau de caractéristiques de sortie, après avoir calculé le courant de base $I_B$ sur la maille d'entrée.
\end{itemize}
\end{aretenir}

\courssub{Points de vigilance : les pièges du Probatoire}

Chaque année, les rapports du jury signalent les mêmes erreurs. Les connaître, c'est déjà les éviter.

\begin{attention}[Pièges classiques au Probatoire]
\begin{enumerate}
\item \textbf{Confusion entre $U$ et $I$ sur les axes.} La caractéristique d'un dipôle est tracée avec $U$ en abscisse et $I$ en ordonnée. La droite de charge s'écrit $I = \dfrac{E - U}{r}$ : son intersection avec l'axe des ordonnées est $E/r$ (courant de court-circuit), avec l'axe des abscisses $E$ (tension à vide). Inverser les axes fausse toute la lecture.
\item \textbf{Oubli de la tension $V_{BE}$ du transistor.} Dans la maille d'entrée, le courant de base est $I_B = \dfrac{V_{CC} - V_{BE}}{R_B}$ et \emph{non} $V_{CC}/R_B$. Avec $V_{CC} = \SI{6}{V}$ et $R_B = \SI{100}{k\Omega}$, oublier $V_{BE} = \SI{0,7}{V}$ donne $I_B = \SI{60}{\micro A}$ au lieu de $\SI{53}{\micro A}$ : une erreur de $13\,\%$ qui se propage sur tout l'exercice via $I_C = \beta I_B$.
\item \textbf{Mauvaise lecture d'échelle.} Vérifie toujours les unités des axes : l'ordonnée est souvent en $\si{mA}$ ou $\si{\micro A}$, l'abscisse parfois graduée de $\SI{2}{V}$ en $\SI{2}{V}$. Si l'axe est logarithmique (caractéristique de diode en semi-log), chaque graduation multiplie la valeur par 10 : on ne lit \emph{jamais} par interpolation linéaire entre deux décades.
\item \textbf{Unités incohérentes dans les calculs de puissance.} La loi de Pouillet $P = RI^2$ exige $I$ en ampères : avec $I = \SI{20}{mA}$, il faut écrire $I = \SI{0,020}{A}$, sinon la puissance est fausse d'un facteur $10^6$.
\item \textbf{Oubli des limites de dissipation.} Un résistor de puissance nominale $P_{\max} = \SI{0,25}{W}$ ne peut pas dissiper $\SI{1}{W}$ : il grille. Vérifie toujours que $R I_0^2 \leq P_{\max}$ ; c'est souvent la dernière question de l'exercice.
\item \textbf{Polarités.} Une diode montée en inverse ne conduit pas : le point de fonctionnement est alors $(I_0 = 0)$ et toute la tension se retrouve à ses bornes. Repère le sens du courant avant tout calcul.
\end{enumerate}
\end{attention}

\courssub{Exercice de synthèse résolu}

\begin{exoresolu}[Alimentation d'une lampe par une batterie de mototaxi]
Une batterie de mototaxi de f.é.m. $E = \SI{12,0}{V}$ et de résistance interne $r = \SI{0,50}{\Omega}$ alimente une lampe dont la caractéristique passe par les points $(\SI{4}{V}\,;\,\SI{1,0}{A})$, $(\SI{8}{V}\,;\,\SI{1,6}{A})$, $(\SI{10}{V}\,;\,\SI{1,8}{A})$, $(\SI{11}{V}\,;\,\SI{1,9}{A})$.
\begin{enumerate}
\item Écrire l'équation de la droite de charge et donner ses points remarquables.
\item Déterminer graphiquement le point de fonctionnement. Vérifier qu'un couple $(U_0 \approx \SI{10,4}{V}\,;\, I_0 \approx \SI{3,2}{A})$ n'est pas acceptable.
\item Calculer la puissance fournie par la batterie, la puissance reçue par la lampe et le rendement de l'alimentation.
\end{enumerate}
\tcblower
\textbf{1. Droite de charge.} La loi des mailles donne $E = rI + U$, soit :
\[ I = \frac{E - U}{r} = \frac{12,0 - U}{0,50}. \]
Points remarquables : court-circuit $(U = 0\,;\,I = E/r = \SI{24}{A})$ et circuit ouvert $(U = E = \SI{12,0}{V}\,;\,I = 0)$.

\textbf{2. Point de fonctionnement.} Le couple $(\SI{10,4}{V}\,;\,\SI{3,2}{A})$ proposé est impossible : la droite de charge impose $I = (12,0 - 10,4)/0,50 = \SI{3,2}{A}$, ce qui est cohérent avec la droite, mais la caractéristique de la lampe donne à $\SI{10,4}{V}$ un courant voisin de $\SI{1,85}{A}$ (interpolation entre $\SI{10}{V}$ et $\SI{11}{V}$), très inférieur à $\SI{3,2}{A}$. Le point doit être sur \emph{les deux} courbes. Cherchons l'intersection : sur la droite de charge, à $I = \SI{1,8}{A}$ on a $U = 12,0 - 0,50 \times 1,8 = \SI{11,1}{V}$ ; or la lampe donne $U = \SI{10}{V}$ pour $\SI{1,8}{A}$. À $I = \SI{1,9}{A}$ : droite de charge $U = 12,0 - 0,95 = \SI{11,05}{V}$, lampe $U = \SI{11}{V}$. Les deux courbes se croisent donc pratiquement en :
\[ \boxed{U_0 \approx \SI{11,0}{V} \quad ; \quad I_0 \approx \SI{1,9}{A}}. \]
Vérification par la loi des mailles : $r I_0 + U_0 = 0,50 \times 1,9 + 11,0 = \SI{11,95}{V} \approx E$. Cohérent.

\textbf{3. Bilan de puissances (loi de Pouillet).}
\begin{align*}
P_g &= E I_0 = 12,0 \times 1,9 = \SI{22,8}{W}, \\
P_u &= U_0 I_0 = 11,0 \times 1,9 = \SI{20,9}{W}, \\
P_J &= r I_0^2 = 0,50 \times (1,9)^2 = \SI{1,8}{W}.
\end{align*}
Contrôle : $P_u + P_J = 20,9 + 1,8 = \SI{22,7}{W} \approx P_g$. Rendement :
\[ \eta = \frac{P_u}{P_g} = \frac{20,9}{22,8} \approx 0,92 \quad \text{soit } 92\,\%. \]
La batterie, de faible résistance interne, alimente la lampe avec un excellent rendement.
\end{exoresolu}

\courssub{Applications concrètes autour de nous}

\textbf{Driver de LED.} Une LED blanche d'éclairage (téléphone, lampe torche rechargeable vendue au marché) présente une caractéristique très raide : sa tension de seuil est voisine de $\SI{3,2}{V}$ et une petite augmentation de tension provoque une forte augmentation de courant, qui peut la détruire. On l'alimente donc à travers un résistor de limitation : la droite de charge fixe le point de fonctionnement à $I_0 = \SI{20}{mA}$, valeur nominale. Avec une alimentation de $\SI{5}{V}$ : $R = \dfrac{5 - 3,2}{0,020} = \SI{90}{\Omega}$, et la puissance dissipée $P = RI_0^2 = \SI{0,036}{W}$ reste très inférieure à la limite $\SI{0,25}{W}$ d'un résistor courant.

\textbf{Étage de puissance audio.} Dans un amplificateur (enceinte de sonorisation d'un bar de Yaoundé), le transistor de sortie est polarisé par sa droite de charge tracée sur le réseau $I_C(U_{CE})$. Le point de repos est choisi au milieu de la droite de charge pour que le signal audio puisse faire osciller le point de fonctionnement symétriquement sans écrêtage : c'est la condition d'un son fidèle, sans distorsion. La loi de Pouillet $P = U_{CE} I_C$ donne la puissance dissipée au repos, qui dimensionne le radiateur de refroidissement.

\textbf{Commande de moteur à courant continu (aperçu).} Un petit moteur (ventilateur, pompe d'irrigation solaire) se comporte comme une résistance $R_M$ en série avec une f.c.é.m. $E'$ proportionnelle à sa vitesse de rotation. La loi des mailles $U = E' + R_M I$ montre que le courant — donc le couple moteur — se règle par la tension d'alimentation. Au démarrage, $E' = 0$ et le courant $I = U/R_M$ est maximal : c'est pourquoi on insère un résistor de démarrage, retiré ensuite. Tu retrouveras cela en classe de Terminale.

\begin{savaistu}[La loi de Pouillet sur les lignes haute tension du Cameroun]
Les ingénieurs de SONATREL (ex-SONEL), qui exploitent le réseau de transport d'électricité du Cameroun, utilisent quotidiennement la loi de Pouillet. Entre la centrale de Song Loulou et Douala, l'énergie transite par des lignes haute tension de plusieurs dizaines de kilomètres. La puissance perdue par effet Joule dans une ligne de résistance $R$ est $P_J = RI^2$ : pour transporter une puissance $P = UI$ donnée, on a intérêt à \emph{augmenter la tension} $U$ afin de \emph{diminuer le courant} $I$, car les pertes varient comme le \emph{carré} de l'intensité. C'est pourquoi l'électricité voyage sous $\SI{90}{kV}$ ou $\SI{225}{kV}$ sur les pylônes que tu vois le long des routes, avant d'être abaissée à $\SI{220}{V}$ dans les quartiers par les transformateurs. Avec les nouveaux barrages de Nachtigal et de Lom Pangar, le dimensionnement des lignes d'évacuation — choix des câbles, échauffement admissible, point de fonctionnement du réseau — repose exactement sur les outils de cette leçon.
\end{savaistu}

\courssub{Pour finir : l'essentiel à retenir}

\begin{aretenir}[L'essentiel de la leçon]
\begin{itemize}
\item Le \cle{point de fonctionnement} d'un circuit est le couple $(U_0, I_0)$, intersection de la caractéristique du dipôle récepteur et de la \cle{droite de charge} $I = \dfrac{E - U}{r}$ imposée par le générateur.
\item Choix de méthode : \textbf{analytique} si tous les dipôles sont linéaires (loi de Pouillet $I = \dfrac{E}{r + R}$), \textbf{graphique} dès qu'apparaît une diode, un transistor ou toute caractéristique courbe.
\item La \cle{loi de Pouillet} pour les puissances : $P_g = EI_0$, $P_u = U_0 I_0$, $P_J = rI_0^2$ (ou $RI_0^2$), avec le bilan $P_g = P_u + P_J$ qui sert de vérification systématique.
\item Vigilance permanente : unités SI (courants en ampères), polarités des composants, tension $V_{BE} \approx \SI{0,7}{V}$ non négligeable, lecture soigneuse des échelles, et respect des puissances maximales admissibles.
\end{itemize}
Entraîne-toi sur les exercices de la série : au Probatoire, un exercice de point de fonctionnement bien mené, avec schéma, droite de charge propre et bilan de puissances vérifié, est l'un des sujets où la méthode paie le plus. Tu as désormais toutes les cartes en main.
\end{aretenir}

\newpage

\coursec{Activité d'intégration}

\courssub{Objectifs de l'activité}

À l'issue de cette activité d'intégration, tu seras capable de :
\begin{itemize}
\item mobiliser la notion de \cle{point de fonctionnement} d'un circuit série simple et le déterminer par deux méthodes complémentaires : calcul (loi des mailles) et méthode graphique (droite de charge) ;
\item appliquer la \cle{loi de Pouillet} et les bilans de puissance $P = UI$ à un dipôle réel ;
\item choisir un composant (résistance de la série normalisée E12) en fonction d'un cahier des charges énergétique ;
\item confronter les prévisions théoriques à des mesures expérimentales et discuter les écarts ;
\item travailler en équipe, rédiger un rapport synthétique et argumenter un choix technique.
\end{itemize}

\courssub{Situation}

\textbf{Un indicateur lumineux pour le groupe électrogène du lycée.}

Le lycée de Ngaoundéré, comme beaucoup d'établissements camerounais, subit des délestages fréquents. Le surveillant général utilise alors un petit groupe électrogène, mais celui-ci ne possède aucun voyant indiquant qu'il est sous tension : plusieurs fois, des élèves ont approché les câbles sans savoir que le courant était rétabli. Le club scientifique du lycée décide de fabriquer un \cle{indicateur lumineux à LED} basse consommation, alimenté par une pile de $9\ \text{V}$ de récupération, qui restera allumé en permanence sur le tableau de commande.

Le technicien du club fournit le matériel suivant :
\begin{itemize}
\item une pile de $9\ \text{V}$, modélisée par un générateur idéal de f.é.m. $E = 9{,}0\ \text{V}$ en série avec une résistance interne $r = 10\ \Omega$ ;
\item des LED rouges dont la notice indique : tension de seuil (tension directe) $V_f \approx 2{,}0\ \text{V}$, courant nominal $I_f = 20\ \text{mA}$, courant maximal admissible $I_{\max} = 30\ \text{mA}$ ;
\item un assortiment de résistances de la série normalisée E12, de puissance maximale admissible $P_{\max} = 0{,}25\ \text{W}$ : valeurs disponibles (en $\Omega$) : 100 ; 120 ; 150 ; 180 ; 220 ; 270 ; 330 ; 390 ; 470 ; 560 ; 680 ; 820 ;
\item un multimètre, des fils de connexion, une plaque d'essai.
\end{itemize}

Pour la modélisation, on admettra que la LED se comporte comme un dipôle dont la caractéristique est assimilable à :
\begin{itemize}
\item un interrupteur ouvert si $U < V_f$ (aucun courant) ;
\item une source de tension parfaite de valeur $V_f = 2{,}0\ \text{V}$ si elle conduit (modèle dit \og tension de seuil \fg{}).
\end{itemize}

Le cahier des charges impose :
\begin{enumerate}
\item un courant dans la LED \textbf{proche de} $20\ \text{mA}$ (luminosité nominale) sans jamais dépasser $30\ \text{mA}$ ;
\item une puissance dissipée dans la résistance \textbf{strictement inférieure à} $0{,}25\ \text{W}$ (sinon la résistance chauffe et se détruit) ;
\item une consommation la plus faible possible pour préserver la pile.
\end{enumerate}

\courssub{Consignes}

Travail en groupes de trois élèves. Chaque groupe rédige un rapport synthétique (une page maximum) répondant aux questions suivantes.

\begin{enumerate}
\item \textbf{Schéma et modélisation.} Dessiner le schéma du circuit série (pile réelle, résistance de protection $R$, LED) en convention récepteur pour les dipôles passifs. Écrire la loi des mailles et montrer que le courant vérifie $E = rI + RI + V_f$.
\item \textbf{Choix de la résistance.} En déduire la valeur théorique de $R$ permettant d'obtenir $I = 20\ \text{mA}$. Choisir dans la série E12 la valeur normalisée la plus proche et justifier ce choix. Calculer le courant réellement obtenu avec cette valeur et vérifier qu'il reste inférieur à $I_{\max} = 30\ \text{mA}$.
\item \textbf{Détermination graphique du point de fonctionnement.} Sur un graphe $I$ en fonction de $U$ (tension aux bornes de l'association LED, axe horizontal gradué de $0$ à $9\ \text{V}$, courant de $0$ à $40\ \text{mA}$) :
\begin{enumerate}
\item tracer la caractéristique de la LED (modèle à seuil) ;
\item tracer la \cle{droite de charge} du générateur réel associé à la résistance $R$ choisie, c'est-à-dire la courbe $I = \dfrac{E - U}{R + r}$ ;
\item localiser le point de fonctionnement $F$ et lire graphiquement $U_F$ et $I_F$. Comparer avec les valeurs calculées à la question 2.
\end{enumerate}
\item \textbf{Bilan énergétique (loi de Pouillet).} Calculer :
\begin{enumerate}
\item la puissance $P_R = RI^2$ dissipée par effet Joule dans la résistance ; conclure quant au risque de surchauffe ;
\item la puissance $P_{\text{LED}} = V_f \cdot I$ reçue par la LED ;
\item la puissance totale $P_g = EI$ fournie par la pile et la puissance $P_r = rI^2$ perdue dans sa résistance interne ;
\item vérifier la conservation de la puissance : $P_g = P_r + P_R + P_{\text{LED}}$.
\end{enumerate}
\item \textbf{Mesures expérimentales.} Monter le circuit, mesurer à l'aide du multimètre la tension réelle $U_{\text{LED}}$ aux bornes de la LED et le courant réel $I_{\text{réel}}$. Comparer aux prévisions et proposer une explication des écarts éventuels (tolérance des résistances E12 : $\pm 10\ \%$, tension de seuil réelle de la LED comprise entre $1{,}8$ et $2{,}2\ \text{V}$, usure de la pile).
\item \textbf{Prolongement.} Le surveillant propose de remplacer la pile par deux piles de $9\ \text{V}$ en série ($18\ \text{V}$). Sans refaire tous les calculs, dire qualitativement si la même résistance convient encore, puis le vérifier par un calcul rapide du courant.
\end{enumerate}

\courssub{Ressources à mobiliser}

\begin{aretenir}[Boîte à outils de la leçon]
\begin{itemize}
\item \textbf{Loi des mailles} : la somme algébrique des tensions le long d'une maille fermée est nulle.
\item \textbf{Générateur réel} : $U_{PN} = E - rI$ (convention générateur).
\item \textbf{Point de fonctionnement} : point d'intersection de la caractéristique $I(U)$ du dipôle et de la droite de charge du circuit d'alimentation ; il donne le couple $(U_F, I_F)$ réellement établi.
\item \textbf{Droite de charge} d'un générateur $(E, r)$ en série avec une résistance $R$ : $I = \dfrac{E - U}{R + r}$ ; elle passe par les points remarquables $(U = E\,;\, I = 0)$ et $(U = 0\,;\, I = \dfrac{E}{R+r})$.
\item \textbf{Loi de Pouillet} (puissance) : puissance reçue par un dipôle en convention récepteur : $P = UI$ ; puissance dissipée par effet Joule dans un résistor : $P = RI^2 = \dfrac{U^2}{R}$.
\item \textbf{Bilan de puissance} : la puissance fournie par le générateur égale la somme des puissances reçues par tous les dipôles du circuit.
\end{itemize}
\end{aretenir}

\courssub{Démarche et corrigé commenté}

\begin{exoresolu}[Corrigé commenté de l'activité]
\textbf{Énoncé résumé.} Pile $E = 9{,}0\ \text{V}$, $r = 10\ \Omega$ ; LED rouge ($V_f = 2{,}0\ \text{V}$, $I_f = 20\ \text{mA}$, $I_{\max} = 30\ \text{mA}$) ; résistances E12 de $0{,}25\ \text{W}$. Déterminer $R$, le point de fonctionnement, le bilan de puissance, et discuter les mesures.
\tcblower
\textbf{1. Schéma et loi des mailles.} Le circuit série comporte le générateur réel $(E, r)$, la résistance $R$ et la LED, tous traversés par le même courant $I$. En parcourant la maille dans le sens du courant :
\[ E - rI - RI - V_f = 0 \quad \Longrightarrow \quad E = (R + r)I + V_f. \]

\textbf{2. Choix de la résistance.} Pour $I = 20\ \text{mA} = 0{,}020\ \text{A}$ :
\[ R = \frac{E - V_f}{I} - r = \frac{9{,}0 - 2{,}0}{0{,}020} - 10 = 350 - 10 = 340\ \Omega. \]
La valeur E12 la plus proche est $R = 330\ \Omega$ (la valeur $390\ \Omega$ donnerait un courant trop faible). Avec $R = 330\ \Omega$ :
\[ I = \frac{E - V_f}{R + r} = \frac{7{,}0}{340} \approx 0{,}0206\ \text{A} \approx 21\ \text{mA} < 30\ \text{mA}. \]
Le cahier des charges est respecté : courant proche de $20\ \text{mA}$ et inférieur à $I_{\max}$.

\textbf{3. Point de fonctionnement graphique.} La caractéristique de la LED est la demi-droite verticale $U = V_f = 2{,}0\ \text{V}$ pour $I > 0$. La droite de charge s'écrit :
\[ I = \frac{E - U}{R + r} = \frac{9{,}0 - U}{340}\ (\text{A}). \]
Elle passe par $(U = 9{,}0\ \text{V}\,;\, I = 0)$ et $(U = 0\,;\, I = 26{,}5\ \text{mA})$. Son intersection avec la droite $U = 2{,}0\ \text{V}$ donne :
\[ F\,(U_F = 2{,}0\ \text{V}\ ;\ I_F \approx 21\ \text{mA}), \]
en parfait accord avec le calcul de la question 2. \emph{Remarque :} avec une LED réelle, la caractéristique n'est pas rigoureusement verticale ; le point $F$ lu graphiquement peut donner $U_F \approx 2{,}1\ \text{V}$ et $I_F \approx 20\ \text{mA}$.

\textbf{4. Bilan énergétique.} Avec $I = 20{,}6\ \text{mA}$ :
\begin{align*}
P_R &= RI^2 = 330 \times (0{,}0206)^2 \approx 0{,}14\ \text{W} < 0{,}25\ \text{W} \quad \text{(pas de surchauffe)} ;\\
P_{\text{LED}} &= V_f \cdot I = 2{,}0 \times 0{,}0206 \approx 0{,}041\ \text{W} ;\\
P_g &= EI = 9{,}0 \times 0{,}0206 \approx 0{,}185\ \text{W} ;\\
P_r &= rI^2 = 10 \times (0{,}0206)^2 \approx 0{,}004\ \text{W}.
\end{align*}
Vérification : $P_r + P_R + P_{\text{LED}} \approx 0{,}004 + 0{,}140 + 0{,}041 = 0{,}185\ \text{W} = P_g$. La conservation de la puissance est bien vérifiée. On note que la résistance dissipe environ trois fois plus de puissance que la LED : c'est le prix de la simplicité de ce montage.

\textbf{5. Mesures.} Les valeurs mesurées typiques sont $U_{\text{LED}} \approx 2{,}1\ \text{V}$ et $I_{\text{réel}} \approx 19\ \text{mA}$. Les écarts s'expliquent par la tolérance $\pm 10\ \%$ de la résistance, la dispersion de $V_f$ d'une LED à l'autre et la tension réelle de la pile, souvent inférieure à $9{,}0\ \text{V}$ si elle est usagée.

\textbf{6. Prolongement.} Avec $E' = 18\ \text{V}$ et la même résistance : $I' = \dfrac{18 - 2{,}0}{340} \approx 47\ \text{mA} > 30\ \text{mA}$ : la LED serait détruite. Il faudrait $R' = \dfrac{18 - 2{,}0}{0{,}020} - 10 = 790\ \Omega$, soit $820\ \Omega$ en E12, et vérifier $P_{R'} = 820 \times (0{,}0195)^2 \approx 0{,}31\ \text{W} > 0{,}25\ \text{W}$ : il faudrait alors une résistance de $0{,}5\ \text{W}$. Le montage à $9\ \text{V}$ est donc le meilleur choix.
\end{exoresolu}

\courssub{Bilan de l'activité}

Cette activité a permis de mettre en évidence les points essentiels de la leçon :
\begin{itemize}
\item le \cle{point de fonctionnement} n'est pas une abstraction : c'est le couple $(U_F, I_F)$ qui s'établit réellement dans le circuit, et on l'obtient indifféremment par le calcul (loi des mailles) ou graphiquement (intersection de la caractéristique et de la droite de charge) ;
\item la \cle{loi de Pouillet} et le bilan de puissance sont des outils de \textbf{conception} : ils permettent de vérifier qu'un composant ne sera pas détruit (ici $P_R < 0{,}25\ \text{W}$) et de quantifier les pertes ;
\item le \cle{choix d'un composant} réel se fait parmi des valeurs normalisées : l'ingénieur ne cherche pas la valeur exacte mais la valeur disponible qui respecte toutes les contraintes, avec une marge de sécurité ;
\item la confrontation théorie--mesure montre que les modèles (LED à seuil, pile idéale) sont des approximations utiles mais perfectibles : c'est toute la démarche scientifique ;
\item enfin, ce mini-projet illustre une compétence directement transférable : dimensionner un circuit simple, fiable et économe en énergie, une préoccupation quotidienne au Cameroun, des voyants d'appareils aux systèmes solaires domestiques.
\end{itemize}

\newpage

\coursec{Méthodes et exercices résolus}

\courssub{Fiches méthodes et exercices résolus}

\begin{methode}[Tracer et exploiter une caractéristique courant-tension $I(U)$]
\begin{enumerate}
\item \textbf{Identifier le dipôle} et relever les couples de mesures $(U\,;I)$ fournis par l'énoncé (tableau de TP ou valeurs lues).
\item \textbf{Choisir les axes} : tension $U$ en abscisse (en V), intensité $I$ en ordonnée (en A ou mA). Adapter l'échelle pour que la courbe occupe l'essentiel du graphique.
\item \textbf{Placer les points} puis tracer la courbe : droite passant par l'origine pour un résistor (loi d'Ohm), courbe avec tension de seuil pour une diode.
\item \textbf{Exploiter} : pour un résistor, le rapport $\dfrac{U}{I}$ (inverse de la pente) est constant et donne la résistance $R = \dfrac{U}{I}$ (en $\Omega$). Vérifier l'homogénéité : $1~\Omega = 1~\mathrm{V \cdot A^{-1}}$.
\item \textbf{Conclure} en précisant si le dipôle est passif/actif, symétrique/non symétrique, linéaire/non linéaire.
\end{enumerate}
\textbf{Réflexe} : toujours convertir les mA en A avant tout calcul ($1~\mathrm{mA} = 10^{-3}~\mathrm{A}$).
\end{methode}

\begin{exoresolu}[Caractéristique d'un résistor de chauffage]
Au cours d'un TP, un élève de Première C mesure, pour un résistor, les couples suivants :

\begin{center}
\begin{tabular}{|l|c|c|c|c|c|}
\hline
\rowcolor{popblueL} $U$ (V) & 0 & 2,0 & 4,0 & 6,0 & 8,0 \\
\hline
$I$ (mA) & 0 & 40 & 80 & 120 & 160 \\
\hline
\end{tabular}
\end{center}

\begin{enumerate}
\item Tracer la caractéristique $I = f(U)$ du résistor.
\item En déduire la valeur de sa résistance $R$.
\item Quelle intensité traverse ce résistor sous $U = \SI{12}{V}$ ?
\end{enumerate}
\tcblower
\textbf{1. Tracé de la caractéristique.} On porte $U$ en abscisse et $I$ en ordonnée. Les points $(0\,;0)$, $(2{,}0\,;40)$, $(4{,}0\,;80)$, $(6{,}0\,;120)$, $(8{,}0\,;160)$ sont alignés avec l'origine : la caractéristique est une \cle{droite passant par l'origine}.

\begin{popfigure}
\begin{tikzpicture}
\begin{axis}[width=0.75\linewidth,height=5.5cm,
xlabel={$U$ (V)},ylabel={$I$ (mA)},
xmin=0,xmax=9,ymin=0,ymax=190,
grid=both,grid style={popline!40},
xtick={0,2,4,6,8},ytick={0,40,80,120,160}]
\addplot[popcurve,domain=0:8.5,samples=50]{20*x};
\addplot[only marks,mark=*,mark size=1.5pt,color=poporange] coordinates {(2,40) (4,80) (6,120) (8,160)};
\end{axis}
\end{tikzpicture}
\legende{Caractéristique $I(U)$ du résistor : droite passant par l'origine, donc dipôle linéaire.}
\end{popfigure}

\textbf{2. Valeur de la résistance.} La caractéristique étant une droite passant par l'origine, le dipôle vérifie la loi d'Ohm $U = RI$. On choisit un point éloigné de l'origine pour limiter l'erreur, par exemple $(U = \SI{8,0}{V}\,;\,I = \SI{160}{mA} = \SI{0,160}{A})$ :
\[ R = \frac{U}{I} = \frac{8{,}0}{0{,}160} = \SI{50}{\ohm}. \]
Vérification dimensionnelle : $\dfrac{\mathrm{V}}{\mathrm{A}} = \Omega$. Tous les points donnent le même rapport, ce qui confirme la linéarité.

\textbf{3. Intensité sous 12 V.} D'après la loi d'Ohm :
\[ I = \frac{U}{R} = \frac{12}{50} = \SI{0,24}{A} = \SI{240}{mA}. \]
\textbf{Conclusion} : le résistor a une résistance $R = \SI{50}{\ohm}$ ; sous $\SI{12}{V}$, il est traversé par un courant d'intensité $I = \SI{0,24}{A}$.
\end{exoresolu}

\begin{methode}[Déterminer graphiquement le point de fonctionnement]
\begin{enumerate}
\item \textbf{Écrire l'équation de la droite de charge} (loi des mailles) pour le circuit générateur $(E, r)$ + dipôle : $U = E - rI$, soit $I = \dfrac{E - U}{r}$.
\item \textbf{Tracer sur le même graphique} la caractéristique $I(U)$ du dipôle ($U$ en abscisse, $I$ en ordonnée) et la droite de charge. Deux points suffisent pour la droite : $(U = 0\,;\,I = E/r)$ et $(U = E\,;\,I = 0)$.
\item \textbf{Lire les coordonnées du point d'intersection} $F(U_F\,;\,I_F)$ : c'est le \cle{point de fonctionnement}, seul couple compatible avec les deux dipôles.
\item \textbf{Vérifier} le résultat par le calcul si les lois sont connues (ex. résistor : résoudre $E - rI = RI$).
\item \textbf{Conclure} en écrivant $U_F$ et $I_F$ avec leurs unités et un nombre cohérent de chiffres significatifs.
\end{enumerate}
\textbf{Réflexe} : la droite de charge ne dépend que du générateur ; la caractéristique ne dépend que du dipôle récepteur. Le point de fonctionnement est leur compromis.
\end{methode}

\begin{exoresolu}[Point de fonctionnement d'une pile et d'un résistor]
Une pile de f.é.m. $E = \SI{9,0}{V}$ et de résistance interne $r = \SI{10}{\ohm}$ alimente un résistor de résistance $R = \SI{80}{\ohm}$.
\begin{enumerate}
\item Établir l'équation de la droite de charge du circuit.
\item Déterminer le point de fonctionnement $F(U_F\,;\,I_F)$ par le calcul, puis vérifier graphiquement.
\item Calculer la tension aux bornes de la pile. Commenter.
\end{enumerate}
\tcblower
\textbf{1. Droite de charge.} La loi des mailles (ou loi d'Ohm aux bornes du générateur) donne la tension aux bornes de la pile :
\[ U = E - rI = 9{,}0 - 10\,I \quad (U \text{ en V},\, I \text{ en A}). \]
C'est l'équation d'une droite décroissante passant par les points $(U = 0\,;\,I = 0{,}90~\mathrm{A})$ et $(U = 9{,}0~\mathrm{V}\,;\,I = 0)$.

\textbf{2. Point de fonctionnement.} Le résistor impose $U = RI = 80\,I$. Le point de fonctionnement vérifie simultanément les deux relations :
\begin{align*}
80\,I_F &= 9{,}0 - 10\,I_F \\
90\,I_F &= 9{,}0 \\
I_F &= \frac{9{,}0}{90} = \SI{0,10}{A}.
\end{align*}
On en déduit $U_F = RI_F = 80 \times 0{,}10 = \SI{8,0}{V}$.

Vérification graphique : la droite de charge $I = \frac{9{,}0 - U}{10}$ et la caractéristique du résistor $I = \frac{U}{80}$ se coupent en $F(U_F = 8{,}0~\mathrm{V}\,;\,I_F = 0{,}10~\mathrm{A})$.

\begin{popfigure}
\begin{tikzpicture}
\begin{axis}[width=0.75\linewidth,height=6cm,
xlabel={$U$ (V)},ylabel={$I$ (A)},
xmin=0,xmax=10,ymin=0,ymax=1,
grid=both,grid style={popline!40},
legend style={font=\small,at={(0.98,0.97)},anchor=north east}]
\addplot[popcurve,domain=0:9,samples=50]{(9-x)/10};
\addlegendentry{Droite de charge $I = (E-U)/r$}
\addplot[popcurve,poporange,domain=0:8.5,samples=50]{x/80};
\addlegendentry{Résistor $I = U/R$}
\addplot[only marks,mark=*,mark size=2pt,color=popdark] coordinates {(8,0.1)};
\node[anchor=south west,font=\small] at (axis cs:8,0.1) {$F(8{,}0\,;\,0{,}10)$};
\end{axis}
\end{tikzpicture}
\legende{Le point de fonctionnement $F$ est l'intersection de la droite de charge et de la caractéristique du résistor (axes $U$ horizontal, $I$ vertical).}
\end{popfigure}

\textbf{3. Tension aux bornes de la pile.} $U_{PN} = E - rI_F = 9{,}0 - 10 \times 0{,}10 = \SI{8,0}{V}$. La tension disponible ($\SI{8,0}{V}$) est inférieure à la f.é.m. ($\SI{9,0}{V}$) à cause de la chute de tension interne $rI_F = \SI{1,0}{V}$.

\textbf{Conclusion} : le point de fonctionnement est $F(U_F = \SI{8,0}{V}\,;\,I_F = \SI{0,10}{A})$ ; la pile débite $\SI{100}{mA}$ sous $\SI{8,0}{V}$.
\end{exoresolu}

\begin{methode}[Appliquer la loi de Pouillet (puissance et énergie)]
\begin{enumerate}
\item \textbf{Identifier le régime} : la loi de Pouillet $\mathcal{P} = RI^2 = \dfrac{U^2}{R}$ ne s'applique qu'à un \cle{résistor} (effet Joule pur).
\item \textbf{Choisir la formule adaptée aux données} : $\mathcal{P} = RI^2$ si l'on connaît $I$ ; $\mathcal{P} = \dfrac{U^2}{R}$ si l'on connaît $U$ ; $\mathcal{P} = UI$ dans tous les cas pour un dipôle quelconque.
\item \textbf{Calculer l'énergie} sur une durée $\Delta t$ : $W = \mathcal{P}\,\Delta t$, avec $W$ en joules (J) si $\mathcal{P}$ en W et $\Delta t$ en s. Pour l'énergie électrique domestique : $W$ en kWh si $\mathcal{P}$ en kW et $\Delta t$ en h ($1~\mathrm{kWh} = 3{,}6 \times 10^{6}~\mathrm{J}$).
\item \textbf{Vérifier l'homogénéité} : $\Omega \cdot \mathrm{A}^2 = \mathrm{V \cdot A} = \mathrm{W}$.
\item \textbf{Conclure} avec une phrase physique (échauffement, coût, rendement...).
\end{enumerate}
\textbf{Réflexe} : pour un générateur, distinguer puissance totale $EI$, puissance utile $UI$ et puissance dissipée $rI^2$, avec le bilan $EI = UI + rI^2$.
\end{methode}

\begin{exoresolu}[Loi de Pouillet et facture d'électricité]
Un fer à repasser camerounais porte les indications « $\SI{220}{V}$ ; $\SI{1100}{W}$ ». On l'utilise sous la tension nominale pendant $1~\mathrm{h}~30~\mathrm{min}$.
\begin{enumerate}
\item Calculer la résistance $R$ du fer et l'intensité $I$ qui le traverse.
\item Calculer l'énergie électrique consommée en kWh puis en joules.
\item Le kWh est facturé $99~\mathrm{FCFA}$ par ENEO. Quel est le coût de cette séance de repassage ?
\end{enumerate}
\tcblower
\textbf{1. Résistance et intensité.} Le fer est un récepteur purement résistif : la loi de Pouillet s'applique. À partir de $\mathcal{P} = \dfrac{U^2}{R}$ :
\[ R = \frac{U^2}{\mathcal{P}} = \frac{220^2}{1\,100} = \frac{48\,400}{1\,100} = \SI{44}{\ohm}. \]
Intensité nominale :
\[ I = \frac{\mathcal{P}}{U} = \frac{1\,100}{220} = \SI{5,0}{A}. \]
Vérification par la loi d'Ohm : $I = \dfrac{U}{R} = \dfrac{220}{44} = \SI{5,0}{A}$. Cohérent.

\textbf{2. Énergie consommée.} Durée : $\Delta t = 1~\mathrm{h}~30~\mathrm{min} = \SI{1,5}{h} = \SI{5400}{s}$.
\[ W = \mathcal{P}\,\Delta t = \SI{1,1}{kW} \times \SI{1,5}{h} = \SI{1,65}{kWh}. \]
Conversion en joules :
\[ W = 1\,100 \times 5\,400 = 5{,}94 \times 10^{6}~\mathrm{J} = \SI{5,94}{MJ}. \]
On retrouve bien $1{,}65 \times 3{,}6 \times 10^{6} = 5{,}94 \times 10^{6}~\mathrm{J}$.

\textbf{3. Coût.} $C = 1{,}65 \times 99 = \SI{163,35}{FCFA} \approx \SI{163}{FCFA}$.

\textbf{Conclusion} : le fer a une résistance de $\SI{44}{\ohm}$, est traversé par $\SI{5,0}{A}$, consomme $\SI{1,65}{kWh}$ ($\SI{5,94}{MJ}$) en $1~\mathrm{h}~30$, soit environ $\SI{163}{FCFA}$.
\end{exoresolu}

\begin{methode}[Point de fonctionnement avec une diode]
\begin{enumerate}
\item \textbf{Modéliser la diode} selon l'énoncé : diode idéale ($U_D = 0$ si passante, interrupteur ouvert si bloquée) ou diode à seuil ($U_D = U_S \approx \SI{0,7}{V}$ pour le silicium si passante, $I = 0$ si $U_D < U_S$).
\item \textbf{Tester l'hypothèse « diode passante »} : supposer la diode conductrice, écrire la loi des mailles $E = rI + U_S + RI$, en tirer $I = \dfrac{E - U_S}{r + R}$.
\item \textbf{Vérifier la cohérence} : si $I > 0$, l'hypothèse est validée ; si le calcul donne $I < 0$, la diode est bloquée et $I = 0$.
\item \textbf{Identifier le point de fonctionnement} : $F(U_F = U_S\,;\,I_F)$ pour la diode à seuil passante ; $F(U_F = E\,;\,0)$ si bloquée (montage en inverse).
\item \textbf{Attention au sens de branchement} : une diode montée en inverse (anode côté borne négative) est bloquée : $I = 0$ dans tout le circuit.
\end{enumerate}
\end{methode}

\begin{exoresolu}[Circuit avec diode à seuil]
Un générateur ($E = \SI{6,0}{V}$, $r = \SI{5,0}{\ohm}$) alimente en série une diode au silicium (tension de seuil $U_S = \SI{0,7}{V}$, supposée idéale en conduction) et un résistor $R = \SI{100}{\ohm}$. La diode est montée dans le sens passant.
\begin{enumerate}
\item Déterminer l'intensité $I_F$ du courant dans le circuit.
\item En déduire les coordonnées du point de fonctionnement de la diode.
\item On inverse le sens de la diode. Que devient l'intensité ?
\end{enumerate}
\tcblower
\textbf{1. Intensité du courant.} Hypothèse : la diode est passante, donc $U_D = U_S = \SI{0,7}{V}$. La loi des mailles s'écrit :
\[ E = rI + U_S + RI. \]
On en tire :
\[ I_F = \frac{E - U_S}{r + R} = \frac{6{,}0 - 0{,}7}{5{,}0 + 100} = \frac{5{,}3}{105} \approx \SI{0,050}{A} = \SI{50}{mA}. \]
Vérification de l'hypothèse : $I_F > 0$, la diode conduit effectivement. L'hypothèse est validée.

\textbf{2. Point de fonctionnement de la diode.} Dans le modèle à seuil, la caractéristique de la diode passante est la demi-droite verticale $U_D = U_S$, $I > 0$. Le point de fonctionnement est donc :
\[ F\left(U_F = \SI{0,7}{V}\,;\,I_F = \SI{50}{mA}\right). \]
Vérification par la droite de charge : $U_D = E - (r+R)I = 6{,}0 - 105 \times \frac{5{,}3}{105} = \SI{0,7}{V}$ (exact).

\textbf{3. Diode inversée.} La diode est alors montée en inverse : elle est bloquée, aucun courant ne peut circuler :
\[ I = \SI{0}{A}. \]
La tension aux bornes de la diode vaut alors $U_D = E - (r+R)\times 0 = \SI{6,0}{V}$ (en polarisation inverse).

\textbf{Conclusion} : dans le sens passant, le circuit est parcouru par $I_F \approx \SI{50}{mA}$ et la diode fonctionne au point $F(\SI{0,7}{V}\,;\,\SI{50}{mA})$ ; dans le sens inverse, le courant est nul : la diode joue le rôle d'interrupteur ouvert.
\end{exoresolu}

\begin{methode}[Point de fonctionnement avec un transistor (montage amplificateur)]
\begin{enumerate}
\item \textbf{Identifier le régime} : en régime linéaire (amplification), le transistor impose $I_C = \beta\, I_B$ où $\beta$ est le gain en courant, et $U_{BE} \approx \SI{0,7}{V}$ (jonction base-émetteur passante).
\item \textbf{Calculer le courant de base} par la maille d'entrée : $I_B = \dfrac{E_B - U_{BE}}{R_B}$.
\item \textbf{En déduire le courant collecteur} : $I_C = \beta I_B$.
\item \textbf{Écrire la maille de sortie} (droite de charge) : $E_C = R_C I_C + U_{CE}$, d'où $U_{CE} = E_C - R_C I_C$.
\item \textbf{Vérifier le régime linéaire} : il faut $U_{CE} > U_{CE,\text{sat}} \approx \SI{0,2}{V}$. Si le calcul donne $U_{CE} \leq 0$, le transistor est \cle{saturé} : $U_{CE} \approx 0$ et $I_C = \dfrac{E_C}{R_C}$ (la relation $I_C = \beta I_B$ n'est alors plus valable).
\item \textbf{Conclure} : le point de fonctionnement est $Q(I_C\,;\,U_{CE})$, appelé aussi point de repos.
\end{enumerate}
\end{methode}

\begin{exoresolu}[Point de repos d'un transistor NPN]
Dans un montage amplificateur, un transistor NPN de gain $\beta = 100$ est alimenté par $E_C = \SI{12}{V}$. La maille de sortie comporte $R_C = \SI{1,0}{k\Omega}$ ; la maille d'entrée comporte $R_B = \SI{470}{k\Omega}$ reliée à $E_C$. On prend $U_{BE} = \SI{0,7}{V}$.
\begin{enumerate}
\item Calculer $I_B$, puis $I_C$ en régime linéaire.
\item Calculer $U_{CE}$ et conclure sur le régime de fonctionnement.
\item Donner les coordonnées du point de fonctionnement $Q$.
\end{enumerate}
\tcblower
\textbf{1. Courants de base et de collecteur.} Maille d'entrée : $E_C = R_B I_B + U_{BE}$, donc :
\[ I_B = \frac{E_C - U_{BE}}{R_B} = \frac{12 - 0{,}7}{470 \times 10^{3}} = \frac{11{,}3}{470 \times 10^{3}} \approx \SI{2,4e-5}{A} = \SI{24}{\micro A}. \]
En régime linéaire :
\[ I_C = \beta I_B = 100 \times \SI{24}{\micro A} = \SI{2,4}{mA}. \]

\textbf{2. Tension collecteur-émetteur.} Maille de sortie (droite de charge) : $E_C = R_C I_C + U_{CE}$, donc :
\[ U_{CE} = E_C - R_C I_C = 12 - 1{,}0 \times 10^{3} \times 2{,}4 \times 10^{-3} = 12 - 2{,}4 = \SI{9,6}{V}. \]
Comme $U_{CE} = \SI{9,6}{V} > \SI{0,2}{V}$, le transistor fonctionne bien en \cle{régime linéaire} : l'hypothèse $I_C = \beta I_B$ est validée.

\textbf{3. Point de fonctionnement.} Le point de repos du transistor est :
\[ Q\left(U_{CE} = \SI{9,6}{V}\,;\,I_C = \SI{2,4}{mA}\right). \]
Graphiquement, $Q$ est l'intersection de la droite de charge $U_{CE} = E_C - R_C I_C$ (reliant $(I_C = 0\,;\,\SI{12}{V})$ à $(I_C = \SI{12}{mA}\,;\,0)$) avec la caractéristique de sortie du transistor correspondant à $I_B = \SI{24}{\micro A}$.

\textbf{Conclusion} : le transistor est polarisé en régime linéaire au point $Q(\SI{9,6}{V}\,;\,\SI{2,4}{mA})$, avec $I_B = \SI{24}{\micro A}$ : un faible courant de base commande un courant collecteur cent fois plus grand, c'est le principe de l'amplification.
\end{exoresolu}

\begin{methode}[Exploiter un bilan de puissance dans un circuit complet]
\begin{enumerate}
\item \textbf{Déterminer le point de fonctionnement} d'abord (calcul ou graphique) : aucun bilan de puissance n'est possible sans connaître $I$ et $U$.
\item \textbf{Écrire les trois puissances} du générateur : totale $\mathcal{P}_{\text{tot}} = EI$, utile $\mathcal{P}_u = UI$, dissipée par effet Joule $\mathcal{P}_J = rI^2$.
\item \textbf{Vérifier le bilan} : $\mathcal{P}_{\text{tot}} = \mathcal{P}_u + \mathcal{P}_J$. Cette vérification est un excellent réflexe de contrôle.
\item \textbf{Calculer le rendement} du générateur : $\eta = \dfrac{\mathcal{P}_u}{\mathcal{P}_{\text{tot}}} = \dfrac{U}{E}$ (grandeur sans unité, souvent exprimée en \%).
\item \textbf{Conclure} en commentant le rendement (proche de 1 si $r$ est faible devant $R$).
\end{enumerate}
\end{methode}

\begin{exoresolu}[Bilan de puissance d'une batterie de mototaxi]
La batterie d'une moto ($E = \SI{12}{V}$, $r = \SI{0,10}{\ohm}$) alimente le démarreur, équivalent à un résistor $R = \SI{0,50}{\ohm}$.
\begin{enumerate}
\item Déterminer l'intensité $I$ du courant de démarrage.
\item Calculer la puissance totale fournie par la batterie, la puissance utile reçue par le démarreur et la puissance dissipée dans la batterie.
\item Calculer le rendement du transfert. Commenter.
\end{enumerate}
\tcblower
\textbf{1. Intensité.} Le circuit est équivalent à un générateur $(E, r)$ débitant sur $R$. La loi des mailles donne $E = (r + R)I$ :
\[ I = \frac{E}{r + R} = \frac{12}{0{,}10 + 0{,}50} = \frac{12}{0{,}60} = \SI{20}{A}. \]

\textbf{2. Bilan de puissance.}
\begin{align*}
\mathcal{P}_{\text{tot}} &= EI = 12 \times 20 = \SI{240}{W}, \\
\mathcal{P}_u &= UI = RI^2 = 0{,}50 \times 20^2 = 0{,}50 \times 400 = \SI{200}{W}, \\
\mathcal{P}_J &= rI^2 = 0{,}10 \times 400 = \SI{40}{W}.
\end{align*}
Vérification du bilan : $\mathcal{P}_u + \mathcal{P}_J = 200 + 40 = \SI{240}{W} = \mathcal{P}_{\text{tot}}$. Le bilan est cohérent.

\textbf{3. Rendement.}
\[ \eta = \frac{\mathcal{P}_u}{\mathcal{P}_{\text{tot}}} = \frac{200}{240} \approx 0{,}83, \quad \text{soit } \eta \approx 83~\%. \]
On vérifie aussi : $\eta = \dfrac{U}{E} = \dfrac{RI}{E} = \dfrac{10}{12} \approx 0{,}83$.

\textbf{Conclusion} : au démarrage, la batterie débite $\SI{20}{A}$ ; $\SI{200}{W}$ sont utiles au démarreur, $\SI{40}{W}$ chauffent la batterie, et le rendement du transfert est d'environ $83~\%$. C'est pourquoi une batterie fatiguée ($r$ augmenté) démarre mal : la puissance utile chute.
\end{exoresolu}

\begin{attention}[Les 5 erreurs qui coûtent des points au Probatoire]
\begin{enumerate}
\item \textbf{Oublier la résistance interne du générateur.} Écrire $U = E$ au lieu de $U = E - rI$ fausse toute la droite de charge et le point de fonctionnement. La tension aux bornes d'un générateur qui débite est \emph{toujours} inférieure à sa f.é.m.
\item \textbf{Négliger les conversions d'unités.} Les mA, $\mu$A, k$\Omega$ doivent être convertis en A et $\Omega$ avant tout calcul : $I = \dfrac{6 - 0{,}7}{105}$ donne $\SI{0,050}{A}$, pas $50~\mathrm{A}$ ! Une erreur de facteur $10^3$ est immédiatement sanctionnée.
\item \textbf{Appliquer la loi de Pouillet à un dipôle non résistif.} $\mathcal{P} = RI^2$ n'est valable que pour un résistor. Pour un moteur ou un électrolyseur, la puissance reçue est $UI$ et seule une partie est dissipée par effet Joule ($r'I^2$).
\item \textbf{Ne pas vérifier l'hypothèse de conduction (diode) ou le régime linéaire (transistor).} Après avoir supposé la diode passante, il faut vérifier $I > 0$ ; après avoir utilisé $I_C = \beta I_B$, il faut vérifier $U_{CE} > \SI{0,2}{V}$. Une hypothèse non validée rend le raisonnement incomplet.
\item \textbf{Lire le point de fonctionnement sur une seule courbe.} Le point de fonctionnement est l'\emph{intersection} de la caractéristique du dipôle et de la droite de charge : il faut tracer les deux sur le même graphique, avec des axes correctement gradués et légendés, et donner les deux coordonnées $U_F$ et $I_F$ avec leurs unités.
\end{enumerate}
\end{attention}

\newpage

\coursec{Exercices}

\courssub{Je vérifie mes connaissances}

\begin{exercice}[QCM : le point de fonctionnement]
Pour chaque affirmation, choisir la (ou les) bonne(s) réponse(s) et justifier brièvement.
\begin{enumerate}
\item Le point de fonctionnement d'un circuit constitué d'un générateur et d'un dipôle est :
\begin{enumerate}
\item le point d'intersection des caractéristiques $I(U)$ du générateur et du dipôle ;
\item le point de la caractéristique du dipôle où la tension est maximale ;
\item le couple $(U_0\,;\,I_0)$ vérifiant simultanément les lois des deux dipôles.
\end{enumerate}
\item La droite de charge d'un générateur $(E, r)$ débitant sur un dipôle a pour équation :
\begin{enumerate}
\item $I = \dfrac{E - U}{r}$ ;
\item $I = \dfrac{U - E}{r}$ ;
\item $I = \dfrac{E}{r} - \dfrac{U}{r}$.
\end{enumerate}
\item La loi de Pouillet s'écrit, pour un circuit série simple comportant un générateur $(E, r)$ et un résistor $R$ :
\begin{enumerate}
\item $I = \dfrac{E}{R}$ ;
\item $I = \dfrac{E}{R + r}$ ;
\item $I = \dfrac{E - U}{R + r}$.
\end{enumerate}
\item La puissance électrique reçue par un dipôle récepteur traversé par un courant $I$ sous une tension $U$ (convention récepteur) vaut :
\begin{enumerate}
\item $P = \dfrac{U}{I}$ ;
\item $P = U \times I$ ;
\item $P = R\,I$.
\end{enumerate}
\end{enumerate}
\end{exercice}

\begin{exercice}[Vrai ou faux, avec justification]
Répondre par vrai ou faux en justifiant chaque réponse par une phrase ou un calcul.
\begin{enumerate}
\item Une diode idéale polarisée en direct se comporte comme un interrupteur fermé.
\item Le point de fonctionnement d'un circuit dépend uniquement du dipôle récepteur.
\item Dans la convention récepteur, si le produit $U \times I$ est positif, le dipôle consomme de l'énergie électrique.
\item Un transistor NPN est saturé lorsque $V_{CE} = V_{CE,\text{sat}} \approx 0{,}2~\text{V}$ et que $I_C < \beta\, I_B$.
\item La résistance interne d'un générateur idéal de tension est infinie.
\end{enumerate}
\end{exercice}

\begin{exercice}[Définitions à compléter]
Recopier et compléter les phrases suivantes avec les mots ou expressions de la liste : \emph{point de fonctionnement ; caractéristique ; droite de charge ; saturation ; blocage ; loi de Pouillet ; convention récepteur}.
\begin{enumerate}
\item La courbe $I = f(U)$ d'un dipôle s'appelle sa \ldots\ courant--tension.
\item Le \ldots\ est le couple $(U_0\,;\,I_0)$ qui décrit l'état électrique réel du circuit.
\item Graphiquement, on détermine ce point en traçant la \ldots\ du générateur sur le même graphique que la caractéristique du dipôle.
\item Un transistor fonctionne en \ldots\ lorsque $I_B = 0$ : aucun courant ne traverse le collecteur.
\item La \ldots\ donne l'expression de l'intensité du courant dans un circuit série à partir des caractéristiques du générateur et des récepteurs.
\item Pour calculer la puissance reçue par un dipôle, on flèche $U$ et $I$ en sens opposés : c'est la \ldots.
\end{enumerate}
\end{exercice}

\begin{exercice}[Calculs directs]
\begin{enumerate}
\item Un générateur de f.é.m. $E = 12~\text{V}$ et de résistance interne $r = 2~\Omega$ alimente un résistor $R = 10~\Omega$. En appliquant la loi de Pouillet, calculer l'intensité $I$ du courant, la tension $U$ aux bornes du résistor et la puissance $P$ qu'il reçoit.
\item Un dipôle est traversé par un courant $I = 0{,}35~\text{A}$ sous une tension $U = 4{,}5~\text{V}$ (convention récepteur). Calculer la puissance reçue et l'énergie consommée en $10$ minutes, exprimée en joules puis en wattheures.
\item Un transistor NPN de gain $\beta = 100$ est parcouru par un courant de base $I_B = 0{,}20~\text{mA}$ en régime linéaire. Calculer le courant de collecteur $I_C$.
\end{enumerate}
\end{exercice}

\courssub{Je m'entraîne}

\begin{exercice}[Droite de charge d'un générateur]
Un générateur de f.é.m. $E = 9~\text{V}$ et de résistance interne $r = 15~\Omega$ alimente un dipôle inconnu.
\begin{enumerate}
\item Établir l'équation de la droite de charge $I = f(U)$ du générateur.
\item Calculer les coordonnées des points d'intersection de cette droite avec les axes (courant de court-circuit et tension à vide).
\item Tracer cette droite dans un repère $(U, I)$ bien choisi : échelles, unités, flèches.
\item Le dipôle est un résistor $R = 30~\Omega$. Déterminer graphiquement puis par le calcul les coordonnées du point de fonctionnement.
\end{enumerate}
\end{exercice}

\begin{exercice}[Diode et résistance de protection]
Une diode électroluminescente (DEL) rouge, assimilée à une diode de tension de seuil $U_S = 1{,}8~\text{V}$ et de résistance dynamique négligeable, est branchée en série avec un résistor $R = 220~\Omega$ aux bornes d'une pile de f.é.m. $E = 4{,}5~\text{V}$ et de résistance interne négligeable.
\begin{enumerate}
\item Faire le schéma du circuit en fléchant les tensions et le courant (convention récepteur pour les dipôles passifs).
\item Montrer que la diode est polarisée en direct.
\item Calculer l'intensité $I_0$ du courant et la tension $U_0$ aux bornes de la diode au point de fonctionnement.
\item Calculer la puissance reçue par la DEL et celle dissipée par le résistor. Vérifier la conservation de la puissance.
\item Que devient le courant si on inverse le sens de la diode ? Justifier.
\end{enumerate}
\end{exercice}

\begin{exercice}[Transistor en commutation]
Un transistor NPN ($\beta = 120$, $V_{BE} = 0{,}7~\text{V}$, $V_{CE,\text{sat}} = 0{,}2~\text{V}$) commande un relais de résistance $R_C = 80~\Omega$ alimenté sous $E = 12~\text{V}$. La base est alimentée à travers un résistor $R_B = 4{,}7~\text{k}\Omega$ par une tension $E_B = 5~\text{V}$.
\begin{enumerate}
\item Calculer le courant de base $I_B$.
\item En supposant le régime linéaire, calculer $I_C = \beta I_B$, puis la tension $V_{CE}$ correspondante. Conclure sur le régime réel du transistor.
\item Déterminer le courant de saturation $I_{C,\text{sat}}$ et vérifier que le transistor est bien saturé.
\item Calculer la puissance dissipée par le transistor en saturation et la puissance reçue par le relais.
\end{enumerate}
\end{exercice}

\begin{exercice}[Bilan énergétique d'un circuit]
Un moteur électrique de f.c.é.m. $E' = 6~\text{V}$ et de résistance interne $r' = 2~\Omega$ est monté en série avec un résistor $R = 8~\Omega$ aux bornes d'un générateur $(E = 18~\text{V}\,;\, r = 2~\Omega)$.
\begin{enumerate}
\item Écrire la loi de Pouillet généralisée et calculer l'intensité $I$ du courant.
\item Calculer la puissance totale fournie par le générateur, la puissance utile du moteur et la puissance dissipée par effet Joule dans l'ensemble du circuit.
\item Vérifier le bilan de puissance et calculer le rendement du circuit (puissance utile / puissance fournie).
\end{enumerate}
\end{exercice}

\courssub{Je me prépare au Probatoire}

\begin{exercice}[Point de fonctionnement d'une diode — type Probatoire]
Au laboratoire du lycée, un groupe d'élèves relève la caractéristique courant--tension d'une diode au silicium et obtient le tableau suivant :
\begin{center}
\begin{tabular}{|l|c|c|c|c|c|c|c|c|}
\hline
\rowcolor{popblueL} $U$ (V) & 0 & 0{,}40 & 0{,}55 & 0{,}60 & 0{,}65 & 0{,}70 & 0{,}75 & 0{,}80 \\
\hline
$I$ (mA) & 0 & 0{,}2 & 1{,}0 & 2{,}5 & 5{,}0 & 10 & 20 & 40 \\
\hline
\end{tabular}
\end{center}
La diode est ensuite branchée en série avec un résistor $R = 220~\Omega$ aux bornes d'un générateur de f.é.m. $E = 5{,}0~\text{V}$ et de résistance interne négligeable.
\begin{enumerate}
\item Tracer sur papier millimétré la caractéristique $I = f(U)$ de la diode. Échelles conseillées : $1~\text{cm}$ pour $0{,}10~\text{V}$ et $1~\text{cm}$ pour $5~\text{mA}$.
\item Établir, en justifiant par la loi des mailles, l'équation de la droite de charge du circuit : $I = \dfrac{E - U}{R}$.
\item Tracer la droite de charge sur le même graphique et déterminer graphiquement les coordonnées $(U_0\,;\,I_0)$ du point de fonctionnement.
\item Calculer la puissance $P_D$ reçue par la diode et la puissance $P_R$ dissipée dans le résistor au point de fonctionnement.
\item Le générateur réel possède en fait une résistance interne $r = 10~\Omega$. Écrire la nouvelle équation de la droite de charge et dire, sans nouveau tracé, dans quel sens se déplace le point de fonctionnement.
\item Un élève affirme avoir trouvé $I_0 = 23~\text{mA}$ et $U_0 = 0{,}7~\text{V}$. Montrer, par un calcul, que sa réponse est incompatible avec la droite de charge et identifier l'erreur probable.
\end{enumerate}
\end{exercice}

\begin{exercice}[Commande d'un ventilateur par transistor — type Probatoire]
Dans une salle informatique de Yaoundé, un transistor NPN commande un ventilateur de refroidissement modélisé par une résistance $R_C = 100~\Omega$, alimenté sous $E = 12~\text{V}$. Données du transistor : $\beta = 150$ ; $V_{BE} = 0{,}7~\text{V}$ ; $V_{CE,\text{sat}} = 0{,}2~\text{V}$. La base est commandée par une tension $E_B = 5~\text{V}$ à travers un résistor $R_B$.
\begin{enumerate}
\item Faire le schéma du montage (émetteur commun, charge au collecteur) en fléchant $I_B$, $I_C$, $V_{BE}$ et $V_{CE}$.
\item Le transistor est saturé. Montrer que le courant de collecteur en saturation vaut $I_{C,\text{sat}} = \dfrac{E - V_{CE,\text{sat}}}{R_C}$ et calculer sa valeur.
\item En déduire le courant de base minimal $I_{B,\min}$ nécessaire pour atteindre la saturation.
\item Pour garantir la saturation, on choisit un coefficient de sursaturation $k = 2$ : $I_B = k\, I_{B,\min}$. Calculer la valeur maximale du résistor $R_B$ à utiliser.
\item Calculer la puissance dissipée par le transistor en saturation ($P_T = V_{BE} I_B + V_{CE,\text{sat}} I_C$) et la puissance reçue par le ventilateur.
\item Calculer le rendement du montage en commutation et commenter l'intérêt de faire fonctionner le transistor en saturation plutôt qu'en régime linéaire.
\end{enumerate}
\end{exercice}

\begin{exercice}[Éclairage d'une maison en zone rurale — problème]
Dans un village non raccordé au réseau ENEO, une installation solaire comprend une batterie de f.é.m. $E = 12{,}6~\text{V}$ et de résistance interne $r = 0{,}10~\Omega$. Elle alimente, par un câble de résistance totale $R_f = 0{,}40~\Omega$, une lampe à DEL dont la caractéristique est assimilable à une droite : $U = U_S + R_d\, I$ avec $U_S = 2{,}8~\text{V}$ et $R_d = 6{,}0~\Omega$.
\begin{enumerate}
\item Faire le schéma du circuit équivalent en série et écrire la loi des mailles.
\item Établir l'expression littérale de l'intensité $I_0$ du courant (loi de Pouillet) puis calculer sa valeur.
\item Déterminer les coordonnées du point de fonctionnement $(U_0\,;\,I_0)$ de la lampe, où $U_0$ est la tension à ses bornes.
\item Calculer la puissance reçue par la lampe, la puissance perdue dans le câble et dans la batterie, puis le rendement global de l'installation.
\item La lampe fonctionne en moyenne $5~\text{h}$ par nuit. Calculer l'énergie qu'elle consomme en une semaine, en wattheures.
\item Le propriétaire souhaite ajouter une deuxième lampe identique en dérivation. Expliquer qualitativement comment se déplace le point de fonctionnement de chaque lampe et pourquoi l'éclairement risque de diminuer.
\end{enumerate}
\end{exercice}

\newpage

\coursec{Corrigés des exercices}

\begin{corrige}[Exercice 1 — QCM : le point de fonctionnement]
\textbf{1. Réponses \textbf{a} et \textbf{c}.} Le point de fonctionnement est le couple $(U_0\,;\,I_0)$ qui vérifie simultanément la loi du générateur et la loi du dipôle ; graphiquement, c'est donc l'intersection des deux caractéristiques $I(U)$. La réponse b est fausse : la tension au point de fonctionnement n'a aucune raison d'être maximale.

\textbf{2. Réponses \textbf{a} et \textbf{c}.} La loi des mailles donne $U = E - rI$, d'où $I = \dfrac{E-U}{r} = \dfrac{E}{r} - \dfrac{U}{r}$ : les deux écritures sont équivalentes. La réponse b donnerait une pente positive, absurde pour un générateur.

\textbf{3. Réponse \textbf{b}.} La loi de Pouillet s'écrit $I = \dfrac{E}{R+r}$ : la f.é.m. divisée par la somme de \emph{toutes} les résistances de la maille, y compris la résistance interne. La réponse a oublie $r$ ; la réponse c mélange tension et f.é.m.

\textbf{4. Réponse \textbf{b}.} En convention récepteur, $P = U \times I$. La réponse a est homogène à une résistance, la réponse c à un courant : l'analyse dimensionnelle les élimine immédiatement.
\end{corrige}

\begin{corrige}[Exercice 2 — Vrai ou faux, avec justification]
\textbf{1. \textbf{Vrai}.} Une diode idéale polarisée en direct présente une tension nulle à ses bornes quel que soit le courant : elle se comporte comme un interrupteur fermé (un fil).

\textbf{2. \textbf{Faux}.} Le point de fonctionnement dépend du générateur \emph{et} du dipôle : c'est l'intersection de la droite de charge du générateur $(E, r)$ avec la caractéristique du récepteur. Modifier $E$ ou $r$ déplace le point de fonctionnement même si le dipôle est inchangé.

\textbf{3. \textbf{Vrai}.} En convention récepteur, $P = UI > 0$ signifie que le dipôle reçoit effectivement de la puissance électrique du reste du circuit : il consomme de l'énergie.

\textbf{4. \textbf{Vrai}.} En saturation, $V_{CE} = V_{CE,\text{sat}} \approx \SI{0,2}{V}$ et le courant de collecteur est imposé par le circuit extérieur : $I_C = I_{C,\text{sat}} < \beta I_B$. La relation $I_C = \beta I_B$ n'est plus vérifiée.

\textbf{5. \textbf{Faux}.} Un générateur idéal de tension a une résistance interne \emph{nulle} : la tension à ses bornes vaut $E$ quel que soit le courant débité. C'est le \emph{générateur idéal de courant} qui possède une résistance interne infinie.
\end{corrige}

\begin{corrige}[Exercice 3 — Définitions à compléter]
\begin{enumerate}
\item La courbe $I = f(U)$ d'un dipôle s'appelle sa \cle{caractéristique} courant--tension.
\item Le \cle{point de fonctionnement} est le couple $(U_0\,;\,I_0)$ qui décrit l'état électrique réel du circuit.
\item Graphiquement, on détermine ce point en traçant la \cle{droite de charge} du générateur sur le même graphique que la caractéristique du dipôle.
\item Un transistor fonctionne en \cle{blocage} lorsque $I_B = 0$ : aucun courant ne traverse le collecteur.
\item La \cle{loi de Pouillet} donne l'expression de l'intensité du courant dans un circuit série à partir des caractéristiques du générateur et des récepteurs.
\item Pour calculer la puissance reçue par un dipôle, on flèche $U$ et $I$ en sens opposés : c'est la \cle{convention récepteur}.
\end{enumerate}
\end{corrige}

\begin{corrige}[Exercice 4 — Calculs directs]
\textbf{1.} Loi de Pouillet : $I = \dfrac{E}{R+r} = \dfrac{12}{10+2} = \SI{1,0}{A}$.

Tension aux bornes du résistor : $U = RI = 10 \times 1,0 = \SI{10}{V}$ (on vérifie : $U = E - rI = 12 - 2 = \SI{10}{V}$).

Puissance reçue : $P = UI = RI^2 = 10 \times 1,0^2 = \SI{10}{W}$.

\textbf{2.} Puissance reçue : $P = UI = 4,5 \times 0,35 = \SI{1,575}{W} \approx \SI{1,6}{W}$.

Énergie sur $\Delta t = 10~\text{min} = \SI{600}{s}$ :
\[ W = P\,\Delta t = 1,575 \times 600 = \SI{945}{J}. \]
En wattheures : $W = \SI{1,575}{W} \times \dfrac{10}{60}~\text{h} = \SI{0,2625}{Wh} \approx \SI{0,26}{Wh}$. (Vérification : $945/3600 = \SI{0,2625}{Wh}$.)

\textbf{3.} En régime linéaire : $I_C = \beta I_B = 100 \times \SI{0,20}{mA} = \SI{20}{mA}$.
\end{corrige}

\begin{corrige}[Exercice 5 — Droite de charge d'un générateur]
\textbf{1.} La loi des mailles s'écrit $E = rI + U$, d'où l'équation de la droite de charge :
\[ I = \frac{E - U}{r} = \frac{9 - U}{15} \qquad (U \text{ en V},\, I \text{ en A}). \]

\textbf{2.} Intersection avec l'axe des courants ($U = 0$, court-circuit) : $I_{CC} = \dfrac{E}{r} = \dfrac{9}{15} = \SI{0,60}{A}$.

Intersection avec l'axe des tensions ($I = 0$, circuit ouvert) : $U = E = \SI{9,0}{V}$ (tension à vide).

\textbf{3.} On place les deux points $(\SI{0}{V}\,;\,\SI{0,60}{A})$ et $(\SI{9}{V}\,;\,\SI{0}{A})$ et on les relie par une droite. Échelles conseillées : $1~\text{cm}$ pour $1~\text{V}$ en abscisse, $1~\text{cm}$ pour $50~\text{mA}$ en ordonnée. La droite est décroissante, de pente $-\dfrac{1}{r} = -\SI{0,067}{A/V}$.

\textbf{4.} \emph{Par le calcul.} Le résistor impose $U = RI$. En égalisant avec la droite de charge :
\[ RI = E - rI \;\Longrightarrow\; I = \frac{E}{R+r} = \frac{9}{30+15} = \frac{9}{45} = \SI{0,20}{A}, \qquad U = RI = 30 \times 0,20 = \SI{6,0}{V}. \]
\emph{Graphiquement}, on trace la droite $I = U/R$ (passant par l'origine et par $(\SI{6}{V}\,;\,\SI{0,20}{A})$) ; son intersection avec la droite de charge redonne le point de fonctionnement $F(\SI{6,0}{V}\,;\,\SI{0,20}{A})$. Les deux méthodes concordent.
\end{corrige}

\begin{corrige}[Exercice 6 — Diode et résistance de protection]
\textbf{1.} Schéma : pile $(E)$ en série avec le résistor $R$ puis la DEL (anode vers le pôle $+$). Le courant $I$ sort du pôle $+$ de la pile ; les flèches $U_R$ et $U_D$ sont orientées en sens opposé de $I$ (convention récepteur), la flèche $E$ dans le même sens que $I$ (convention générateur).

\textbf{2.} L'anode de la diode est reliée, à travers $R$, au pôle positif de la pile : la diode est soumise à une tension directe. Comme $E = \SI{4,5}{V} > U_S = \SI{1,8}{V}$, le générateur peut imposer une tension supérieure au seuil : la diode conduit, elle est bien polarisée en direct.

\textbf{3.} La diode passante est modélisée par $U_D = U_S = \SI{1,8}{V}$ (résistance dynamique négligeable). La loi des mailles donne :
\[ E = RI_0 + U_S \;\Longrightarrow\; I_0 = \frac{E - U_S}{R} = \frac{4,5 - 1,8}{220} = \frac{2,7}{220} = \SI{1,23e-2}{A} \approx \SI{12,3}{mA}. \]
Point de fonctionnement de la diode : $U_0 = \SI{1,8}{V}$, $I_0 \approx \SI{12,3}{mA}$.

\textbf{4.} Puissance reçue par la DEL : $P_D = U_S I_0 = 1,8 \times 0,0123 = \SI{2,2e-2}{W} \approx \SI{22}{mW}$.

Puissance dissipée par le résistor : $P_R = RI_0^2 = 220 \times (0,0123)^2 = \SI{3,3e-2}{W} \approx \SI{33}{mW}$.

Vérification de la conservation : puissance fournie par la pile $P_E = E I_0 = 4,5 \times 0,0123 = \SI{55}{mW}$, et $P_D + P_R = 22 + 33 = \SI{55}{mW}$. Le bilan est vérifié.

\textbf{5.} Si on inverse la diode, elle est polarisée en inverse : elle est bloquée et se comporte comme un interrupteur ouvert. Le courant devient nul ($I = 0$), la DEL ne s'allume pas ; toute la tension du générateur se retrouve aux bornes de la diode ($U_D = -\SI{4,5}{V}$ en convention récepteur).
\end{corrige}

\begin{corrige}[Exercice 7 — Transistor en commutation]
\textbf{1.} Maille d'entrée (base--émetteur) : $E_B = R_B I_B + V_{BE}$, d'où
\[ I_B = \frac{E_B - V_{BE}}{R_B} = \frac{5 - 0,7}{4,7\times 10^{3}} = \frac{4,3}{4700} = \SI{9,15e-4}{A} \approx \SI{0,92}{mA}. \]

\textbf{2.} Hypothèse du régime linéaire : $I_C = \beta I_B = 120 \times \SI{0,915}{mA} = \SI{109,8}{mA} \approx \SI{0,11}{A}$.

Tension collecteur--émetteur correspondante (maille de sortie) :
\[ V_{CE} = E - R_C I_C = 12 - 80 \times 0,1098 = 12 - 8,78 = \SI{3,2}{V}. \]
Comme $V_{CE} = \SI{3,2}{V} > V_{CE,\text{sat}} = \SI{0,2}{V}$, l'hypothèse est cohérente : \textbf{le transistor fonctionne en régime linéaire} (amplification), pas en saturation.

\textbf{3.} Le courant de saturation est imposé par le circuit extérieur :
\[ I_{C,\text{sat}} = \frac{E - V_{CE,\text{sat}}}{R_C} = \frac{12 - 0,2}{80} = \frac{11,8}{80} = \SI{0,1475}{A} \approx \SI{148}{mA}. \]
Condition de saturation : $\beta I_B \geq I_{C,\text{sat}}$. Or $\beta I_B = \SI{109,8}{mA} < \SI{147,5}{mA}$ : la condition n'est \textbf{pas} remplie, ce qui confirme le régime linéaire. Pour saturer le transistor, il faudrait $I_B \geq \dfrac{I_{C,\text{sat}}}{\beta} = \dfrac{147,5}{120} = \SI{1,23}{mA}$, donc diminuer $R_B$ en dessous de $\dfrac{4,3}{1,23\times 10^{-3}} \approx \SI{3,5}{k\Omega}$.

\textbf{4.} Le transistor étant en régime linéaire (et non saturé), on calcule avec les valeurs réelles :
\[ P_T = V_{BE} I_B + V_{CE} I_C = 0,7 \times 0,915\times 10^{-3} + 3,2 \times 0,1098 = \SI{6,4e-4}{W} + \SI{0,353}{W} \approx \SI{0,35}{W}. \]
Puissance reçue par le relais : $P_R = R_C I_C^2 = 80 \times (0,1098)^2 = \SI{0,96}{W}$.

\emph{Remarque :} cet exercice illustre qu'en régime linéaire le transistor dissipe une puissance importante ($\SI{0,35}{W}$) ; c'est précisément pour éviter cela qu'en commutation on choisit $R_B$ assez faible pour garantir la saturation.
\end{corrige}

\begin{corrige}[Exercice 8 — Bilan énergétique d'un circuit]
\textbf{1.} Loi de Pouillet généralisée (les f.é.m. et f.c.é.m. se soustraient, toutes les résistances s'ajoutent) :
\[ I = \frac{E - E'}{r + r' + R} = \frac{18 - 6}{2 + 2 + 8} = \frac{12}{12} = \SI{1,0}{A}. \]

\textbf{2.} Puissance totale fournie par le générateur : $P_g = E I = 18 \times 1,0 = \SI{18}{W}$.

Puissance utile (mécanique) du moteur : $P_u = E' I = 6 \times 1,0 = \SI{6,0}{W}$.

Puissance dissipée par effet Joule dans l'ensemble du circuit :
\[ P_J = (r + r' + R) I^2 = 12 \times 1,0^2 = \SI{12}{W}. \]

\textbf{3.} Bilan : $P_u + P_J = 6 + 12 = \SI{18}{W} = P_g$. La conservation de la puissance est vérifiée.

Rendement du circuit : $\eta = \dfrac{P_u}{P_g} = \dfrac{6}{18} = 0,33$, soit environ $33~\%$. Les deux tiers de l'énergie fournie sont perdus par effet Joule, dont $R I^2 = \SI{8}{W}$ dans le seul résistor.
\end{corrige}

\begin{corrige}[Exercice 9 — Point de fonctionnement d'une diode (type Probatoire)]
\textbf{1.} On reporte les huit points $(U\,;\,I)$ du tableau et on les relie par une courbe lisse : la caractéristique est quasi nulle en dessous de $\SI{0,5}{V}$, puis croît très vite au-delà de $\SI{0,6}{V}$ (allure exponentielle typique d'une diode au silicium).

\textbf{2.} Le circuit est une maille unique : générateur $(E)$, résistor $R$, diode. En convention récepteur pour $R$ et la diode, la loi des mailles s'écrit :
\[ E = RI + U \;\Longrightarrow\; \boxed{I = \frac{E - U}{R}} \]
avec $E = \SI{5,0}{V}$ et $R = \SI{220}{\ohm}$. C'est l'équation de la droite de charge.

\textbf{3.} Points remarquables de la droite de charge : $(U = \SI{0}{V}\,;\, I = \frac{5}{220} = \SI{22,7}{mA})$ et $(U = \SI{5}{V}\,;\, I = 0)$. On la trace sur le même graphique. L'intersection avec la caractéristique se situe entre $U = \SI{0,70}{V}$ (droite : $\SI{19,5}{mA}$ ; courbe : $\SI{10}{mA}$) et $U = \SI{0,75}{V}$ (droite : $\SI{19,3}{mA}$ ; courbe : $\SI{20}{mA}$). Lecture graphique :
\[ U_0 \approx \SI{0,74}{V}\text{--}\SI{0,75}{V}, \qquad I_0 \approx \SI{19}{mA}\text{--}\SI{20}{mA}. \]
Toute valeur cohérente avec le tracé dans cet intervalle est acceptée.

\textbf{4.} Avec $U_0 = \SI{0,745}{V}$ et $I_0 = \SI{19,3}{mA}$ :
\[ P_D = U_0 I_0 = 0,745 \times 0,0193 = \SI{1,44e-2}{W} \approx \SI{14}{mW}, \]
\[ P_R = RI_0^2 = 220 \times (0,0193)^2 = \SI{8,2e-2}{W} \approx \SI{82}{mW}. \]
(Vérification : $P_E = EI_0 = 5 \times 0,0193 = \SI{96,5}{mW} \approx P_D + P_R = \SI{96,4}{mW}$.)

\textbf{5.} Avec $r = \SI{10}{\ohm}$, la loi des mailles devient $E = (R + r)I + U$, d'où la nouvelle droite de charge :
\[ I = \frac{E - U}{R + r} = \frac{5 - U}{230}. \]
La pente diminue légèrement en valeur absolue et le courant de court-circuit passe de $\SI{22,7}{mA}$ à $\SI{21,7}{mA}$ : la droite pivote légèrement vers le bas autour du point $(\SI{5}{V}\,;\,0)$. Le point de fonctionnement se déplace donc vers le \textbf{bas et vers la gauche} : $I_0$ diminue légèrement, $U_0$ aussi (très peu, car la caractéristique de la diode est presque verticale dans cette zone).

\textbf{6.} Testons la réponse de l'élève dans l'équation de la droite de charge : si $I_0 = \SI{23}{mA}$, alors
\[ U_0 = E - RI_0 = 5 - 220 \times 0,023 = 5 - 5,06 = -\SI{0,06}{V}, \]
ce qui est incompatible avec $U_0 = \SI{0,7}{V}$ (et physiquement absurde pour une diode en direct). L'erreur probable : l'élève a calculé $I = \dfrac{E}{R} = \dfrac{5}{220} = \SI{22,7}{mA} \approx \SI{23}{mA}$, c'est-à-dire le \emph{courant de court-circuit}, en négligeant la tension aux bornes de la diode — autrement dit il a pris le point d'intersection de la droite de charge avec l'axe des ordonnées au lieu du point d'intersection avec la caractéristique.

\textbf{Barème indicatif (sur 10) :} tracé soigné de la caractéristique (2 pts) ; loi des mailles justifiée et équation de la droite (2 pts) ; tracé de la droite et lecture de $(U_0\,;\,I_0)$ cohérente (2 pts) ; puissances avec unités (1,5 pt) ; nouvelle droite et sens de déplacement (1,5 pt) ; vérification par le calcul et identification de l'erreur (1 pt).

\textbf{Points de vigilance :} citer explicitement la loi des mailles avant d'écrire $I = \dfrac{E-U}{R}$ ; convertir les mA en A dans les calculs de puissance ; la lecture graphique doit être accompagnée des pointillés de construction sur la copie ; ne jamais confondre courant de court-circuit et courant au point de fonctionnement.
\end{corrige}

\begin{corrige}[Exercice 10 — Commande d'un ventilateur par transistor (type Probatoire)]
\textbf{1.} Schéma : émetteur à la masse ; collecteur relié au $+E$ à travers le ventilateur $R_C$ ; base alimentée par $E_B$ à travers $R_B$. On flèche $I_B$ entrant dans la base, $I_C$ entrant dans le collecteur, $V_{BE}$ de la base vers l'émetteur et $V_{CE}$ du collecteur vers l'émetteur (flèches tension opposées aux flèches courant, convention récepteur).

\textbf{2.} En saturation, $V_{CE} = V_{CE,\text{sat}}$. La loi des mailles sur le circuit collecteur s'écrit :
\[ E = R_C I_C + V_{CE,\text{sat}} \;\Longrightarrow\; I_{C,\text{sat}} = \frac{E - V_{CE,\text{sat}}}{R_C} = \frac{12 - 0,2}{100} = \frac{11,8}{100} = \SI{0,118}{A} = \SI{118}{mA}. \]

\textbf{3.} La saturation exige $\beta I_B \geq I_{C,\text{sat}}$, d'où le courant de base minimal :
\[ I_{B,\min} = \frac{I_{C,\text{sat}}}{\beta} = \frac{118}{150} = \SI{0,787}{mA} \approx \SI{0,79}{mA}. \]

\textbf{4.} Avec le coefficient de sursaturation $k = 2$ : $I_B = k\, I_{B,\min} = 2 \times 0,787 = \SI{1,57}{mA}$.

Maille d'entrée : $E_B = R_B I_B + V_{BE}$, d'où
\[ R_B \leq \frac{E_B - V_{BE}}{I_B} = \frac{5 - 0,7}{1,57\times 10^{-3}} = \frac{4,3}{1,57\times 10^{-3}} = \SI{2,73e3}{\ohm}. \]
On choisit donc $R_{B,\max} \approx \SI{2,7}{k\ohm}$ (valeur normalisée : $\SI{2,7}{k\ohm}$, ou $\SI{2,2}{k\ohm}$ pour plus de marge).

\textbf{5.} Puissance dissipée par le transistor en saturation :
\[ P_T = V_{BE} I_B + V_{CE,\text{sat}} I_C = 0,7 \times 1,57\times 10^{-3} + 0,2 \times 0,118 = \SI{1,1e-3}{W} + \SI{2,36e-2}{W} \approx \SI{2,5e-2}{W} = \SI{25}{mW}. \]
Puissance reçue par le ventilateur :
\[ P_V = R_C I_C^2 = 100 \times (0,118)^2 = \SI{1,39}{W} \quad (\text{ou } (E - V_{CE,\text{sat}}) I_C = 11,8 \times 0,118 = \SI{1,39}{W}). \]

\textbf{6.} Puissance fournie par l'alimentation principale : $P_E = E I_C = 12 \times 0,118 = \SI{1,42}{W}$. Rendement du montage :
\[ \eta = \frac{P_V}{P_E} = \frac{1,39}{1,42} = 0,983, \quad \text{soit environ } 98~\%. \]
Commentaire : en saturation, $V_{CE}$ est très faible ($\SI{0,2}{V}$), donc la puissance dissipée par le transistor est négligeable ($\SI{25}{mW}$ contre $\SI{1,39}{W}$ pour la charge) : presque toute l'énergie va au ventilateur. En régime linéaire, $V_{CE}$ serait de plusieurs volts et le transistor dissiperait une fraction importante de la puissance, d'où échauffement, besoin d'un radiateur et gaspillage d'énergie. C'est pourquoi, en commutation (tout-ou-rien), on fait toujours fonctionner le transistor en blocage ou en saturation.

\textbf{Barème indicatif (sur 10) :} schéma fléché correct (1,5 pt) ; démonstration de $I_{C,\text{sat}}$ par la loi des mailles (2 pts) ; $I_{B,\min}$ (1 pt) ; $R_{B,\max}$ avec inégalité justifiée (2 pts) ; puissances $P_T$ et $P_V$ (2 pts) ; rendement et commentaire pertinent (1,5 pt).

\textbf{Points de vigilance :} la condition de saturation $\beta I_B \geq I_{C,\text{sat}}$ doit être écrite explicitement ; pour $R_B$, il s'agit d'une valeur \emph{maximale} (un résistor plus petit donne un $I_B$ plus grand, donc une saturation plus sûre) ; ne pas utiliser $I_C = \beta I_B$ en saturation ; bien distinguer la puissance dissipée par le transistor de celle reçue par la charge.
\end{corrige}

\begin{corrige}[Exercice 11 — Éclairage d'une maison en zone rurale (problème)]
\textbf{1.} Circuit équivalent en série : batterie $(E, r)$, câble $R_f$, lampe modélisée par la f.c.é.m. $U_S$ en série avec $R_d$. En parcourant la maille dans le sens du courant, la loi des mailles s'écrit :
\[ E = rI + R_f I + U_S + R_d I. \]

\textbf{2.} On regroupe : $E - U_S = (r + R_f + R_d) I$, d'où la loi de Pouillet :
\[ I_0 = \frac{E - U_S}{r + R_f + R_d} = \frac{12,6 - 2,8}{0,10 + 0,40 + 6,0} = \frac{9,8}{6,5} = \SI{1,51}{A}. \]

\textbf{3.} Tension aux bornes de la lampe :
\[ U_0 = U_S + R_d I_0 = 2,8 + 6,0 \times 1,508 = 2,8 + 9,05 = \SI{11,85}{V} \approx \SI{11,8}{V}. \]
Point de fonctionnement de la lampe : $(U_0 \approx \SI{11,8}{V}\,;\, I_0 \approx \SI{1,51}{A})$. On remarque que la lampe reçoit moins que les $\SI{12,6}{V}$ de la batterie à cause des chutes de tension dans le câble et la résistance interne.

\textbf{4.} Puissance reçue par la lampe : $P_L = U_0 I_0 = 11,85 \times 1,508 = \SI{17,9}{W}$.

Puissance perdue dans le câble : $P_f = R_f I_0^2 = 0,40 \times (1,508)^2 = \SI{0,91}{W}$.

Puissance perdue dans la batterie : $P_r = r I_0^2 = 0,10 \times (1,508)^2 = \SI{0,23}{W}$.

Vérification : puissance fournie $P_E = E I_0 = 12,6 \times 1,508 = \SI{19,0}{W}$, et $P_L + P_f + P_r = 17,9 + 0,91 + 0,23 = \SI{19,0}{W}$. Bilan vérifié.

Rendement global : $\eta = \dfrac{P_L}{P_E} = \dfrac{17,9}{19,0} = 0,941$, soit environ $94~\%$.

\textbf{5.} Énergie hebdomadaire : $\Delta t = 5~\text{h} \times 7 = \SI{35}{h}$,
\[ W = P_L \, \Delta t = 17,9 \times 35 = \SI{6,3e2}{Wh} \approx \SI{0,63}{kWh}. \]

\textbf{6.} En ajoutant une deuxième lampe identique en dérivation, la résistance équivalente de l'ensemble des lampes diminue (deux branches en parallèle). Le courant total débité par la batterie augmente donc, et avec lui les chutes de tension $rI$ et $R_f I$ dans la batterie et le câble. La tension disponible aux bornes de l'ensemble des lampes diminue : le point de fonctionnement de chaque lampe glisse le long de sa caractéristique vers des valeurs plus faibles de $U$ et de $I$. Chaque lampe reçoit alors une puissance $P = UI$ plus faible : l'éclairement de chacune diminue. C'est le phénomène bien connu d'« affaissement » de la tension quand on branche trop d'appareils sur une petite installation solaire : d'où l'importance du dimensionnement des câbles (section suffisante pour réduire $R_f$) et de la batterie.
\end{corrige}

\newpage

\coursec{Fiche bilan}

\begin{popfigure}
\begin{tikzpicture}[mindmap, grow cyclic, every node/.style={concept, text=white, font=\small\bfseries},
  root concept/.append style={concept color=popdark, font=\large\bfseries},
  level 1/.append style={level distance=3.4cm, sibling angle=60},
  level 2/.append style={level distance=2.2cm, font=\footnotesize\bfseries}]
\node[root concept]{Point de fonctionnement et loi de Pouillet}
  child[concept color=popblue]{ node{Composants de base}
    child{ node{R\'esistor : $U=RI$} }
    child{ node{Diode, transistor} }
  }
  child[concept color=poporange]{ node{Caract\'eristique $I(U)$}
    child{ node{Trac\'e point par point} }
    child{ node{Dip\^oles passifs / actifs} }
  }
  child[concept color=popgreen]{ node{Point de fonctionnement}
    child{ node{Intersection des courbes} }
    child{ node{$Q(U_0\,;\,I_0)$} }
  }
  child[concept color=poppurple]{ node{Loi de Pouillet}
    child{ node{$I=\dfrac{E}{R+r}$} }
    child{ node{Circuit s\'erie complet} }
  }
  child[concept color=poppink]{ node{Puissance et \'energie}
    child{ node{$P=UI$ ; $W=UIt$} }
    child{ node{$P_{\text{utile}}=P_{\text{tot}}-rI^2$} }
  }
  child[concept color=popgold]{ node{M\'ethodes}
    child{ node{M\'ethode graphique} }
    child{ node{M\'ethode alg\'ebrique} }
  };
\end{tikzpicture}
\legende{Carte mentale de la le\c{c}on : point de fonctionnement d'un circuit et loi de Pouillet.}
\end{popfigure}

\begin{bilan}
\textbf{D\'efinitions cl\'es}
\begin{itemize}
  \item \cle{Caract\'eristique} courant--tension d'un dip\^ole : courbe $I=f(U)$ (ou $U=g(I)$) donnant l'intensit\'e qui traverse le dip\^ole en fonction de la tension \`a ses bornes.
  \item \cle{Point de fonctionnement} d'un circuit : point $Q(U_0\,;\,I_0)$, intersection de la caract\'eristique du dip\^ole r\'ecepteur et de la droite d'attaque du g\'en\'erateur ; il donne les valeurs r\'eelles $U_0$ et $I_0$ dans le circuit.
  \item \cle{Loi de Pouillet} : dans un circuit s\'erie ferm\'e comportant un g\'en\'erateur $(E,r)$ et des r\'esistors, l'intensit\'e est le quotient de la f.\'e.\'m. totale par la r\'esistance totale.
\end{itemize}

\textbf{Formules \`a conna\^itre par c\oe ur}
\begin{center}
\begin{tabularx}{\linewidth}{|l|X|l|}\hline
\rowcolor{popblueL} \textbf{Grandeur / loi} & \textbf{Formule} & \textbf{Unit\'es SI} \\ \hline
Loi d'Ohm (r\'esistor) & $U=RI$ & $U$ en V, $R$ en $\Omega$, $I$ en A \\ \hline
G\'en\'erateur $(E,r)$ & $U_{PN}=E-rI$ & $E$ en V, $r$ en $\Omega$ \\ \hline
Loi de Pouillet & $I=\dfrac{E}{\sum R + r}$ & A \\ \hline
Puissance \'electrique & $P=UI$ & W \\ \hline
\'Energie \'electrique & $W=UIt=P\,t$ & J (ou kWh) \\ \hline
Puissance Joule & $P_J=RI^2$ & W \\ \hline
Bilan g\'en\'erateur & $P_{\text{fournie}}=EI$ ; $P_{\text{utile}}=EI-rI^2$ & W \\ \hline
Rendement & $\eta=\dfrac{P_{\text{utile}}}{P_{\text{totale}}}=\dfrac{U}{E}$ & sans unit\'e \\ \hline
\end{tabularx}
\end{center}

\textbf{Composants connus}
\begin{itemize}
  \item \cle{R\'esistor} : caract\'eristique droite passant par l'origine, de pente $\dfrac{1}{R}$ ; dip\^ole passif sym\'etrique.
  \item \cle{Diode} : dip\^ole non sym\'etrique ; bloqu\'ee si $U<U_s$ (tension seuil $U_s\approx\num{0,6}$ V pour le silicium), passante si $U\geq U_s$ ; mod\`ele id\'ealis\'e : $U=U_s$ d\`es qu'elle conduit.
  \item \cle{Transistor} (npn) : fonctionne en interrupteur command\'e ; bloqu\'e si $I_B=0$ ($I_C=0$), satur\'e si $I_B$ suffisant ($V_{CE}\approx 0$) ; en r\'egime lin\'eaire $I_C=\beta\,I_B$.
\end{itemize}

\textbf{M\'ethodes express}
\begin{itemize}
  \item \emph{Point de fonctionnement graphique} : (1) tracer la caract\'eristique $I(U)$ du dip\^ole ; (2) tracer la droite du g\'en\'erateur $I=\dfrac{E-U}{r}$ (points $U=0$, $I=E/r$ et $I=0$, $U=E$) ; (3) lire les coordonn\'ees $(U_0\,;\,I_0)$ de l'intersection.
  \item \emph{Point de fonctionnement alg\'ebrique} : \'ecrire la loi des mailles $E-rI=U_{\text{dip\^ole}}(I)$, remplacer $U_{\text{dip\^ole}}$ par son expression ($RI$, $U_s$, \dots) et r\'esoudre.
  \item \emph{Loi de Pouillet} : v\'erifier que le circuit est \textbf{s\'erie}, sommer toutes les r\'esistances et la r\'esistance interne, puis $I=\dfrac{E}{R_{\text{totale}}}$.
\end{itemize}
\end{bilan}

\begin{aretenir}[Checklist avant le Probatoire]
\begin{itemize}
  \item[$\square$] Je sais d\'efinir la caract\'eristique courant--tension d'un dip\^ole et la tracer \`a partir d'un tableau de mesures.
  \item[$\square$] Je distingue un dip\^ole passif (courbe passant par l'origine) d'un dip\^ole actif.
  \item[$\square$] Je sais que le point de fonctionnement $Q(U_0\,;\,I_0)$ est l'intersection de la caract\'eristique du dip\^ole et de la droite d'attaque du g\'en\'erateur.
  \item[$\square$] Je sais tracer la droite $U=E-rI$ \`a partir de deux points : $(I=0\,;\,U=E)$ et $(U=0\,;\,I=E/r)$.
  \item[$\square$] J'\'enonce et j'applique la loi de Pouillet $I=\dfrac{E}{\sum R+r}$ uniquement pour un circuit s\'erie.
  \item[$\square$] Je connais le mod\`ele de la diode : bloqu\'ee si $U<U_s$, passante avec $U\approx U_s$ ; je v\'erifie toujours si elle conduit avant de calculer.
  \item[$\square$] Je connais les trois r\'egimes du transistor : bloqu\'e ($I_B=0$), lin\'eaire ($I_C=\beta I_B$), satur\'e ($V_{CE}\approx 0$).
  \item[$\square$] Je calcule une puissance $P=UI$, une \'energie $W=Pt$ et je convertis correctement entre joules et kilowattheures ($1$ kWh $=3,6\times10^6$ J).
  \item[$\square$] Je sais faire le bilan de puissance d'un g\'en\'erateur : $EI=UI+rI^2$, et calculer le rendement $\eta=U/E$.
  \item[$\square$] Je v\'erifie l'homog\'en\'eit\'e de chaque formule et j'exprime chaque r\'esultat avec son unit\'e SI et des chiffres significatifs coh\'erents.
  \item[$\square$] Je lis correctement un point sur un graphique en respectant les \'echelles des axes (attention aux graduations !).
  \item[$\square$] Je r\'edige proprement : sch\'ema du circuit, fl\`eches de tension et de courant, loi des mailles \'ecrite avant tout calcul.
\end{itemize}
\end{aretenir}

\findecours

\end{document}
